\documentclass[TQFT_main]{subfiles}

\begin{document}

% \setcounter{}{}
\chapter{層状化空間・因子化ホモロジー}

~\cite{AFT2014stratified}, ~\cite{AFT2014FH}のレビュー

\section{conically smoothな層状化空間}

\subsection{層状化空間}

\begin{mydef}[label=def:topo-poset]{半順序集合の位相}
    $(P,\, \le)$ を半順序集合とする．
    $P$ 上の位相 $\mathscr{O}_{\le} \subset 2^P$ を以下で定義する：
    \begin{align}
        U \in \mathscr{O}_{\le} \DEF \forall x \in U,\, \forall y \in P,\; \bigl[\, x \le y \IMP y \in U \,\bigr]
    \end{align}
\end{mydef}

実際，空集合の定義から $\emptyset \in \mathscr{O}_{\le}$ であり，$\forall U_1,\, U_2 \in \mathcal{O}_{\le}$ に対して $x \in U_1 \cap U_2$ であることは
\begin{align}
     \forall y \in P,\; x \le y \IMP y \in U_1 \AND y \in U_2
\end{align}
と同値なので $U_1 \cap U_2 \in \mathscr{O}_{\le}$ であり，
さらに勝手な開集合族 $\Familyset[\big]{U_\lambda \in \mathscr{O}_{\le}}{\lambda \in \Lambda}$ に対して
$x \in \bigcup_{\lambda \in \Lambda} U_\lambda$ は
\begin{align}
    \exists \alpha \in \Lambda,\; \forall y \in P,\; x \le y \IMP y \in U_\alpha \subset \bigcup_{\lambda \in \Lambda} U_\lambda 
\end{align}
と同値であるから $\bigcup_{\lambda \in \Lambda} U_\lambda \in \mathscr{O}_{\le}$ であり，$\mathscr{O}_{\le}$ は集合 $P$ の位相である．

\begin{myexample}[label=ex:topo-poset]{{$[n]$} の位相}
    半順序集合 $[2] \coloneqq \{0 \le 1 \le 2\}$ を考える．このとき，\hyperref[def:topo-poset]{位相 $\mathscr{O}_{\le}$}とは
    \begin{align}
        \mathscr{O}_{\le} = \bigl\{\, \emptyset,\, \{2\},\, \{1,\, 2\},\, \{0,\, 1,\, 2\} \,\bigr\} 
    \end{align}
    のことである．同様に，半順序集合 $[n] \coloneqq \{0 \le 1 \le \cdots \le n\}$ に対して
    \begin{align}
        \mathscr{O}_{\le} = \bigl\{\, \emptyset,\, \{n\},\, \{n-1,\, n\},\, \dots,\, \{0,\, \dots,\, n\}  \,\bigr\}
    \end{align}
    が成り立つ．
\end{myexample}

\begin{mydef}[label=def:stratified-space]{層状化空間・層状化写像}
    $(P,\, \le)$ を半順序集合とし，定義\ref{def:topo-poset}の位相を入れて位相空間にする．
    
    このとき，位相空間 $X$ が $\bm{P}$\textbf{-層状化}されている ($P$-stratified) とは，連続写像 $s \colon X \lto P$ が存在することを言う\footnote{$s$ のことを，$X$ の\textbf{層状化} (stratification) と呼ぶ．}．
    組 $(X,\, s \colon X \lto P)$ のことを\textbf{$\bm{P}$-層状化空間} ($P$-stratified space) と呼ぶ．
    また，$i \in P$ の逆像 $\bm{X_i} \coloneqq s^{-1}(\{i\}) \subset X$ のことを\textbf{$\bm{i}$-層} ($i$-strata) と呼ぶ．

    \tcblower

    層状化空間 $(X,\, s \colon X \lto P),\; (X',\, s' \colon X' \lto P')$ の間の\textbf{層状化写像} (stratified map) とは，連続写像の組み $(f \colon X \lto X',\; \tilde{f} \colon P \lto P')$ であって以下の図式を可換にするもののこと：
    \begin{center}
        % https://tikzcd.yichuanshen.de/#N4Igdg9gJgpgziAXAbVABwnAlgFyxMJZABgBpiBdUkANwEMAbAVxiRAA0QBfU9TXfIRQBGclVqMWbdgHJuvEBmx4CRMsPH1mrRCAAK8vssFFRG6lqm69cruJhQA5vCKgAZgCcIAWyRkQOBBIAEwWkjogAMog1Ax0AEYwDHr8KkIgHliOABY4hiCePkiiAUGIAMxh2myRtgqFvoj+gcVVVgX5DSHULRWxCUkpxqq6mTl5bRGO3BRcQA
\begin{tikzcd}
X \arrow[d, "s"'] \arrow[r, red,"f"] & X' \arrow[d, "s'"] \\
P \arrow[r, red,"\tilde{f}"']                & P'                
\end{tikzcd}
    \end{center}
\end{mydef}

\begin{myexample}[label=ex:strat-n]{{$[n]$}-層状化空間}
    半順序集合 $[n] \coloneqq \{0 \le \, \cdots \le n\}$ に対して\exref{ex:topo-poset}の位相を入れる．
    まず，
    \begin{align}
        X_0 = s^{-1} \bigl( [n] \setminus \{1,\, \dots,\, n\} \bigr) 
    \end{align}
    でかつ $\{1,\, \dots,\, n\}$ は $[n]$ の開集合であるから，$s$ の連続性から $X$ の部分空間 $X_0 \subset X$ は閉集合だとわかる．
    さらに
    \begin{align}
        X_0 \cup X_1 &= s^{-1} \bigl( [n] \setminus \{2,\, \dots,\, n\} \bigr), \\
        X_0 \cup X_1 \cup X_2 &= s^{-1} \bigl( [n] \setminus \{3,\, \dots,\, n\} \bigr), \\
        &\vdots \\
        X_0 \cup \cdots \cup X_n &= X
    \end{align}
    が成り立つことから，$s$ の連続性より $X$ の部分空間 $X_0 \cup \cdots \cup X_{m \le n}$ は閉集合だと分かる．
\end{myexample}

\begin{myexample}[label=ex:str-CW]{CW複体}
    CW複体 $X$ を与える．$X_{\le k}$ を $X$ の $k$-骨格とするとき，$X_k \setminus X_{k-1}$ を $k \in \mathbb{Z}_{\ge 0}$ に写す写像 $s \colon X \lto \mathbb{Z}_{\ge 0}$ は $X$ の\hyperref[def:stratified-space]{層状化}を与える．
\end{myexample}

直観的には，層状化空間とはdefect付き位相多様体の一般化である．
特に $X$ を位相多様体とするとき，$[n]$-層状化空間 $(X,\, s \colon X \lto [n])$ の $i$-層 $X_i$ とは，多様体 $X$ 上の余次元 $d-i$ のdefectを全て集めてきたものだと見做せる．

\begin{mydef}[label=def:strat-emb]{層状化開埋め込み}
    \hyperref[def:stratified-space]{層状化写像} $(f,\, \tilde{f}) \colon (X,\, s \colon X \lto P) \lto (X',\, s' \colon X' \lto P')$ が\textbf{層状化開埋め込み} (stratified open embedding) であるとは，以下の3条件を充たすことを言う：
    \begin{enumerate}
        \item 連続写像 $f \colon X \lto X'$ は位相的開埋め込みである\footnote{i.e. $f \colon X \lto f(X)$ が同相写像かつ $f(X) \subset Y$ が開集合}
        \item 連続写像 $\bar{f} \colon P \lto P'$ は位相的開埋め込みである．
        \item $\forall p \in P$ に対して，$f$ の \hyperref[def:stratified-space]{$p$-strata}への制限
        \footnote{\hyperref[def:stratified-space]{層状化写像}の定義に登場する図式の可換性より，$\forall x \in X_p$ に対して $s' \bigl( f(x) \bigr) = s' \circ f (x) = \tilde{f} \circ s (x) = \tilde{f}(p)$，i.e. $f(x) \in s'{}^{-1} \bigl( \{\tilde{f}(p)\} \bigr) = X'_{\tilde{f}(p)}$ が分かる．}
        \begin{align}
            f|_{X_p} \colon X_p \lto X'_{\tilde{f}(p)}
        \end{align}
        は位相的開埋め込みである．
    \end{enumerate}
\end{mydef}

\begin{marker}
    以下では混乱が生じにくい場合，\hyperref[def:stratified-space]{層状化空間} $(X,\, s \colon X \lto P)$ のことを $(X \xrightarrow{s} P)$ や $(X \to P)$ と略記する．
    さらに，\hyperref[def:stratified-space]{層状化写像} $(f,\, \tilde{f}) \colon (X,\, s \colon X \lto P) \lto (X',\, s' \colon X' \lto P')$ のことを $f \colon (X \to P) \lto (X' \to P')$ と略記し，連続写像 $\tilde{f} \colon P \lto P'$ のことも $f$ と書く場合がある．
\end{marker}

\begin{mydef}[label=def:StTop]{層状化空間と層状化開埋め込みの成す {$(1,\, 1)$} 圏}
    $(1,\, 1)$-圏 $\StTOP$ を，
    \begin{itemize}
        \label{def:StTop}
        \item 第2可算なHausdorff空間であるような\hyperref[def:stratified-space]{層状化空間}を対象とする
        \item \hyperref[def:strat-emb]{層状化開埋め込み}を射とする
    \end{itemize}
    ことで定義する．
\end{mydef}


\subsection{$C^0$ 級層状化空間}


\begin{mydef}[label=def:str-cone,breakable]{コーン}
    \hyperref[def:stratified-space]{層状化空間} $(X \xrightarrow{s} P)$ を与える．
    $X$ の\textbf{コーン} (cone) とは，以下のようにして構成される\hyperref[def:stratified-space]{層状化空間} $\bigl( \Cone{X},\, \Cone{s} \colon \Cone{X} \lto \Cone{P} \bigr)$ のこと：
    \begin{itemize}
        \item 位相空間 $\Cone{X}$ を，\hyperref[def:pullback-pushout]{押し出し}位相空間
        \begin{align}
            \Cone{X} \coloneqq \{\mathrm{pt}\} \amalg_{\{0\} \times X} (\mathbb{R}_{\ge 0} \times X)
        \end{align}
        と定義する：
        \begin{center}
            % https://tikzcd.yichuanshen.de/#N4Igdg9gJgpgziAXAbVABwnAlgFyxMJZABgBpiBdUkANwEMAbAVxiRAB13hjOBfTvAFt4AAgAaIXqXSZc+QigCM5KrUYs2nQXRwALAEb7gAJV4B9YJwDmMEcV4iBWYXHGTpIDNjwEiZRar0zKyIHFxaOroAToLAaDj87LzuMt7yRMoB1EEaoZyW7Np6MXEJfI7sdNoMVhb5PEkVQqJiDgAUEXqGJuYFNnYOTi7iAJSSqjBQNggooABmURCCSGQgOBBIAEzUDHT6MAwACrI+CiBRWFa6OCkgC0tIymsbiADMUvOLy4ir64-Z6hCYW45SGok60ViWCg5gkHzuXy21D+bx2ewOxzSvlCFyuN14FF4QA
        \begin{tikzcd}
\{0\}\times X \arrow[d] \arrow[r, hookrightarrow, "\{0\} \times \mathrm{id}_X"] & \mathbb{R}_{\ge 0} \times X \arrow[d,red]                                 \\
\{\mathrm{pt}\} \arrow[r,hookrightarrow,red]                                       & \textcolor{red}{\{\mathrm{pt}\} \amalg_{\{0\} \times X} (\mathbb{R}_{\ge 0} \times X)}
        \end{tikzcd}
        \end{center}
        \item 半順序集合 $\Cone{P}$ を，$P$ に最小の要素 $-\infty$ を付け足すことで定義する．これは半順序集合の圏における押し出し
        \begin{align}
            \Cone{P} \coloneqq \{-\infty\} \amalg_{\{0\} \times P} \bigl( [1] \times P \bigr) 
        \end{align}
        である．
        \item 連続写像
        \begin{align}
            \mathbb{R}_{\ge 0} \times X &\lto [1] \times P, \\
            (t,\, x) &\lmto 
            \begin{cases}
                \bigl( 0,\, s(x) \bigr), &t=0, \\
                \bigl( 1,\, s(x) \bigr), &t>0
            \end{cases}
        \end{align}
        が誘導する連続写像 $\Cone{X} \lto \Cone{P}$ を $\Cone{s}$ と書く．
    \end{itemize}
\end{mydef}

位相空間の圏における押し出しの公式から，位相空間 $\Cone{X}$ とは
\begin{align}
    i_1 \colon \{0\} \times X &\lto \mathbb{R}_{\ge 0} \times X,\; x \lmto (0,\, x), \\
    i_2 \colon \{0\} \times X &\lto \{\mathrm{pt}\},\; x \lmto \mathrm{pt}
\end{align}
とおいたときの\hyperref[def:eq-coeq]{コイコライザ}
\begin{center}
    % https://tikzcd.yichuanshen.de/#N4Igdg9gJgpgziAXAbVABwnAlgFyxMJZABgBpiBdUkANwEMAbAVxiRAB13hjOBfTvAFt4AAgAaIXqXSZc+QigCM5KrUYs2nYJ0F0cACwBOg4Ghz92FuroYBzEQAode-QCNXwAEq8A+tva2MCLEvCICWMJw4gCUktIgGNh4BEQATCrU9MysiBzsugZwAGbAAMK8DmKxvKowUIEIKKBFhhCCSGQgOBBIyiAMWGA5IHAQA1Ag1PowdBOIYEwMDNQ4dFgMbJBDkyP6WEU4SAC0fVkauVg+inHNre2Ind291HB7B8-9dK4wDAAKsskFCBDFhbPpDpl1MNLqkbiAWm0Pk9EKkarwgA
    \begin{tikzcd}
    \{0\}\times X \arrow[r, "i_1", shift left] \arrow[r, "i_2"', shift right] & \{\mathrm{pt}\}\sqcup (\mathbb{R}_{\ge 0} \times X) \arrow[r, "q"] & \mathsf{C}(X),
    \end{tikzcd}
\end{center}
i.e. 商位相空間
\begin{align}
    \frac{\mathbb{R}_{\ge 0} \times X}{i_1 (x) \sim i_2 (x)} = \frac{\mathbb{R}_{\ge 0} \times X}{\{0\} \times X}
\end{align}
のこと．従って $\Cone{s} \colon \Cone{X} \lto \Cone{P}$ とは，連続写像\footnote{$\Cone{P}$ の\hyperref[def:topo-poset]{位相 $\mathscr{O}_{\Cone{P}}$}は，$P$ の位相 $\mathscr{O}_{P}$ に1つの開集合 $\{-\infty\} \cup P$ を加えたものである．$\forall U \in \mathscr{O}_{P}$ に対して $\Cone{s}^{-1} (U) = \mathbb{R}_{\textcolor{red}{>0}} \times s^{-1}(U) \in \mathscr{O}_{\Cone{X}}$ で，かつ $\Cone{s}^{-1}(\{-\infty\} \cup P) = \Cone{X} \in \mathscr{O}_{\Cone{X}}$ なので $\Cone{s}$ は連続である．}
\begin{align}
    \label{def:stratification-cone}
    \frac{\mathbb{R}_{\ge 0} \times X}{\{0\} \times X} &\lto \Cone{P},\;
    [(t,\, x)] \lmto \begin{cases}
        -\infty, &t = 0 \\
        s(x), &t > 0
    \end{cases}
\end{align}
のことである．
% また，コーンポイントのみからなる1点集合 $\{\mathrm{pt}\} \subset X$ は $q^{-1} \bigl( \{\mathrm{pt}\} \bigr) = \{0\} \times X$ を充たすが，$\{0\} \times X$ は 

\begin{marker}
    以下では，混乱の恐れがない限り\hyperref[def:stratified-space]{層状化空間} $(X \xrightarrow{s} P)$ の\hyperref[def:str-cone]{コーン}を $\Cone{X \xrightarrow{s} P}$ と略記する．
    さらに，1点 $[0,\, *] \in \Cone{Z}$ のことを\textbf{コーンポイント}と呼び，$\mathrm{pt}$ と略記する．
\end{marker}

素材となる層状化空間 $(Z \xrightarrow{s} P)$ が連結でなくとも，位相空間 $\Cone{Z}$ においてはコーンポイントを経由することができるので，弧状連結になる．

\begin{mylem}[label=lem:cone-connected]{層状化空間のコーンは連結}
    \hyperref[def:stratified-space]{層状化空間} $(Z \xrightarrow{s} P)$ の\hyperref[def:str-cone]{コーン} $\Cone{Z \xrightarrow{s} P}$ において，位相空間 $\Cone{Z}$ は弧状連結である\footnote{従って連結でもある}．
\end{mylem}

\begin{proof}
    $Z$ を弧状連結成分のdisjoint unionに分割し，$Z = \coprod_{i \in \pi_0(Z)} Z_i$ と書く．
    いま，位相空間 $\Cone{Z}$ の互いに相異なる2点 $[t_i,\, z_i],\, [t_j,\, z_j] \in \Cone{Z}$ を任意にとる．
    ただし $z_i \in Z_i,\, z_j \in Z_j$ とする．

    $i=j$ ならば，$z_i$ と $z_j$ を繋ぐ $Z_i = Z_j$ 上の連続曲線 $\varphi \colon [0,\, 1] \lto Z_i$ が存在する．よって連続写像
    \begin{align}
        \bar{\varphi} \colon [0,\, 1] &\lto \Cone{Z},\; 
        s \lmto \bigl[ (1-s)t_i + st_j,\, \varphi(s) \bigr] 
    \end{align}
    が $[t_i,\, z_i],\, [t_j,\, z_j]$ を繋ぐ．

    $i \neq j$ とする．このとき，$\Cone{Z}$ 上の連続曲線
    \begin{align}
        \varphi_i \colon [0,\, 1] &\lto \Cone{Z},\; s \lmto \bigl[ (1-s)t_i,\, z_i \bigr] \\
        \varphi_j \colon [0,\, 1] &\lto \Cone{Z},\; s \lmto \bigl[ (1-s)t_j,\, z_j \bigr]
    \end{align}
    はそれぞれ $[t_i,\, z_i],\, [t_j,\, z_j]$ とコーンポイント $\mathrm{pt} = [0,\, *] \in \Cone{Z}$ を繋ぐ．故に
    $\varphi_i$ と $\varphi_j^{-1}$ を繋ぎ合わせて得られる $\Cone{Z}$ 上の連続曲線
    \begin{align}
        \varphi_i * \varphi_j^{-1} \colon [0,\, 1] &\lto \Cone{Z},\; 
        s \lmto \begin{cases}
            \varphi_i (2s), &s \in [0,\, \frac{1}{2}] \\
            \varphi_i (2-2s), &s \in [\frac{1}{2},\, 1]
        \end{cases}
    \end{align}
    がちょうど $[t_i,\, z_i],\, [t_j,\, z_j]$ を繋ぐ．
\end{proof}

\begin{myexample}[label=ex:Cone-dim1]{次元1のコーン}
    \hyperref[def:covering-dim]{次元} $0$ のコンパクトな\hyperref[def:stratified-space]{層状化空間}とは，有限個の点のdisjoint unionのことである．
    これらのコーンを図示すると，
    \begin{align}
        \Cone{\underset{1}{\bullet} \cdots \underset{n}{\bullet}}
        &=\quad
        \begin{tikzpicture}[baseline={([yshift=-.5ex]current bounding box.center)}]
            \node[bullet] (O) at (0,0) {};
            \foreach \q in {0,60,...,300} {
                \draw[very thick] (O) -- ++(\q:2);
            }
            \node[below right] at (-120:2) {$1$};
            \node[below left] at (-60:2) {$n$};
            \foreach \q in {-30,30,...,210} {
                \node[rotate={\q + 90}] at ($(O)+(\q:1.5)$) {$\cdots$};
            }
        \end{tikzpicture}
    \end{align}
    のようになる．ただし，コーンポイントは中心の点である．
\end{myexample}

\begin{myexample}[label=ex:Cone-dim2]{次元2のコーン}
    \hyperref[def:covering-dim]{次元} $1$ のコンパクトな\hyperref[def:stratified-space]{層状化空間}とは，閉じた非空な無向グラフのことである．
    特に，頂点のない連結かつコンパクトな無向グラフは $S^1$ と同相であり，これの\hyperref[def:str-cone]{コーン}をとることで得られる2次元の層状化空間は次のように図示できる：
    % これらのコーンを図示すると，
    \begin{align}
        \label{fig:Cone-S1}
        \Cone{S^1}
        &=\quad
        \begin{tikzpicture}[baseline={([yshift=-.5ex]current bounding box.center)}]
            \path coordinate[bullet] (O)
            ;
            \draw[dashed] (O.center) circle (0.5);
            \begin{scope}[on background layer]
                \fill[gray!20] (O.center) circle (1);
            \end{scope}
        \end{tikzpicture}
    \end{align}
    さらに，$n$ 個の頂点が付いた $S^1$ のコーンを図示すると次のようになる：
    \begin{align}
        \Cone{
            \begin{tikzpicture}[baseline={([yshift=-.5ex]current bounding box.center)}]
                \path coordinate (O)
                ;
                % contour
                \draw[gray] (O) circle (1);
                % spoke
                \foreach \q in {0,60,...,300} {
                    \node[vertex] at ($(O)+(\q:1)$) {};
                }
                \node[below left] at (-120:1) {$1$};
                \node[below right] at (-60:1) {$n$};
                \foreach \q in {-30,30,...,210} {
                    \node[rotate={\q + 90}] at ($(O)+(\q:1.2)$) {$\cdots$};
                }
            \end{tikzpicture}
        }
        &= \quad
        \begin{tikzpicture}[baseline={([yshift=-.5ex]current bounding box.center)}]
            \path coordinate[bullet] (O)
            ;
            % contour
            \draw[dashed] (O.center) circle (1);
            % spoke
            \foreach \q in {0,60,...,300} {
                \draw[very thick] (O.center) -- ++(\q:2);
            }
            \node[below right] at (-120:2) {$1$};
            \node[below left] at (-60:2) {$n$};
            \foreach \q in {-30,30,...,210} {
                \node[rotate={\q + 90}] at ($(O.center)+(\q:1.2)$) {$\cdots$};
            }
            % background
            \begin{scope}[on background layer]
                \fill[gray!20] (O.center) circle (2);
            \end{scope}
        \end{tikzpicture}
    \end{align}
    他にも，例えば4面体の1-骨格のコーンを図示すると次のようになる：
    \begin{align}
        \Cone{ 
            \begin{tikzpicture}[baseline={([yshift=-.5ex]current bounding box.center)}]
                \pgfmathsetmacro{\a}{2}
                \path coordinate (O)
                +(90:\a) coordinate (P_1)
                +(0:{0.8*\a}) coordinate (P_2)
                +(-160:\a) coordinate (P_3)
                +(-60:\a) coordinate (P_4)
                ;
                % edges
                \draw[gray,name path=e_12] (P_1) -- (P_2);
                \draw[gray,name path=e_13] (P_1) -- (P_3);
                \draw[gray,name path=e_14,spath/save=e_14] (P_1) -- (P_4);
                \draw[gray,name path=e_23] (P_2) -- (P_3);
                \draw[gray,name path=e_24] (P_2) -- (P_4);
                \draw[gray,name path=e_34] (P_3) -- (P_4);
                \begin{scope}
                    % \draw[gray,save named path=e_14] (P_1) -- (P_4);
                    % \draw[gray,save named path=e_23] (P_2) -- (P_3);
                    \clip[name intersections={of=e_14 and e_23}] (intersection-1) circle (5pt);
                    \draw[spath/use=e_14,white,ultra thick,double=gray,double distance=0.4pt];
                \end{scope}
                % vertices
                \foreach \i in {1,...,4} {
                    \node[vertex] at (P_\i) {};
                }
            \end{tikzpicture}
        } 
        &= \quad
        \begin{tikzpicture}[baseline={([yshift=-.5ex]current bounding box.center)}]
            \pgfmathsetmacro{\a}{4}
            \pgfmathsetmacro{\b}{2}
            \path coordinate (O)
            +(90:\a) coordinate (P_1)
            +(0:{0.8*\a}) coordinate (P_2)
            +(-160:\a) coordinate (P_3)
            +(-60:\a) coordinate (P_4)
            +(90:\b) coordinate (Q_1)
            +(0:{0.8*\b}) coordinate (Q_2)
            +(-160:\b) coordinate (Q_3)
            +(-60:\b) coordinate (Q_4)
            ;
            % Link
            \draw[gray,dashed] (Q_1) -- (Q_2);
            \draw[gray,dashed] (Q_1) -- (Q_3);
            \draw[gray,dashed] (Q_1) -- (Q_4);
            \draw[gray,dashed] (Q_2) -- (Q_3);
            \draw[gray,dashed] (Q_2) -- (Q_4);
            \draw[gray,dashed] (Q_3) -- (Q_4);
            % edges
            \draw[gray,name path=e_12] (P_1) -- (P_2);
            \draw[gray,name path=e_13] (P_1) -- (P_3);
            \draw[gray,name path=e_14,spath/save=e_14] (P_1) -- (P_4);
            \draw[gray,dotted,name path=e_23] (P_2) -- (P_3);
            \draw[gray,name path=e_24] (P_2) -- (P_4);
            \draw[gray,name path=e_34] (P_3) -- (P_4);
            \draw[very thick,name path=e_1] (O) -- (P_1);
            \draw[very thick,name path=e_2,spath/save=e_2] (O) -- (P_2);
            \draw[very thick,name path=e_3] (O) -- (P_3);
            \draw[very thick,name path=e_4,spath/save=e_4] (O) -- (P_4);
            \begin{scope}
                \clip (P_1) -- (O) -- (P_4) -- cycle;
                \draw[spath/use=e_2,white,ultra thick];
                \draw[spath/use=e_2,dotted,very thick];
            \end{scope}
            \begin{scope}
                \clip[name intersections={of=e_14 and e_23}] (intersection-1) circle (5pt);
                \draw[spath/use=e_14,white,ultra thick,double=gray,double distance=0.4pt];
            \end{scope}
            \begin{scope}
                \clip[name intersections={of=e_14 and e_2}] (intersection-1) circle (5pt);
                \draw[spath/use=e_14,white,ultra thick,double=gray,double distance=0.4pt];
            \end{scope}
            \begin{scope}
                \clip[name intersections={of=e_23 and e_4}] (intersection-1) circle (5pt);
                \draw[spath/use=e_4,white,ultra thick,double=black,double distance=1.2pt];
            \end{scope}
            % vertex
            \node[bullet] at (O) {};
            % faces
            \fill[fill=green,opacity=0.1,save named path=f_12] (P_1) -- (O) -- (P_2) -- cycle;
            \fill[fill=red,opacity=0.1,save named path=f_13] (P_1) -- (O) -- (P_3) -- cycle;
            \fill[fill=blue,opacity=0.1,save named path=f_14] (P_1) -- (O) -- (P_4) -- cycle;
            \fill[fill=DarkGreen,opacity=0.1,save named path=f_23] (P_2) -- (O) -- (P_3) -- cycle;
            \fill[fill=DarkOrange,opacity=0.1,save named path=f_24] (P_2) -- (O) -- (P_4) -- cycle;
            \fill[fill=LightSlateBlue,opacity=0.1,save named path=f_34] (P_3) -- (O) -- (P_4) -- cycle;
        \end{tikzpicture}
    \end{align}
    なお，$\Cone{Z \xrightarrow{s} P}$ において，$Z$ の連結成分毎の次元が互いに異なる場合を考えることもできる．例えば $Z = S^1 \amalg \{\mathrm{pt}\}$ のとき，
    \begin{align}
        \Cone{
            \begin{tikzpicture}[baseline={([yshift=-.5ex]current bounding box.center)}]
                \path coordinate (O)
                +(2,0) coordinate (p)
                ;
                \draw[gray] (O) circle (1);
                \node[vertex] at (p) {};
            \end{tikzpicture}
        } 
        &= \quad
        \begin{tikzpicture}[baseline={([yshift=-.5ex]current bounding box.center)}]
            \path coordinate (O)
            +(2,0) coordinate (p)
            ;
            % line
            \draw[very thick] (O) -- (p);
            \begin{scope}
                \clip (O) ellipse (1 and 2);
                \draw[white, ultra thick] (O) -- (p);
                \draw[dotted, very thick] (O) -- (p);
            \end{scope}
            % cone
            \begin{scope}
                \fill[gray!20,opacity=0.6] (O) ellipse (1 and 2);
            \end{scope}
            \draw[gray,dashed] (O) ellipse (0.5 and 1);
            % cone point
            \node[bullet] at (O) {};
        \end{tikzpicture}
    \end{align}
    のように図示できる．
\end{myexample}

\begin{mydef}[label=def:Snglr-C0,breakable]{$C^0$ 級層状化空間}
    以下を充たす \hyperref[def:StTop]{$\StTOP$} の\hyperref[def:faithful]{充満部分圏} $\Cat{T}$ のうち最小のものを $\bigcap_{\Cat{T}} \Cat{T} \subset \StTOP$ と書く：
    \begin{description}
        \item[\textbf{(Snglr-1)}] $(\emptyset \to \emptyset) \in \Obj{\Cat{T}}$
        \item[\textbf{(Snglr-2)}]　
        
        $(X \to P) \in \Obj{\Cat{T}}$ かつ $X,\, P$ が位相空間としてコンパクト 
        
        $\IMP$ $\Cone{X \to P} \in \Obj{\Cat{T}}$
        
        \item[\textbf{(Snglr-3)}]　
        
        $(X \to P) \in \Obj{\Cat{T}}$ $\IMP$ $(X \times \mathbb{R} \to P) \in \Obj{\Cat{T}}$\footnote{$X \times \mathbb{R}$ の層状化は，連続写像 $X \times \mathbb{R} \lto X,\; (x,\, t) \lmto x$ を前もって合成することにより定める．}
        
        \item[\textbf{(Snglr-4)}]　
        
        $(X \to P) \in \Obj{\Cat{T}}$ かつ $\Hom{\StTOP} \bigl( (U \to P_U),\, (X \to P) \bigr) \neq \emptyset$ 
        
        $\IMP$ $(U \to P_U) \in \Obj{\Cat{T}}$
        
        \item[\textbf{(Snglr-5)}]　 
        
        $(X \to P) \in \Obj{\StTOP}$ が開被覆 $\Familyset[\big]{(U_\lambda \to P_\lambda) \lto (X \to P)}{\lambda \in \Lambda}$\footnote{i.e. $\Familyset[\big]{U_\lambda}{\lambda \in \Lambda},\; \Familyset[\big]{P_\lambda}{\lambda \in \Lambda}$ が，それぞれ位相空間 $X,\, P$ の開被覆を成す．} を持ち，かつ $\forall \lambda \in \Lambda$ に対して $(U_\lambda \to P_\lambda) \in \Obj{\Cat{T}}$ 
        
        $\IMP$ $(X \to P) \in \Obj{\Cat{T}}$
    \end{description}
    \tcblower
    $(1,\, 1)$-圏 $\bigcap_{\Cat{T}} \Cat{T}$ の対象のうち，位相空間としてパラコンパクト
    \footnote{この仮定により，任意の $C^0$ 級層状化空間 $(X \to P)$ において位相空間 $X$ が\textbf{完全正規} (completely normal/hereditarily normal) であることが分かる．この事実は\hyperref[def:covering-dim]{被覆次元}および\hyperref[def:dim-depth]{深さ}を使用する際に有用である．特に，補題\ref{lem:strat-basic}の証明で暗に使用している．}
    であるもの全体が成す充満部分圏を $\SnglrC \subset \bigcap_{\Cat{T}} \Cat{T} \subset \StTOP$ と書く．$\SnglrC$ の対象を\textbf{$\bm{C^0}$ 級層状化空間} ($C^0$ stratified space) と呼ぶ．
\end{mydef}

\begin{myexample}[label=ex:topomfld]{位相多様体は $C^0$ 級層状化空間}
    \hyperref[def:Snglr-C0]{\textsf{\textbf{(Snglr-1)}}}より，$* \coloneqq \Cone{\emptyset \to \emptyset} \in \Obj{\SnglrC}$ である．
    \hyperref[def:Snglr-C0]{\textsf{\textbf{(Snglr-3)}}}より，$\forall n \ge 0$ に対して $\mathbb{R}^n = * \times (\mathbb{R}^n \to [0]) \in \Obj{\SnglrC}$ であることが帰納的に分かる．
    $\mathbb{R}^n$ の任意の開集合 $U \hookrightarrow \mathbb{R}^n$ に対して，
    \begin{center}
            \begin{tikzcd}
        U \arrow[d] \arrow[r, hookrightarrow] & \mathbb{R}^n \arrow[d] \\
        {[0]} \arrow[r, "="']                & {[0]}               
    \end{tikzcd}
    \end{center}
    は\hyperref[def:strat-emb]{層状化埋め込み}であり，従って \hyperref[def:Snglr-C0]{\textsf{\textbf{(Snglr-4)}}}より $U \coloneqq (U \to [0]) \in \Obj{\SnglrC}$ が分かる．
    以上の考察と\hyperref[def:Snglr-C0]{\textsf{\textbf{(Snglr-5)}}}を併せて，任意の位相多様体 $M$ は\footnote{より正確には，$M$ を\hyperref[def:stratified-space]{層状化空間} $(M \to [0])$ と同一視している．}圏 $\SnglrC$ の対象である．
\end{myexample}

\exref{ex:topomfld}の意味で，\hyperref[def:Snglr-C0]{$C^0$ 級層状化空間}は位相多様体の一般化と見做せる．しかしまだそこには $C^\infty$ 構造を一般化した構造は入っておらず，$C^\infty$ 多様体の一般化とは見做せない．

% \begin{mylem}[label=lem:strat-surj]{$C^0$ 級層状化空間の層状化は全射}
%     \textcolor{red}{$\forall (X \xrightarrow{s} P) \in \Obj{\SnglrC}$ に対して，連続写像 $s$ は全射である．}\footnote{(2025/12/6) ~\cite[Remark 2.1.18.]{AFT2014stratified}の証明のスケッチには反例があると思われる．}
% \end{mylem}

% \begin{proof}
%     $s$ が全射であるような\hyperref[def:stratified-space]{層状化空間} $(X \to P)$ 全体が成す充満部分圏を $\Cat{T} \subset \StTOP$ と書く．$\Cat{T}$ が条件 \hyperref[def:Snglr-C0]{\textsf{\textbf{(Snglr-1)}}-\textsf{\textbf{(Snglr-5)}}}を充たすことを示せば十分である\footnote{これらの条件を充たす最小の $(1,\, 1)$-圏のことを $\SnglrC$ と定義したため，$\SnglrC \subset \Cat{T}$ が言えるのである．}．

%     \begin{description}
%         \item[\textbf{(Snglr-1)}] $(\emptyset \to \emptyset) \in \Obj{\Cat{T}}$ は自明である．
%         \item[\textbf{(Snglr-2)}] $\forall (X \xrightarrow{s} P) \in \Obj{\Cat{T}}$ に対して，\eqref{def:stratification-cone}よりその\hyperref[def:str-cone]{コーン}の層状化は $\Cone{s} \bigl( \Cone{X} \bigr) = s(X) \cup \{-\infty\} = P \cup \{-\infty\} = \Cone{P}$ を充たし，全射である．
%         i.e. $\Cone{X \xrightarrow{s} P} \in \Obj{\Cat{T}}$ である．
%         \item[\textbf{(Snglr-3)}] $\forall (X \xrightarrow{s} P) \in \Obj{\Cat{T}}$ に対して，$(X \times \mathbb{R} \to P) \in \Obj{\StTOP}$ の層状化をあからさまに書くと $(x,\, t) \lmto s(x)$ であり，全射である．i.e. $(X \times \mathbb{R} \to P) \in \Obj{\Cat{T}}$ である．
%         \item[\textbf{(Snglr-4)}] 
%         % \textcolor{red}{反例が存在する．}
%         % $X \coloneqq \{x\},\; U \coloneqq \{u\},\; P \coloneqq \{p\},\, P_U \coloneqq \{p_0,\; p_1\}$ とおき，
%         % \begin{align}
%         %     s \colon X &\lto P,\; x \lmto p \\
%         %     s_U \colon U &\lto P_U,\; u \lmto p_0
%         % \end{align}
%         % とおく．$P_U$ の半順序集合としての構造は，$\le\; \coloneqq \bigl\{ (p_0,\, p_0),\, (p_1,\, p_1)\bigr\} \subset P_U \times P_U$ と定める．
%         % このとき $s,\, s_U$ は定義\ref{def:topo-poset}の位相に関してどちらも連続写像であり，$(X \xrightarrow{s} P) \in \Obj{\Cat{T}},\; (U \xrightarrow{s_U} P_U) \in \Obj{\StTOP}$ である．
        
%         % ここで，連続写像
%         % \begin{align}
%         %     i \colon U &\lto X,\; u \lmto x \\
%         %     \bar{i} \colon P_U &\lto P,\; q \lmto p
%         % \end{align}
%         % は $(i,\, \bar{i}) \in \Hom{\StTOP} \bigl( (U \xrightarrow{s_U} P_U),\, (X \xrightarrow{s} P) \bigr)$ を充たすので $\Hom{\StTOP} \bigl( (U \xrightarrow{s_U} P_U),\, (X \xrightarrow{s} P) \bigr) \neq \emptyset$ だが，明らかに $(U \xrightarrow{s_U} P_U) \notin \Obj{\Cat{T}}$ である．
%         % $(X \xrightarrow{s} P) \in \Obj{\Cat{T}}$ および $(U \xrightarrow{s_U} P_U) \in \Obj{\SnglrC}$ について 
%         % $\Hom{\StTOP} \bigl( (U \to P_U),\, (X \to P) \bigr) \neq \emptyset$ が成り立っているとする．
%         % このとき\hyperref[def:strat-emb]{層状化開埋め込み}
%         % % https://tikzcd.yichuanshen.de/#N4Igdg9gJgpgziAXAbVABwnAlgFyxMJZABgBpiBdUkANwEMAbAVxiRAFUQBfU9TXfIRQBGclVqMWbABrdeIDNjwEio4ePrNWiEAAU5fJYKJl11TVJ26A+py7iYUAObwioAGYAnCAFskZEBwIJFFAuiwGNgALCAgAaxBzSW0QLAMQL18kAGZqIKQAJjzwyJ0Y+MSJLTYAHRqAIzpPYCwuSoY6epgGXX5lIRBPLCconHTMv0QA-MRcqssQOFt2zu7eoxUdIZGxng9vSdCZovmUhHsuIA
%         % \begin{tikzcd}
%         % U \arrow[r, "i", hook] \arrow[d, "s_U"'] & X \arrow[d, "s"] \\
%         % P_U \arrow[r, "\bar{i}"', hook]          & P               
%         % \end{tikzcd}
%         % が存在する．

%         \item[\textbf{(Snglr-5)}] 
%         $(X \xrightarrow{s} P) \in \Obj{\StTOP}$ が開被覆 $\Familyset[\Big]{(i_\lambda,\, \bar{i}_\lambda) \in \Hom{\StTOP} \bigl( (U_\lambda \xrightarrow{s_\lambda} P_\lambda),\, (X \xrightarrow{s} P)\bigr)}{\lambda \in \Lambda}$ を持ち，かつ $\forall \lambda \in \Lambda$ に対して $(U_\lambda \xrightarrow{s_\lambda} P_\lambda) \in \Obj{\Cat{T}}$ であるとする．
%         このとき，族 $\Familyset[\big]{\bar{i}_\lambda \colon P_\lambda \lto P}{\lambda \in \Lambda}$ は $P$ の開被覆であり，$\forall p \in P$ に対してある $\lambda \in \Lambda$ が存在して，$p \in \bar{i}_\lambda (P_\lambda)$ を充たす．
%         さらに仮定より $s_\lambda \colon U_\lambda \lto P_\lambda$ は全射なので，ある $u_\lambda \in U_\lambda$ が存在して $p = \bar{i}_\lambda \circ s_\lambda (u_\lambda)$ を充たす．
%         $(1,\, 1)$-圏 \hyperref[def:StTop]{$\StTOP$ の定義}から $(i_\lambda,\, \bar{i}_\lambda)$ は\hyperref[def:strat-emb]{層状化開埋め込み}なので，
%         $p = \bar{i}_\lambda \circ s_\lambda (u_\lambda) = s \circ i_\lambda (u_\lambda) \in s(X)$ だと分かる．i.e. $P = s(X)$ であり，$(X \xrightarrow{s} P) \in \Obj{\Cat{T}}$ が示された．
%     \end{description}
% \end{proof}

\subsection{$C^0$ basic}

\begin{mydef}[label=def:C0-basic]{$C^0$ basic}
    \hyperref[def:Snglr-C0]{$\bm{C^0}$ 級層状化空間} $(X \to P) \in \Obj{\SnglrC}$ が $\bm{C^0}$\textbf{-basic}であるとは，ある $n \in \mathbb{Z}_{\ge 0}$ およびコンパクトな\hyperref[def:Snglr-C0]{$\bm{C^0}$ 級層状化空間} $(Z \to Q) \in \Obj{\SnglrC}$ が存在して
    $(X \to P) = \bigl(\mathbb{R}^n \to [0]\bigr) \times \Cone{Z \to Q}$ が成り立つことを言う．
\end{mydef}

いま，\hyperref[def:C0-basic]{$C^0$ basic}な $(U \to P_U) = (\mathbb{R}^n \to [0]) \times \Cone{Z \to P} \in \Obj{\SnglrC}$ を1つとる．
\hyperref[def:str-cone]{コーンの定義}から，$U$ の点を $(v,\, [t,\, z]) \in \mathbb{R}^n \times \frac{\mathbb{R}_{\ge 0} \times Z}{\{0\} \times Z}$ と表示することができる．
この表示の下で自己同相
\begin{align}
    \label{eq:auto-basic}
    \gamma \colon \mathbb{R}_{> 0} \times T \mathbb{R}^n \times \Cone{Z} &\lto \mathbb{R}_{> 0} \times T \mathbb{R}^n \times \Cone{Z}, \\
    \bigl( t,\, (v,\, p),\, [s,\, z] \bigr) &\lmto \bigl( t,\, (tv + p,\, p),\, [ts,\, z] \bigr)
\end{align}
を考える\footnote{接束 $T\mathbb{R}^n$ は $\mathbb{R}^{2n}$ と微分同相である．~\cite[p.23]{AFT2014stratified}の記法に合わせて底空間 $\mathbb{R}^n$ の点を $p$，$p$ 上のファイバーの元を $v$ としたとき $(v,\, p) \in T \mathbb{R}^n$ と書いた．命題\ref{prop:tangentbundle}の記法と順番が逆なので注意．}．

さらに，もう1つの\hyperref[def:C0-basic]{$C^0$ basic}な $(U' \to P_{U'}) = (\mathbb{R}^{n'} \to [0]) \times \Cone{Z' \to P'} \in \Obj{\SnglrC}$ および
$f \in \Hom{\SnglrC} \bigl( (U \to P_{U}),\, (U' \to P_{U'})  \bigr)$ をとる．ただし，$f$ はコーンポイントをコーンポイントへ写す，
i.e. $f (\mathbb{R}^n \times \{\mathrm{pt}\}) \in \mathbb{R}^{n'} \times \{\mathrm{pt}\}$ が成り立つことを仮定する．
$f|_{\mathbb{R}^n} \colon \mathbb{R}^n \times \{\mathrm{pt}\} \lto \mathbb{R}^{n'} \times \{\mathrm{pt}\}$ を $f$ のコーンポイントへの制限として，
\begin{align}
    f_{\Delta} \colon \mathbb{R}_{\textcolor{red}{>0}} \times T \mathbb{R}^n \times \Cone{Z} &\lto \mathbb{R}_{\textcolor{red}{>0}} \times T \mathbb{R}^{n'} \times \Cone{Z'}, \\
    \bigl( t,\, v,\, p,\, [s,\, z] \bigr) &\lmto \bigl( t,\, f|_{\mathbb{R}^n}(v),\, f(p,\, [s,\, z]) \bigr)
\end{align}
とおこう．

\begin{myexample}[label=ex:cone-diff]{}
    $Z = Z' = \emptyset$ のとき，$f$ とは単に連続関数 $f \colon \mathbb{R}^n \lto \mathbb{R}^{n'}$ のことである．
    このとき，
    \begin{align}
        (\gamma^{-1} \circ f_{\Delta} \circ \gamma)(t,\, v,\, p)
        &= \gamma^{-1} \circ f_\Delta (t,\, tv+p,\, p) \\
        &= \gamma^{-1} \bigl(t,\, f(tv+p),\, f(p)\bigr) \\
        &= \left( t,\, \frac{f(tv+p) - f(p)}{t},\, f(p) \right)
    \end{align}
    と計算できるため，$f$ が $C^1$ 級であることと $\forall (v,\, p) \in T \mathbb{R}^n$ に対して $t \to +0$ の極限，i.e. $v$ に沿った片側方向微分が存在することは同値である．
\end{myexample}

\exref{ex:cone-diff}をもとに，\hyperref[def:C0-basic]{$C^0$ basic}な\hyperref[def:Snglr-C0]{$C^0$ 級層状化空間}の間の\hyperref[def:strat-emb]{層状化開埋め込み}の\textbf{conically smoothness}を定義する．
$C^\infty$ 多様体の $C^\infty$ 構造の定義においては，チャート $(U,\, \varphi \colon \mathbb{R}^n \to U),\, (V,\, \psi \colon \mathbb{R}^n \to V)$ の間の変換関数 $\psi^{-1} \circ \varphi \colon \mathbb{R}^n \lto \mathbb{R}^{n}$ が $C^\infty$ 級であることを要請した．
次の小節で\textbf{conically smooth structure}の定義を行うが，その際にチャートに対応するものは\textbf{basic} $U = \mathbb{R}^n \times \Cone{Z}$ から着目している\hyperref[def:Snglr-C0]{$C^0$-級層状化空間} $X$ への\hyperref[def:strat-emb]{層状化開埋め込み} $\varphi \colon U \hookrightarrow X$ であり，
概ね\footnote{コーンポイントをコーンポイントに写さない変換関数も存在しうるので，これだけではいけない．}2つのチャート $\varphi \colon U \hookrightarrow X,\; \psi \colon V \hookrightarrow X$ の間の変換関数 $\psi^{-1} \circ \varphi \colon U \lto V$ に対して\hyperref[def:c-smooth-along]{conically smooth (along $\mathbb{R}^n$)}であることを要請する．

\begin{mydef}[label=def:c-smooth-along,breakable]{$\mathbb{R}^n$ に沿ってconically smooth}
    \begin{itemize}
        \item \hyperref[def:C0-basic]{$C^0$ basic}な $(U \to P_{U}) = (\mathbb{R}^{n} \to [0]) \times \Cone{Z \to P} \in \Obj{\SnglrC}$
        \item \hyperref[def:C0-basic]{$C^0$ basic}な $(U' \to P_{U'}) = (\mathbb{R}^{n'} \to [0]) \times \Cone{Z' \to P'} \in \Obj{\SnglrC}$
        \item $f \in \Hom{\SnglrC} \bigl( (U \to P_{U}),\, (U' \to P_{U'})  \bigr)$ であって，コーンポイントを保存するもの
    \end{itemize}
    を与える．このとき，$f$ が\textbf{$\bm{\mathbb{R}^n}$ に沿って $C^1$ 級} ($C^1$ along $\mathbb{R}^n$) であるとは，以下の図式を可換にする連続写像
    \begin{align}
        \textcolor{red}{\tilde{D}f} \colon \mathbb{R}_{\textcolor{red}{\ge 0}} \times T \mathbb{R}^n \times \Cone{Z} &\lto \mathbb{R}_{\textcolor{red}{\ge 0}} \times T \mathbb{R}^{n'} \times \Cone{Z'}
    \end{align}
    が存在することを言う\footnote{写像 $\gamma^{-1} \circ f_{\Delta} \circ \gamma$ の連続性から，拡張 $\tilde{D}f$ は存在すれば一意かつ連続である．}：
    \begin{center}
        % https://tikzcd.yichuanshen.de/#N4Igdg9gJgpgziAXAbVABwnAlgFyxMJZABgBpiBdUkANwEMAbAVxiRAB12BbOnACwBGA4ACUAvgH1gnHDAAeOAMYQGEAE7A1MKGOnsA5jAAExMWKMysXeEYAqF7r0HDxAPTAO81uA4DCBGGAALTEQMVJ0TFx8QhQARnIqWkYWNk4efiFRST1ZBWVVDS0dPUMTM08rG3t0pyy3YDAAcnNLbz8A4JawiJAMbDwCIjI4pPpmVkQOR0yXHIA+csr2mpnnbPdlm05-MECQnsiBmKIE0epx1Kna2eypRdMtn1WM9Ybm1vYvbfZd-e6xEltIYEChQAAzNQQLhIABM1BwECQAGZwhCoTDEPCQIikKZepDoSiEUjEHE0SBCZjsbjEMiLilJtN9HQuDxXMAALTkziKLBqRRGcFSTgAERgDBwdE+fIFDhZbLoIGoDDoAglAAUooNYiA1Fh9HwcIdKRi8SSkAlkhM0l8sAxYMBRWJwcqQAwsHs2FA6HA+NowhQxEA
        \begin{tikzcd}
        \mathbb{R}_{\textcolor{red}{\ge 0}} \times T \mathbb{R}^n \times \Cone{Z} \arrow[r, "\tilde{D}f", red, dashed]                         & \mathbb{R}_{\textcolor{red}{\ge 0}} \times T \mathbb{R}^{n'} \times \Cone{Z'} \\
        \mathbb{R}_{> 0} \times T \mathbb{R}^n \times \Cone{Z} \arrow[r] \arrow[u] \arrow[r, "\gamma^{-1}\circ f_{\Delta} \circ \gamma"'] & \mathbb{R}_{> 0} \times T \mathbb{R}^{n'} \times \Cone{Z'} \arrow[u]         
        \end{tikzcd}
    \end{center}
    \tcblower
    このような拡張が存在するとき，第一変数を $t=0$ に制限して得られる連続写像を
    \begin{align}
        \bm{Df} \colon T \mathbb{R}^n \times \Cone{Z} \lto T \mathbb{R}^{n'} \times \Cone{Z'}
    \end{align}
    と書く．$f$ が\textbf{$\bm{\mathbb{R}^n}$ に沿って $C^r$ 級} であるとは，$Df$ が $\mathbb{R}^n$ に沿って $C^{r-1}$ 級であることを言う．
    $f$ が\textbf{$\bm{\mathbb{R}^n}$ に沿ってconically smooth}であるとは，$\forall r \ge 1$ について $C^r$ 級であることを言う．
\end{mydef}

\subsection{conically smoothな層状化空間}

次に行うべきは，与えられた\hyperref[def:Snglr-C0]{$\bm{C^0}$ 級層状化空間} $(X \to P) \in \Obj{\SnglrC}$ の上の\textbf{conically smooth structure} i.e. 変換関数が\hyperref[def:c-smooth-along]{conically smooth}であるような\textbf{極大アトラス}を定義することである．
この手続きは，次で定義する次元と深さに関する帰納法によって構成される．

\begin{mydef}[label=def:covering-dim]{被覆次元}
    $X$ を位相空間とする．以下の条件を充たす最小の $d \in \mathbb{Z}_{\ge -1}$ のことを（存在すれば）$X$ の\textbf{被覆次元} (covering dimension) と呼び，$\bm{\dim X}$ と書く：
    
    \begin{description}
        \item[\textbf{(covering)}]　
        
        $X$ の任意の開被覆 $\mathscr{U}$ に対して，十分細かい細分 $\mathscr{V}_{\mathscr{U}} \prec \mathscr{U}$ が存在して，任意の互いに異なる $\forall m > d+1$ 個の開集合 $V_1,\, \dots,\, V_{m} \in \mathscr{V}_{\mathscr{U}}$ の共通部分が空になるようにできる．特に，$\emptyset$ の被覆次元は $-1$ と定義する．
    \end{description}
    
    \tcblower

    点 $x \in X$ における\textbf{被覆次元}を以下で定義する：
    \begin{align}
        \dim_x X \coloneqq \inf\, \bigl\{\, \dim U \ge -1 \bigm| x \in U \underset{\text{open}}{\subset} X \,\bigr\}
    \end{align}
\end{mydef}


\begin{mydef}[label=def:dim-depth,breakable]{次元と深さ}
    空でない\hyperref[def:Snglr-C0]{$\bm{C^0}$ 級層状化空間} $(X \to P) \in \Obj{\SnglrC}$ を与える．
    \begin{itemize}
        \item $(X \to P)$ の点 $x \in X$ における\textbf{局所的次元} (local dimension) とは，点 $x$ における $X$ の\hyperref[def:covering-dim]{被覆次元}
        % \footnote{以下の条件を充たす最小の $d \in \mathbb{Z}_{\ge -1}$ のことを（存在すれば）$X$ の\textbf{被覆次元} (covering dimension) と呼ぶ：$X$ の任意の開被覆 $\mathscr{U}$ に対して，十分細かい細分 $\mathscr{V}_{\mathscr{U}} \prec \mathscr{U}$ をとると，任意の互いに異なる $\forall m > d+1$ 個の開集合 $V_1,\, \dots,\, V_{m} \in \mathscr{V}_{\mathscr{U}}$ の共通部分が空になるようにできる．特に，$\emptyset$ の被覆次元は $-1$ と定義する．} 
        $\bm{\dim_x(X)}$ のことを言う．
        \item $(X \to P)$ の\textbf{次元} (dimension) とは
        \begin{align}
            \bm{\dim (X \to P)} \coloneqq \sup_{x \in X} \dim_x (X)
        \end{align}
        のこと．
        \item $(X \xrightarrow{s} P)$ の点 $x \in X$ における\textbf{局所的深さ} (local depth) とは，
        \begin{align}
            \bm{\depth_x(X \to P)} \coloneqq \dim_x (X) - \dim_x (X_{s(x)})
        \end{align}
        のこと．
        \item $(X \to P)$ の\textbf{深さ} (depth) とは，
        \begin{align}
            \bm{\depth(X \to P)} \coloneqq \sup_{x \in X} \depth_x(X \to P)
        \end{align}
        のこと．ただし，$\depth (\emptyset) \coloneqq -1$ と定義する．
    \end{itemize}
\end{mydef}

\begin{myexample}[label=ex:depth-Cone]{コーンの深さ}
    $n$ 次元位相多様体 $Z$ について，定義から $\forall x \in Z$ に対して $\dim_x(Z) = n$ が成り立つ．
    $Z$ を\exref{ex:topomfld}により\hyperref[def:Snglr-C0]{$C^0$ 級層状化空間} $(Z \xrightarrow{s} [0]) \in \Obj{\SnglrC}$ と見做すと，これの\hyperref[def:str-cone]{コーン} $\Cone{Z \xrightarrow{s} [0]}$ について
    \begin{align}
        \depth_x \bigl( \Cone{Z \xrightarrow{s} [0]} \bigr) 
        =
        \begin{cases}
            n+1, &x = \mathrm{pt}, \\
            0, &\text{otherwise}
        \end{cases}
    \end{align}
    であることがわかる．実際 $\Cone{Z}_{\Cone{s}(\mathrm{pt})} = \{\mathrm{pt}\}$ であるが，1点からなる位相空間の\hyperref[def:covering-dim]{被覆次元}は $0$ 次元なので $\dim_{\mathrm{pt}} (\Cone{Z}_{\Cone{s}(\mathrm{pt})}) = 0$ である．
        % ところで， $\dim_{\mathrm{pt}} \bigl( \Cone{Z} \bigr) = n+1$ である．
    一方，コーンポイント以外の点 $x \in \Cone{Z}$ に対して\hyperref[def:stratified-space]{$\Cone{s}(x)$-層}は
    $\Cone{Z}_{\Cone{s}(x)} = \mathbb{R}_{\textcolor{red}{> 0}} \times Z \approx \mathbb{R} \times Z$ であるから，
    $\dim_{x} (\Cone{Z}_{\Cone{s}(x)}) = n + 1$ と計算できる\footnote{さらに，$\forall x \in \Cone{Z}$ に対して $\dim_{x} \Cone{Z} = n+1$ である．}．
    
    　また，$\forall (X \to P) \in \Obj{\SnglrC}$ に対して
    \begin{align}
        \dim \bigl( (\mathbb{R}^m \to [0]) \times (X \to P) \bigr) &= m + \dim \bigl( X \to P \bigr), \\
        \depth \bigl( (\mathbb{R}^m \to [0]) \times (X \to P) \bigr) &= \depth \bigl( X \to P \bigr) 
    \end{align}
    が成り立つ．従って，\hyperref[def:C0-basic]{$C^0$ basic}な $(U \to P_U) = (\mathbb{R}^n \to [0]) \times \Cone{Z \to P} \in \Obj{\SnglrC}$ に対して
    \begin{align}
        \depth (U \to P_U) 
        &= \depth \bigl( \Cone{Z \to P} \bigr)  \\
        &= \dim (Z \to P) + 1
    \end{align}
    が成り立つ．
\end{myexample}

次元と深さに関する帰納法を実行する前に，構成したい $(1,\, 1)$-圏を表す記号の整理をしておこう：
\begin{itemize}
    \item conically smoothチャートの素材となる，\textbf{basic}が成す圏
    \begin{align}
        \Bsc    
    \end{align}
    これは，$C^\infty$ 多様体の圏 $\Mfld$ において $\mathbb{R}^n\; (\forall n \ge -1)$ 全体が成す充満部分圏に相当するものである．
    
    \item 与えられた\hyperref[def:Snglr-C0]{$C^0$ 級層状化空間} $(X \to P) \in \Obj{\SnglrC}$ に対して，その上に入る\textbf{極大アトラス}\footnote{存在するか分からないし，存在したとして一意であるとは限らない．実際，例えば $C^\infty$ 多様体の段階においてさえ $\mathbb{R}^4$ の上の極大アトラス（i.e. $C^\infty$ 構造）は非可算無限個存在する~\cite{taubes1987gauge}．}全体が成す集合を返す\hyperref[def:presheaf]{前層}
    \begin{align}
        \Sm \colon \OP{(\SnglrC)} \lto \SETS
    \end{align}
    この対応が前層であることの直観は，\hyperref[def:strat-emb]{層状化開埋め込み} $f \in \Hom{\SnglrC} \bigl((X \to P),\, (Y \to Q)\bigr)$ が与えられると，$(X \to P)$ 上の極大アトラス $\Sm(X \to P)$ が $(Y \to Q)$ 上の極大アトラス $\Sm(Y \to Q)$ を「制限」する写像 $\Sm(f) \colon \Sm(Y \to Q) \lto \Sm(X \to P)$ によって得られるということである．
    
    \item \hyperref[def:dim-depth]{深さ}が $k$ 以下，かつ\hyperref[def:dim-depth]{次元}が $n$ 以下であるような\hyperref[def:Snglr-C0]{$C^0$ 級層状化空間}全体が成す $\SnglrC$ の充満部分圏を
    \begin{align}
        \SnglrC{}_{\underbrace{\le k}_{\text{depth}},\, \underbrace{\le n}_{\text{dimension}}}
    \end{align}
    と書く．同様に
    \begin{align}
        \Bsc_{\le k,\, \le n},\qquad \Sm_{\le k,\, \le n} \colon \OP{(\SnglrC_{\le k,\, \le n})} \lto \SETS
    \end{align}
    と書く．
    
    \item \textbf{conically smoothな層状化空間}の圏
    \begin{align}
        \Snglr
    \end{align}
    これを作ることが本小節の最終目標である．
\end{itemize}
帰納法により，$\forall k \ge -1$ に対して $\Bsc_{\le k,\, \le \infty}$ および $\Sm_{\le k,\, \le \infty} \colon \OP{(\SnglrC_{\le k,\, \le \infty})} \lto \SETS$ が構成される．

\begin{mydef}[label=def:induction-init]{帰納法の出発点}
    \hyperref[def:Snglr-C0]{\textsf{\textbf{(Snglr-1)}}}より $(\emptyset \to \emptyset) \in \Obj{\SnglrC_{\le -1,\, \le \infty}}$ である．
    \begin{enumerate}
        \item $\Bsc_{\le -1,\, \le \infty} \coloneqq \emptyset$
        \item $\Sm_{\le -1,\, \le \infty}(\emptyset) \coloneqq \{*\}$
    \end{enumerate}
    と定義する．
\end{mydef}

\begin{myassump}[label=assump:induction-Bsc-Sm]{帰納法の仮定}
    与えられた $k \ge -1$ に対して以下の構成が完了していると仮定する：
    \begin{enumerate}
        \item 圏 $\Bsc_{\le k,\, \le \infty}$
        \item 前層 $\Sm_{\le k,\, \le \infty} \colon \OP{(\SnglrC_{\le k,\, \le \infty})} \lto \SETS$
        \item 関手
        \begin{align}
            \mathbb{R} \times (\, \mhyphen \,) \colon \Bsc_{\le k,\, \le \infty} &\lto \Bsc_{\le k,\, \le \infty}, \\
            U &\lmto \mathbb{R} \times U, \\
            \bigl( U \xrightarrow{f} V \bigr) &\lmto \bigl( \mathbb{R} \times U \xrightarrow{\mathrm{id} \times f} \mathbb{R} \times V \bigr) 
        \end{align}
        およびそれが誘導する自然変換\footnote{$X$ の極大アトラス $\Familyset[\big]{U_\alpha,\, \varphi_\alpha}{\alpha \in \Lambda}$ に対して，$\Familyset[\big]{\mathbb{R} \times U_\alpha,\, \mathrm{id} \times \varphi_\alpha}{\alpha \in \Lambda}$ を対応づける．}
        % \begin{align}
        %     \OP{(\SnglrC_{\le k,\, \le \infty})} 
        % \end{align}
        \begin{center}
        \begin{tikzcd}[row sep=large, column sep=large]
            \OP{(\SnglrC_{\le k,\, \le \infty})}  \ar[bend left=50,r, "{\Sm_{\le k,\, \le  \infty}(\, \mhyphen \,)}"{name=U, above}] \ar[bend right=50,r, "{\Sm_{\le k,\, \le  \infty}(\mathbb{R} \times \, \mhyphen \,)}"{name=D, below}] &\SETS
            \ar[Rightarrow, from=U, to=D]
        \end{tikzcd}
        \end{center}
    \end{enumerate}
\end{myassump}

\begin{mydef}[label=def:Bsc-induction,breakable]{圏 $\Bsc_{\le k+1,\, \le \infty}$}
    帰納法の仮定\ref{assump:induction-Bsc-Sm}がある $k \ge -1$ において成立しているとする．
    また，\hyperref[def:C0-basic]{$C^0$ basic}を $U^n_Z \coloneqq (\mathbb{R}^n \to [0]) \times \Cone{Z \to P} \in \Obj{\SnglrC}$ と書く．
    このとき，圏 $\Bsc_{\le k+1,\, \le \infty}$ を以下で定義する：
    \begin{description}
        \item[(\textbf{対象})]　
        
        \hyperref[def:C0-basic]{$C^0$ basic}
        \footnote{\hyperref[def:dim-depth]{depthの定義}から $\depth (Z \to P) \le \dim (Z \to P)$ である．故に\exref{ex:depth-Cone}から，$\depth (Z \to P) \le \dim (Z \to P) = \depth U^n_Z - 1 \le k$ であること，i.e. $(Z \to P) \in \Obj{\SnglrC_{\le k,\, \le \infty}}$ が分かる．} $U^n_Z \in \Obj{\SnglrC_{\le k+1,\, \le \infty}}$ および，極大アトラス $\mathcal{A}_Z \in \Sm_{\textcolor{red}{\le k},\, \le \infty} (Z \to P)$ の組み
        \begin{align}
            (U^n_Z,\, \mathcal{A}_Z)
        \end{align}
        を対象とする．これを\textbf{basic}と呼ぶ．

        \item[(\textbf{射})]　
        
        任意の2つの対象 $(U^n_Z,\, \mathcal{A}_Z),\; (U^m_W,\, \mathcal{A}_W) \in \Obj{\Bsc_{\le k+1,\, \le \infty}}$ に対して，以下の条件を充たす\hyperref[def:strat-emb]{層状化開埋め込み} $f \in \Hom{\SnglrC_{\le k+1,\, \le \infty}} \bigl(U^n_Z,\, U^m_W \bigr)$ を射とする：
        \begin{description}
            \item[\textbf{$\bm{f}$ がコーンポイントを保存しない場合}]　  
            
            ある\hyperref[def:strat-emb]{層状化開埋め込み} $f_0 \in \Hom{\SnglrC_{\le k+1,\, \le \infty}} \bigl(U^n_Z,\; \mathbb{R}^m \times \mathbb{R}_{> 0} \times W\bigr)$ が存在して
            \begin{align}
                f \colon U^n_Z \xrightarrow{f_0} \mathbb{R}^m \times (\mathbb{R}_{> 0} \times W) \hookrightarrow U^m_W = \mathbb{R}^m \times \Cone{W}
            \end{align}
            と書けて，かつ $(U^n_Z,\, f_0) \in \mathcal{A}_{\mathbb{R}^m \times \mathbb{R}_{> 0} \times W} \in \Sm(\mathbb{R}^m \times \mathbb{R}_{> 0} \times W)$
            
            \item[\textbf{$\bm{f}$ がコーンポイントを保存する場合}]　 
            
            $f$ は\hyperref[def:c-smooth-along]{$\mathbb{R}^n$ に沿ってconically smooth}であって，かつ $Df \colon \mathbb{R}^n \times U^n_Z \lto \mathbb{R}^m \times U^m_W$ が単射であり，
            かつ
            \begin{align}
                \mathcal{A}_{f^{-1}(U^m_W \setminus \mathbb{R}^m)} = \Sm_{\le k,\, \le \infty} \bigl( f|_{f^{-1}(U^m_W \setminus \mathbb{R}^m)} \bigr) \bigl(\mathcal{A}_{U^m_W \setminus \mathbb{R}^m}\bigr)
            \end{align}
            を充たす\footnote{ここで帰納法の仮定\ref{assump:induction-Bsc-Sm}-(3) を暗に使っている．}．ただし，$U^m_W \setminus \mathbb{R}^m \coloneqq U^m_W \setminus (\mathbb{R}^m \times \{\mathrm{pt}\}) = \mathbb{R}^{m+1} \times W$ と略記した．
        \end{description}
        
    \end{description}
\end{mydef}

\begin{mydef}[label=def:Sm-induction,breakable]{前層 $\Sm_{\le k+1,\, \le \infty}$}
    帰納法の仮定\ref{assump:induction-Bsc-Sm}がある $k \ge -1$ において成立しているとする．さらに定義\ref{def:Bsc-induction}によって $\Bsc_{\le k+1,\, \le \infty}$ が完成しているとする．
    \begin{itemize}
        \item \hyperref[def:Snglr-C0]{$C^0$ 級層状化空間} $\forall (X \to P) \in \Obj{\SnglrC_{\le k+1,\, \le \infty}}$ に対して，$X \to P$ の\textbf{アトラス} (atlas) を族
        \begin{align}
            \mathcal{A} \coloneqq \Familyset[\Big]{(U_\alpha,\, \varphi_\alpha)  \in \Obj{\Bsc_{\le k+1,\, \le \infty}} \times \Hom{\SnglrC_{\le k+1,\, \le \infty}}\bigl(U_\alpha,\, (X \to P)\bigr) }{\alpha \in \Lambda}
            %  \in \Sm_{\le k+1,\, \le \infty} (X \to P)
        \end{align}
        であって以下の条件を充たすものとして定義する：
        \begin{description}
            \item[\textbf{(Atlas-1)}]　
            
            $\mathcal{A}$ は $(X \to P)$ の開被覆である．
    
            \item[\textbf{(Atlas-2)}]　
            
            $\forall \alpha,\, \beta \in \Lambda$ および $\forall x \in \varphi_\alpha (U_\alpha) \cap \varphi_\beta (U_\beta)$ に対して，
            圏 $\SnglrC_{\le k+1,\, \le \infty}$ の可換図式
            \begin{center}
                % https://tikzcd.yichuanshen.de/#N4Igdg9gJgpgziAXAbVABwnAlgFyxMJZABgBpiBdUkANwEMAbAVxiRAHUQBfU9TXfIRRkAjFVqMWbAKoB9YAB0FjNAAs6XbrxAZseAkRGkx1es1aIQADS189gw+XFmpluYoUAjGDg3dxMFAA5vBEoABmAE4QALZIZCA4EEgAzKaSFiDh8krevpo8EdFxiAlJSEYS5mzZHirqmtQMdN4MAAr8+kIgkVhBqji2WcUV1OWIAExNLTDtnQ6Wvf2D6dWWSvSRalg5ygxqfoXDsaljyZPTrR32BpYMMOErVa4gG3RbqjseeYcUXEA
                \begin{tikzcd}
                \exists \textcolor{red}{W} \arrow[r, hookrightarrow, red,"f_{\beta}"] \arrow[d, hookrightarrow, red,"f_{\alpha}"'] & U_{\beta} \arrow[d, hookrightarrow, "\varphi_{\beta}"] \\
                U_{\alpha} \arrow[r, hookrightarrow, "\varphi_{\alpha}"']         & X                                     
                \end{tikzcd}
            \end{center}
            が存在して $x \in \varphi_\alpha \circ \textcolor{red}{f_\alpha} (\textcolor{red}{W}) = \varphi_\beta \circ \textcolor{red}{f_\beta} (\textcolor{red}{W})$ を充たす．
            ただし，可換図式中の赤色の部分は全て\hyperref[def:Bsc-induction]{圏 $\Bsc_{\le k+1,\, \le \infty}$} の対象および射からなる．

        \end{description}
        アトラス $\mathcal{A}$ の元 $(U_\alpha,\, \varphi_\alpha) \in \mathcal{A}$ のことを\textbf{チャート} (chart) と呼ぶ．

        \item \hyperref[def:Snglr-C0]{$C^0$ 級層状化空間} $\forall (X \to P) \in \Obj{\SnglrC_{\le k+1,\, \le \infty}}$ の2つのアトラス $\mathcal{A},\, \mathcal{B}$ が\textbf{同値}であるとは，$\mathcal{A} \cup \mathcal{B}$ が $(X \to P)$ のアトラスであることを言う．
        これは $(X \to P)$ のアトラス全体の集合の上に同値関係を定める\footnote{同値関係であることの証明は~\cite[Lemma 3.2.11.]{AFT2014stratified}を参照．}．
        $(X \to P)$ の\textbf{極大アトラス} (maximal atlas) とは，この同値関係によるアトラス $\mathcal{A}$ の同値類 $[\mathcal{A}]$ のことを言う．

        \item 前層
        \begin{align}
            \Sm_{\le k+1,\, \le \infty} \colon \OP{(\SnglrC_{\le k+1,\, \le \infty})} \lto \SETS
        \end{align}
        を以下のように定義する：
        \begin{description}
            \item[\textbf{(対象)}]　
            
            任意の\hyperref[def:Snglr-C0]{$C^0$ 級層状化空間} $(X \to P) \in \Obj{\SnglrC_{\le k+1,\, \le \infty}}$ に対して
            \begin{align}
                \Sm_{\le k+1,\, \le \infty} (X \to P) \coloneqq \bigl\{\, [\mathcal{A}] \bigm| \mathcal{A}\; \text{is an atlas of}\; (X \to P) \,\bigr\} 
            \end{align}
            
            \item[\textbf{(射)}]　
            
            任意の\hyperref[def:strat-emb]{層状化開埋め込み} $f \in \Hom{\SnglrC_{\le k+1,\, \le \infty}}\bigl((X \to P),\, (Y \to Q)\bigr)$ に対して，$f$ によるアトラスの引き戻しを対応付ける．
        \end{description}
    \end{itemize}
\end{mydef}

以上の帰納法をまとめて，conically smoothな層状化空間と\hyperref[def:strat-emb]{層状化開埋め込み}の圏 $\Snglr$ を得る．

\begin{mydef}[label=def:Snglr]{圏 $\Snglr$}
    \begin{itemize}
        \item \hyperref[def:Bsc-induction]{\textbf{basic}}のなす圏 $\Bsc$ を以下で定義する：
        \begin{align}
            \Bsc \coloneqq \bigcup_{k \ge -1} \Bsc_{\le k,\, \le \infty}
        \end{align}
        \item 極大アトラスの集合を与える関手 $\Sm \colon \OP{(\SnglrC)} \lto \SETS$ を以下の\hyperref[def:Kanext]{右Kan拡張}として定義する：
        \begin{center}
            % https://tikzcd.yichuanshen.de/#N4Igdg9gJgpgziAXAbVABwnAlgFyxMJZABgBpiBdUkANwEMAbAVxiRAB12B5ABWAApOAZTABzBgCcAwgH1gAHgAEnLGABmOAJ6lOpZewYx9qjZoC+ASjMgzpdJlz5CKMgEYqtRizadeA4WKSUlY2diAY2HgERK7kHvTMrIgc7EIAogAqQjYeMFCi8ESgahIQALZIZCA4EEgATNQJ3snCZXJKKupaOux6nIbGXeahxaUViFU1SK62o+XT1FOIDSAMqkkgUHRwABZ5II1eG60Hq3QARjAMPA5RziASWKI7ODlmQA
            \begin{tikzcd}[column sep=large]
            {\OP{(\SnglrC_{< \infty,\, \le \infty})}} \arrow[r, "{\Sm_{< \infty,\, \le \infty}}"] \arrow[d,hookrightarrow] & \SETS \\
            \OP{(\SnglrC)} \arrow[ru,red, "\Sm"', dashed]                                                       &      
            \end{tikzcd}
        \end{center}
        ただし，$\SnglrC_{< \infty,\, \le \infty} \coloneqq \bigcup_{k \ge -1} \SnglrC_{\le k,\, \le \infty}$ とおいた．
        \item \textbf{conically smoothな層状化空間} (conically smooth stratified space) と\hyperref[def:strat-emb]{層状化開埋め込み}の圏 $\Snglr$ を以下で定義する：
        
        \begin{description}
            \item[\textbf{(対象)}]　
            
            \hyperref[def:Snglr-C0]{$C^0$ 級層状化空間} $(X \to P) \in \Obj{\SnglrC}$ およびその\hyperref[def:Sm-induction]{極大アトラス} $\mathcal{A}_X \in \Sm (X \to P)$ の組み $\bigl((X\to P),\, \mathcal{A}_X\bigr)$ を対象とする．
        
            \item[\textbf{(射)}]　
            
            \hyperref[def:strat-emb]{層状化開埋め込み} $f \in \Hom{\SnglrC} \bigl((X \to P),\, (Y \to Q)\bigr)$ であって，$f^* \mathcal{A}_Y = \mathcal{A}_X$ を充たすものを射とする．
        \end{description}
    \end{itemize}
    
\end{mydef}

\subsection{conically smooth map}

ここまでは\hyperref[def:strat-emb]{層状化開埋め込み}のみを考えていたため，一般の\hyperref[def:stratified-space]{層状化写像}のconically smoothnessを定義しなくてはいけない．

\begin{mydef}[label=def:c-smooth-map]{conically smooth map}
    2つの\hyperref[def:Snglr]{\underline{\textbf{basic}}}\footnote{\hyperref[def:C0-basic]{$C^0$ basic}を $U^n_Z \coloneqq (\mathbb{R}^n \to [0]) \times \Cone{Z \to P} \in \Obj{\SnglrC}$ と書く．} $X = (U^n_Z,\, \mathcal{A}_Z),\, Y = (U^m_W,\, \mathcal{A}_W)\in \Obj{\Bsc}$ の間の\hyperref[def:stratified-space]{層状化写像} $f \colon U^n_Z \lto U^m_W$ が\textbf{conically smooth}であることを，
    \hyperref[def:dim-depth]{$\depth (Y)$} に関する帰納法によって定義する：
    \begin{enumerate}
        \item まず，$\depth (Y) = -1$ のときは $X = Y = \emptyset$ であり，一意的に定まる $X,\, Y$ 間の層状化写像がconically smoothであると定義する．
        \item 深さ $k \ge -1$ のbasicに対して定義が完了しているとする．
        $Y \in \Obj{\Bsc}$ の深さが高々 $k+1$ であるとき，層状化写像 $f \colon X \lto Y$ がconically smoothであることを以下で定義する：
        \begin{description}
            \item[\textbf{$\bm{f}$ がコーンポイントを保存しない場合}]　  
            
            あるconically smoothな層状化写像 $f_0 \colon X \lto \mathbb{R}^m \times \mathbb{R}_{> 0} \times W$ が存在して
            \begin{align}
                f \colon X \xrightarrow{f_0} \mathbb{R}^m \times (\mathbb{R}_{> 0} \times W) \hookrightarrow Y = \mathbb{R}^m \times \Cone{W}
            \end{align}
            と書ける\footnote{\exref{ex:depth-Cone}より$\depth(W) < k+1$であり，帰納法の仮定が使える．}．
            
            \item[\textbf{$\bm{f}$ がコーンポイントを保存する場合}]　 
            
            $f$ は\hyperref[def:c-smooth-along]{$\mathbb{R}^n$ に沿ってconically smooth}であって，かつ制限
            \begin{align}
                f|_{f^{-1}(Y \setminus \mathbb{R}^m)} \colon f^{-1}(Y \setminus \mathbb{R}^m) \lto Y \setminus \mathbb{R}^m
            \end{align}
            がconically smooth．ただし，$U^m_W \setminus \mathbb{R}^m \coloneqq U^m_W \setminus (\mathbb{R}^m \times \{\mathrm{pt}\}) = \mathbb{R}^{m+1} \times W$ と略記した．
        \end{description}
    \end{enumerate}
    
    \tcblower
    % \hyperref[def:C0-basic]{$C^0$ basic}を $U^n_Z \coloneqq (\mathbb{R}^n \to [0]) \times \Cone{Z \to P} \in \Obj{\SnglrC}$ と書く．
    \hyperref[def:Snglr]{conically smoothな層状化空間} $\bigl( (X \to P),\, \mathcal{A}_X \bigr),\, \bigl( (Y \to Q),\, \mathcal{A}_Y \bigr) \in \Obj{\Snglr}$ の間の\hyperref[def:stratified-space]{層状化写像} $f \colon (X \to P) \lto (Y \to Q)$ が\textbf{conically smooth}であるとは，
    任意のチャートの組み合わせ $(U,\, \varphi) \in \mathcal{A}_X,\, (V,\, \psi) \in \mathcal{A}_Y$ に対して
    \begin{align}
        \psi^{-1} \circ f \circ \varphi \colon U \lto V
    \end{align}
    がconically smooth (for \hyperref[def:Snglr]{basic}s) であることを言う．
\end{mydef}

\begin{myprop}[label=prop:c-smooth-map]{conically smooth mapの基本性質}
    2つの\hyperref[def:c-smooth-map]{conically smooth map}の合成もconically smoothである．
\end{myprop}

\begin{proof}
    ~\cite[Proposition 3.3.5]{AFT2014stratified}
\end{proof}

命題\ref{prop:c-smooth-map}より，\textbf{conically smoothな層状化空間}の圏を定義できる．

\begin{mydef}[label=def:c-smooth]{conically smoothな層状化空間の圏 $\Strat$ }
    \textbf{conically smoothな層状化空間}の圏 $\Strat$ を以下で定義する：
    \begin{description}
        \item[\textbf{(対象)}]　
        
        圏 \hyperref[def:Snglr]{$\Snglr$}と全く同じ対象を持つ：
        \begin{align}
            \Obj{\Strat} \coloneqq \Obj{\Snglr}
        \end{align}
        \item[\textbf{(射)}]　
        
        \hyperref[def:c-smooth-map]{conically smooth map}を射とする．
    \end{description}
\end{mydef}

定義から明らかに $\Snglr \subset \Strat$ である．
ここで，圏 $\Strat$ における特別な射に名前をつけておこう：

\begin{mydef}[label=def:cbl,breakable]{constructuble bundle}
    \begin{itemize}
        \item \hyperref[def:c-smooth-map]{conically smoothな層状化写像} $\pi \in \Hom{\Strat}\bigl( (E \to P),\, (B \to Q) \bigr)$ が\textbf{層状化ファイバー束} (conically smooth fiber bundle) であるとは，conically smoothな\hyperref[def:strat-emb]{層状化開埋め込み}の族 $\Familyset[\big]{U_\alpha \hookrightarrow B}{\alpha \in \Lambda},\; \Familyset[\big]{\varphi_\alpha \colon U_\alpha \times F_\alpha \hookrightarrow E}{\alpha \in \Lambda}$ が存在して以下を充たすことを言う：
        \begin{description}
            \item[\textbf{(Bun-1)}]　
            
            $\forall \alpha \in \Lambda$ に対して，圏 $\Strat$ における\hyperref[def:pullback-pushout]{引き戻し}の図式
            
            \begin{center}
                % https://tikzcd.yichuanshen.de/#N4Igdg9gJgpgziAXAbVABwnAlgFyxMJZABgBpiBdUkANwEMAbAVxiRAFUB9YAHR8bQALOgF8ABHzwBbeGIBi3PgOEiQI0uky58hFGQCMVWoxZsuvfgyGi1GkBmx4CRfaUPV6zVohAAhW5qOOi7kRp6mPgCiakYwUADm8ESgAGYAThBSSGQgOBBIAMzUOHRYDGyCEBAA1iAeJt4gfPRpQliKltaq6qkZWYiuufmIAEzFpeU+lTUBIOmZ2cXDg+GNfFJ0OIJpUsBoGQBWIpz6dSAMdABGMAwAClpOuiBpWPGCOLPz-UVDSGPGXjYfDQWBiIiAA
                \begin{tikzcd}
                U_{\alpha} \times F_{\alpha} \arrow[dr, phantom," " {pullback}, very near start, color=black] \arrow[r, "\varphi_{\alpha}", hook] \arrow[d, "\mathrm{proj}_1"'] & E \arrow[d, "\pi"] \\
                U_{\alpha} \arrow[r, hook]                                                                     & B                 
                \end{tikzcd}
            \end{center}
            が成り立つ．

            \item[\textbf{(Bun-2)}]　
            
            族 $\Familyset[\big]{U_\alpha}{\alpha \in \Lambda}$ は $B$ の開基である．
        \end{description}
        
        \item \hyperref[def:c-smooth-map]{conically smoothな層状化写像} $\pi \in \Hom{\Strat}\bigl( (E \to P),\, (B \to Q) \bigr)$ が\textbf{弱構成可能束} (weakly constructuble bundle) であるとは，$\forall q \in Q$ に対して，$\pi$ の\hyperref[def:stratified-space]{$q$-層}への制限
        \begin{align}
            \pi|_{\pi^{-1}(B_q)} \colon \pi^{-1}(B_q) \lto B_q
        \end{align}
        が層状化ファイバー束であることを言う．

        \item  \hyperref[def:c-smooth-map]{conically smoothな層状化写像} $\pi \in \Hom{\Strat}\bigl( (E \to P),\, (B \to Q) \bigr)$ が\textbf{構成可能束} (constructuble bundle) であることを，\hyperref[def:dim-depth]{$\depth (E)$} に関する帰納法によって定義する：
        \begin{enumerate}
            \item $\depth (E) = 0$ のとき，$\pi$ が構成可能束であるとは，$\pi$ が $C^\infty$ ファイバー束であることを言う．
            \item 深さ $k \ge 0$ までの定義が完了しているとする．$\depth(E) \le k+1$ のとき，$\pi$ が構成可能束であるとは，以下の2条件を充たすことを言う：
            \begin{description}
                \item[\textbf{(cBun-1)}] $\pi$ は弱構成可能束である．
                \item[\textbf{(cBun-2)}] $\forall q \in Q$ に対して，$\pi$ が誘導する層状化写像
                \begin{align}
                    \Link_{\pi^{-1}(B_q)} (E) \lto \pi^{-1}(B_q) \times_{B_q} \Link_{B_q} (B)
                \end{align}
                が構成可能束である．
            \end{description}
        \end{enumerate}
    \end{itemize}
\end{mydef}


\subsection{管状近傍・ハンドル分解}

\section{層状化空間の接構造}

\subsection{$\mathrm{Kan}$-豊穣化}

圏 $\Kan$ を，
\begin{itemize}
    \item \hyperref[def:KanCplx]{Kan複体}を対象に持つ
    \item Kan複体の間の自然変換を射に持つ
\end{itemize}
$(1,\, 1)$-圏とする．$\Kan$ は\hyperref[def:SimpSet]{単体的集合の圏 $\sSet$}の充満部分圏であり，直積\eqref{eq:sSet-tensor}をテンソル積とするモノイダル圏になる．

\begin{mydef}[label=def:standard-cosimplicial-mfd]{余単体的多様体}
    以下で定義する関手
    \begin{align}
        \bm{\Delta_e} \colon \Delta \lto \Strat
    \end{align}
    のことを\textbf{余単体的多様体} (standard cosimplicial manifold) と呼ぶ\footnote{\hyperref[def:simplicial-top]{幾何学的 $n$-単体}に似ているが，$x^i \ge 0$ の領域で切り取っていない．}：
    \begin{itemize}
        \item $[n] \in \Obj{\Delta}$ を，\hyperref[def:c-smooth]{conically smoothな層状化空間}
        \begin{align}
            \Delta_e^n \coloneqq \bigl\{\, (x^0,\, \dots,\, x^n) \in \mathbb{R}^{n+1} \bigm| \sum_{i=0}^n x^i = 1 \,\bigr\} 
        \end{align}
        に対応付ける．
        \item $\alpha \in \Hom{\Delta}([m],\, [n])$ を，\hyperref[def:c-smooth-map]{conically smoothな層状化写像}
        \begin{align}
            \Delta_e(\alpha) \colon \Delta_e^m &\lto \Delta_e^n, \\
            (x^0,\, \dots,\, x^m) &\lmto \Bigl( \sum_{j,\, \alpha(j) = 0} x^j,\, \dots,\, \sum_{j,\, \alpha(j) = n} x^j \Bigr)
        \end{align}
        に対応付ける．
    \end{itemize}
\end{mydef}

$\PshSETS{\Strat}$ から $\sSet$ への関手を
\begin{align}
    (\, \mhyphen \,)|_{\Delta} \colon \PshSETS{\Strat} &\lto \sSet, \\
    F &\lmto F \circ \Delta_e
\end{align}
で定義する．さらに，$\forall X,\, Y \in \Obj{\Strat}$ に対して前層 $\widetilde{\Hom{\STRAT}} (X,\, Y) \in \PshSETS{\Strat}$ を
\begin{align}
    \widetilde{\Hom{\STRAT}} (X,\, Y) \colon \OP{\Strat} &\lto \SETS, \\
    \textcolor{red}{Z} &\lmto \bigl\{\, f \in \Hom{\Strat}(\textcolor{red}{Z} \times X,\, \textcolor{red}{Z} \times Y) \bigm| \mathrm{proj}_Z \circ f = \mathrm{proj}_Z \,\bigr\}, \\
    (Z \xrightarrow{\alpha} W) &\lmto \Biggl(\substack{\widetilde{\Hom{\STRAT}(X,\, Y)}(Z) \lto \widetilde{\Hom{\STRAT}(X,\, Y)}(W), \\ f \lmto \biggl((w,\, x) \mapsto \bigl(w,\, \mathrm{proj}_Y \circ f (\alpha(w),\, x)\bigr)\biggr)}\Biggr)
\end{align}
で定義する．ただし，\hyperref[def:c-smooth-map]{conically smoothな層状化写像} $\mathrm{proj}_Z \in \Hom{\Strat}(Z \times X,\, Z)$ とは第一成分への射影のことである．
同様にして前層 $\widetilde{\Hom{\SNGLR}} (X,\, Y) \in \PshSETS{\Strat}$ を
\begin{align}
    \widetilde{\Hom{\SNGLR}} (X,\, Y) \colon \OP{\Strat} &\lto \SETS, \\
    \textcolor{red}{Z} &\lmto \bigl\{\, f \in \Hom{\Snglr}(\textcolor{red}{Z} \times X,\, \textcolor{red}{Z} \times Y) \bigm| \mathrm{proj}_Z \circ f = \mathrm{proj}_Z \,\bigr\}, \\
    (Z \xrightarrow{\alpha} W) &\lmto \Biggl(\substack{\widetilde{\Hom{\SNGLR}(X,\, Y)}(Z) \lto \widetilde{\Hom{\SNGLR}(X,\, Y)}(W), \\ f \lmto \biggl((w,\, x) \mapsto \bigl(w,\, \mathrm{proj}_Y \circ f (\alpha(w),\, x)\bigr)\biggr)}\Biggr)
\end{align}
で定義する．

\begin{mylem}[label=lem:Kan-enriched]{}
    $\forall X,\, Y \in \Obj{\Strat}$ に対して定まる\hyperref[def:SimpSet]{単体的集合}
    \begin{align}
        \Hom{\STRAT}(X,\, Y) &\coloneqq \eval{\widetilde{\Hom{\STRAT}}(X,\, Y)}_{\Delta}, \\
        \Hom{\SNGLR}(X,\, Y) &\coloneqq \eval{\widetilde{\Hom{\SNGLR}}(X,\, Y)}_{\Delta}
    \end{align}
    は\hyperref[def:KanCplx]{Kan複体}である．
\end{mylem}

\begin{proof}
    ~\cite[Lemma 4.1.4.]{AFT2014stratified}．
\end{proof}

\begin{mydef}[label=def:Strat-infty]{{$(\infty,\, 1)$}-圏 {$\STRAT,\, \SNGLR,\, \BSC$}}
    \hyperref[def:enriched]{$\Kan$-豊穣圏} $\STRAT$ を以下で定義する：
    \begin{itemize}
        \item $\Obj{\STRAT} \coloneqq \Obj{\Snglr}$
        \item 補題\ref{lem:Kan-enriched}で構成した $\Hom{\STRAT}(X,\, Y) \in \Obj{\Kan}$ をHom対象とする．
    \end{itemize}
    同様に，$\Kan$-豊穣圏 $\SNGLR$ を以下で定義する：
    \begin{itemize}
        \item $\Obj{\SNGLR} \coloneqq \Obj{\Snglr}$
        \item 補題\ref{lem:Kan-enriched}で構成した $\Hom{\SNGLR}(X,\, Y) \in \Obj{\Kan}$ をHom対象とする．
    \end{itemize}
    
    \tcblower

    $\Kan$-豊穣圏 $\SNGLR$ の対象を \hyperref[def:Snglr]{$\Obj{\Bsc}$} に制限して得られる充満部分圏を $\BSC$ と書く．
    % \tcblower

    % \hyperref[def:infinity-1]{$(\infty,\, 1)$-圏} $\STRAT,\, \SNGLR$ を以下で定義する：
    % \begin{align}
    %     \STRAT &\coloneqq \hcNer  (\STRAT), \\
    %     \SNGLR &\coloneqq \hcNer  (\SNGLR)
    % \end{align}
    
\end{mydef}

\begin{marker}
    % homotopy hypothesisより $\Kan$-豊穣圏は $(\infty,\, 1)$-圏のモデルと見做せる．故に
    $\Kan$-豊穣圏を\hyperref[def:nerve-hc]{homotopy coherent nerve functor} $\hcNer  \colon \sCat \lto \sSet$ で単体的集合の圏 $\sSet$ へ埋め込んだものは $(\infty,\, 1)$-圏である~\cite[Proposition 1.1.5.10.]{lurie2008higher}．
    故に以下では $\Kan$-豊穣圏 $\STRAT,\, \SNGLR$ と $(\infty,\, 1)$-圏 $\hcNer  (\STRAT),\, \hcNer  (\SNGLR)$ を区別しない．
\end{marker}

\subsection{$(\infty,\, 1)$-圏におけるファイブレーション}

\begin{mydef}[label=def:infty-fib]{{$(\infty,\, 1)$}-ファイブレーション}
    $p \colon \Cat{E} \lto \Cat{B}$ を\hyperref[def:infinity-1]{$(\infty,\, 1)$-圏の関手}とする．
    \begin{description}
        \item[\textbf{(lifting property)}]　
        
        包含 $\iota \in \Hom{\sSet}(\Lambda^n_j,\, \Delta^n)$ に関して $p \circ \textcolor{blue}{f_0} = \textcolor{blue}{f} \circ \iota$ を充たす任意の $(\textcolor{blue}{f_0},\, \textcolor{blue}{f}) \in \Hom{\sSet}(\Lambda^n_j,\, \Cat{E}) \times \Hom{\sSet}(\Delta^n,\, \Cat{B})$ に対して，以下の図式を可換にする $\textcolor{red}{\bar{f}} \in \Hom{\sSet}(\Delta^n,\, \Cat{E})$ が存在する：
        \begin{center}
            % https://tikzcd.yichuanshen.de/#N4Igdg9gJgpgziAXAbVABwnAlgFyxMJZABgBpiBdUkANwEMAbAVxiRAB12AZOgWwCModAHpoA+gCsQAX1LpMufIRRkAjFVqMWbTgBEYDHCLQy5IDNjwEiq0uur1mrRB3YBhOjmAAhaafmWSjbkGo7aLpweXgCiftIaMFAA5vBEoABmAE4QvEhkIDgQSADMDlrOIEn+IFk5SLYFRYgATGVObOkg1Ax0-AYACgpWyiCZWEkAFjjVtbmIpY1IrZrtLiayGdlz+YX11EZYDGwTEBAA1l0r4a74Rpc9fQyDgdYuY5PTGzVbe4vz3VgwBUhHAJolLmEKpx+HRMsB0nEKNIgA
            \begin{tikzcd}[column sep=large,row sep=large]
            \Lambda^n_j \arrow[r, blue, "\forall f_0"] \arrow[d, "\iota"', hook]   & \Cat{E} \arrow[d, "p"] \\
            \Delta^n \arrow[r, blue, "\forall f"'] \arrow[ru, red, "\exists \bar{f}", dashed] & \Cat{B}               
            \end{tikzcd}
        \end{center}
    \end{description}
    \tcblower
    \begin{itemize}
        \item $p$ が\textbf{内的ファイブレーション} (inner fibration) であるとは，$0 \,\textcolor{red}{<}\, \forall j \textcolor{red}{<} \forall n$ に対して \textsf{\textbf{(lifting property)}} を充たすことを言う．
        \item $p$ が\textbf{右ファイブレーション} (right fibration) であるとは，$0 \,\textcolor{red}{<}\, \forall j \le \forall n$ に対して \textsf{\textbf{(lifting property)}} を充たすことを言う．
        \item $p$ が\textbf{左ファイブレーション} (left fibration) であるとは，$0 \le \forall j \,\textcolor{red}{<}\, \forall n$ に対して \textsf{\textbf{(lifting property)}} を充たすことを言う．
        \item $p$ が\textbf{Kanファイブレーション} (Kan fibration) であるとは，$0 \le \forall j \le \forall n$ に対して \textsf{\textbf{(lifting property)}} を充たすことを言う．
    \end{itemize}
\end{mydef}

系\ref{col:horn-coeq}によると，\hyperref[def:infty-fib]{\textsf{\textbf{(lifting property)}}}は，\underline{$(\infty,\, 1)$-圏 $\Cat{B}$ における}\hyperref[def:horn]{角}の図式
$\bigl(p_{[n-1]}(\textcolor{blue}{f_0}{}_0),\, \dots,\, \underbrace{\bullet}_{j},\, \dots,\, p_{[n-1]}(\textcolor{blue}{f_0}{}_n)\bigr) \in (\Cat{B}_{n-1})^{\times n}$ を $n$-射 $\textcolor{blue}{f} \in \Cat{B}_n$ が埋めているならば，
\underline{$(\infty,\, 1)$-圏 $\Cat{E}$ における}\hyperref[def:horn]{角}の図式
$\bigl(\textcolor{blue}{f_{0}}{}_0,\, \dots,\, \underbrace{\bullet}_{j},\, \dots,\, \textcolor{blue}{f_0}{}_n\bigr) \in (\Cat{E}_{n-1})^{\times n}$ を埋める $n$-射 $\textcolor{red}{\bar{f}} \in \Cat{E}_n$ が存在することを主張している．

\begin{mydef}[label=def:isofibration]{アイソファイブレーション}
    $p \colon \Cat{E} \lto \Cat{B}$ を\hyperref[def:infty-fib]{内的ファイブレーション}とする．
    $p$ が\textbf{アイソファイブレーション} (isofibration) であるとは，
    $\Cat{E}$ における任意の対象 $x \in \Cat{E}_0$ および $\Cat{B}$ における任意の\hyperref[def:isom-infty]{同型射} $\bigl(y \xrightarrow{f} p_{[0]}(x)\bigr) \in \mathcal{B}_1$ に対して，$\Cat{E}$ の同型射 $\bigl( \bar{y} \xrightarrow{\bar{f}} x\bigr) \in \mathcal{E}_1$ が存在して $p_{[1]}(\bar{f}) = f$ を充たすこと\footnote{$\Cat{B}$ の同型射が $\Cat{E}$ の同型射に持ち上がるということ．}．
\end{mydef}

一般に，定義\ref{def:triv-Kan-fib}, \ref{def:infty-fib}, \ref{def:isofibration}の間には次の関係がある~\cite[\href{https://kerodon.net/tag/01GP}{Tag 01GP}]{kerodon}：
\begin{figure}[H]
    \centering
    % https://tikzcd.yichuanshen.de/#N4Igdg9gJgpgziAXAbVABwnAlgFyxMJZARgBoAGAXVJADcBDAGwFcYkQAdDnGADx2BcAFgE80MAE4SYAM2SwZiHBKy0AtAGl6YNTKwAjSsGWqsTAARaw5vfon08BAL5OQT0uky58hFGWLUdEys7Fw8-IIcouJSsvKyiFhgMjgiugZGVjYG9o5gLm4eIBjYeUQATKTlgQwsbIic3HwCwmKS0nIKicmp6YbAKgDmQjjZdg4+Be6epT5E5FU1wfWN4S1RbbGdCUkpabZGjLKjtrmTrtPFXmV+pADMS3WhTRGtMR3xilhwEKcTBEZvr8cv98hciiVvAQiGQACyPEINMLNSLRdpxLq7XoHYBJMCSMZnZwXQIwKCDeBEUAyCQQAC2SAWIBwECQZBAjCSKygEGY+iOIBoQhg9CgSDAzEYjBoOHoWEY7Eg+MK1NpDMQ7JZSDuNE5yoaPL5AqFIrFiAlUplcoVDSVbEuNPpbJlrMQlQ5XPYhv5bBNovFkulzOtioI9qKjvVOuZrthus9Bt5PsFIGF-vNgat8tDyodaqQ7q1iDjHv1IG9xtTpoDluD2dtYZVIEjSBLRYArPGyxXfVX0xag7L6+BG05KE4gA
    \begin{tikzcd}
                                                                            & {\text{\hyperref[def:triv-Kan-fib]{trivial Kan fibration}}} \arrow[d, Rightarrow]              &                                                                           \\
                                                                            & {\text{\hyperref[def:infty-fib]{Kan fibration}}} \arrow[ld, Rightarrow] \arrow[rd, Rightarrow] &                                                                           \\
    {\text{\hyperref[def:infty-fib]{left fibration}}} \arrow[rd, Rightarrow] &                                                                                                & {\text{\hyperref[def:infty-fib]{right fibration}}} \arrow[ld, Rightarrow] \\
                                                                            & {\text{\hyperref[def:isofibration]{isofibration}}} \arrow[d, Rightarrow]                       &                                                                           \\
                                                                            & {\text{\hyperref[def:infty-fib]{inner fibration}}}                                             &                                                                          
    \end{tikzcd}
    \caption{$(\infty,\, 1)$-圏における種々のファイブレーションの関係．$A \IMP B$ は $A$ ならば $B$ を意味する．}
    \label{fig:infty-fib}
\end{figure}%

\begin{mydef}[label=def:fullsub-infty,breakable]{充満部分{$(\infty,\, 1)$-圏}}
    \begin{itemize}
        \item \hyperref[def:infinity-1]{$(\infty,\, 1)$-圏} $\Cat{C}$ の\textbf{部分 $\bm{(\infty,\, 1)}$-圏} (sub $(\infty,\, 1)$-category) とは，
        \hyperref[def:SimpSet]{単体的部分集合} $\Cat{S} \subset \Cat{C}$ であって，その\hyperref[def:SimpSet]{包含写像} $i \colon \Cat{S} \hookrightarrow \Cat{C}$ が\hyperref[def:infy-fib]{内的ファイブレーション}であるようなもののこと\footnote{このとき $\Cat{S}$ は $(\infty,\, 1)$-圏になる~\cite[\href{https://kerodon.net/tag/01CG}{Tag 01CG}]{kerodon}}．
        \item 部分 $(\infty,\, 1)$-圏 $\Cat{S} \subset \Cat{C}$ が\textbf{充満部分 $\bm{(\infty,\, 1)}$-圏} (full sub $(\infty,\, 1)$-category) であるとは，
        $\forall n \ge 0$ に対して以下の条件を充たすことを言う：
        \begin{description}
            \item[\textbf{(fullsub)}] 
            $\forall \sigma \in \Cat{C}_n \cong \Hom{\sSet} (\Delta^n,\, \Cat{C})$ に対して，
            \begin{align}
                \sigma_{[0]}(\Delta^n_0) \subset \Cat{S}_0 \IMP \sigma \in \Cat{S}_n
            \end{align}
            が成り立つ
            \footnote{すなわち，$\Cat{C}$ の勝手な $n$-射 $\sigma \in \Cat{C}_n$ を図\ref{fig:1-simp}の方法で図示した際に，そこに出現する全ての頂点が $\Cat{S}_0$ の元であるならば $\sigma \in \Cat{S}_n$ である．}．
        \end{description}
    \end{itemize}
\end{mydef}

2つの\hyperref[def:infty-fib]{右ファイブレーション} $\mathcal{E} \xrightarrow{p} \mathcal{B},\; \mathcal{E}' \xrightarrow{p'} \mathcal{B}$ が与えられたとき，これらの間の射とは集合
\begin{align}
    \Hom{\Rfib_{\Cat{B}}} (\Cat{E},\, \Cat{E}') &\coloneqq 
    \bigl\{\, f \in \Hom{\sSet}(\mathcal{E},\, \mathcal{E}') \bigm| p' \circ f = p \,\bigr\}
\end{align}
のことである．$\Hom{\Rfib_{\Cat{B}}} (\Cat{E},\, \Cat{E}')$ の元を $\sSet$ における可換図式として表すと以下の通り：
\begin{center}
    \begin{tikzcd}
        \Cat{E} \arrow[rd, "p"'] \arrow[rr, "f"] &         & \Cat{E}' \arrow[ld, "p'"] \\
                                            & \Cat{B} &                          
    \end{tikzcd}
\end{center}
$\Hom{\Rfib_{\Cat{B}}} (\Cat{E},\, \Cat{E}') \subset \Hom{\sSet}(\mathcal{E},\, \mathcal{E}')$ を\exref{def:Kan}の方法で\hyperref[def:SimpSet]{単体的集合}と見做せる．
このようにして得られる単体的集合 $\Hom{\Rfib_{\Cat{B}}} (\Cat{E},\, \Cat{E}')$ の最大の\hyperref[def:fullsub-infty]{部分Kan複体}を $\bm{\Hom{\RFIB_{\Cat{B}}}(\Cat{E},\, \Cat{E}')}$ と書く．

\begin{mydef}[label=def:RFIB]{右ファイブレーションの成す {$(\infty,\, 1)$}-圏}
    $\Cat{B}$ を $(\infty,\, 1)$-圏とする．
    $\Kan$-豊穣圏 $\RFIB_{\Cat{B}}$ を
    \begin{itemize}
        \item \hyperref[def:infty-fib]{右ファイブレーション}を対象とする
        \item $\Hom{\RFIB_{\Cat{B}}}(\Cat{E},\, \Cat{E}')$ をHom対象とする
    \end{itemize}
    ことで定義する．
    以降では $(\infty,\, 1)$-圏 $\hcNer (\RFIB_{\Cat{B}}) \in \Obj{\sSet}$ のことも $\RFIB_{\Cat{B}}$ と書き，区別しない．
\end{mydef}

\subsection{{$(\infty,\, 1)$}-圏におけるスライス圏}

\begin{mydef}[label=def:Simp-Join]{単体的集合のjoin}
    2つの単体的集合 $S,\, T \in \Obj{\sSet}$ の\textbf{join}とは，単体的集合
    \begin{align}
        \bm{S \star T} \colon \OP{\Delta} &\lto \SETS, \\
        [n] &\lmto \coprod_{[i];\, -1 \le i \le n} (S_i \times T_{n-i-1}), \\
        \biggl( [m] \xrightarrow{\alpha} [n] \biggr) &\lmto \biggl( \Bigl( [i];\, (x,\, y) \Bigr) \mapsto \Bigl( \alpha^{-1}([i]);\, \bigl(S(\alpha|_{\alpha^{-1}([i])})(x),\, T(\alpha|_{\alpha^{-1}([m]\setminus[i])})(y) \bigr)  \Bigr)  \biggr) 
    \end{align}
    のこと．ただし $S_{-1} = T_{-1} \coloneqq \{*\},\; [-1] \coloneqq \emptyset$ とおいた．
\end{mydef}
$d^n_j \in \Hom{\OP{\Delta}}([n],\, [n-1])$ に対して
\begin{align}
    (d^n_j)^{-1} ([i]) &= 
    \begin{cases}
        [i], &-1 \le i < j \\
        [i-1], &j \le i \le n
    \end{cases} 
    \\
    (d^n_j)^{-1} ([n]\setminus [i]) &= \begin{cases}
        [n-1]\setminus [i], &-1 \le i < j \\
        [n-1]\setminus [i-1], &j \le i \le n
    \end{cases}
\end{align}
であるから，
$S \star T$ の\hyperref[def:SimpSet]{面写像}は $n \ge 1,\, 0 \le j \le n$ に対して
\begin{align}
    \label{eq:face-join}
    \partial^n_j \colon \coprod_{-1 \le i \le n} (S_i \times T_{n-i-1}) &\lto \coprod_{-1 \le i \le n-1} (S_i \times T_{n-i-2}), \\
    \bigl([i];\, (x,\, y) \bigr) &\lmto 
    \begin{cases}
        % \bigl(-1;\,y \bigr), &i=0 \\
        \bigl([-1];\, (*,\, \partial^n_j y) \bigr), &i=-1 \\
        \bigl([i];\, (x,\, \partial^{n-i-1}_{j-i-1} y) \bigr), &0 \le i < j,\; (i,\, j) \neq (n-1,\, n) \\
        \bigl([i-1];\, (\partial^i_j x,\, y) \bigr), &j \le i \le n-1,\; (i,\, j) \neq (0,\, 0) \\
        \bigl([n-1];\, (\partial^n_j x,\, *) \bigr), &i=n \\
        \bigl([n-1];\, (x,\, *) \bigr), &(i,\, j) = (n-1,\, n) \\
        \bigl([-1];\, (*,\, y) \bigr), &(i,\, j) = (0,\, 0)
        % \bigl(n;\, x\bigr), &j > i=0
        % \bigl(-1;\,\partial^n_j y \bigr), &i=-1
    \end{cases}
\end{align}
% であり，\hyperref[def:SimpSet]{縮退写像}は
% \begin{align}
%     \label{eq:degen-join}
%     \sigma^n_j \colon \coprod_{-1 \le i \le n} (S_i \times T_{n-i-1}) &\lto \coprod_{-1 \le i \le n+1} (S_i \times T_{n-i}), \\
%     \bigl( i;\, (x,\, y) \bigr) &\lmto 
%     \begin{cases}
%         % \bigl(n;\,\sigma^n_j x \bigr), &i=n \\
%         \bigl(i;\, (\sigma_j^i x,\, y) \bigr), &j < i,\; i\neq 0 \\
%         \bigl(i+1;\, (\sigma_{j}^i x,\, y) \bigr), &j = i \\
%         \bigl(i+1;\, (x,\, \sigma^{n-i-1}_{j-i-1} y) \bigr), &j > i,\; n-i-1\neq 0 \\
%         % \bigl(-1;\,\sigma^n_j y \bigr), &i=-1
%     \end{cases}
% \end{align}
となる．

\begin{myexample}[label=ex:join-0]{join {$\Delta^0 \star \Delta^0$}}
    $\textcolor{red}{\Delta^0} \star \textcolor{blue}{\Delta^0}$ を計算してみよう\footnote{左右の区別を付けるために色を付けた．}．まず対象は
    \begin{align}
        (\Delta^0 \star \Delta^0)_0 &= \textcolor{red}{\Delta^0}_0 \sqcup \textcolor{blue}{\Delta^0}_0 
        = \left\{ 
            \begin{tikzpicture}[baseline={([yshift=-.5ex]current bounding box.center)}]
                \path coordinate[bullet,blue,label=below:$\{0\}$] (x)
                ++(0,1) coordinate[bullet,red,label=above:$\{0\}$] (y)
                ;
            \end{tikzpicture}
        \right\}
    \end{align}
    である．1-射は
    \begin{align}
        (\Delta^0 \star \Delta^0)_1 &= \textcolor{red}{\Delta^0}_1 \sqcup (\textcolor{DarkGreen}{\Delta^0_0 \times \Delta^0_0}) \sqcup \textcolor{blue}{\Delta^0}_1
    \end{align}
    であるが，\eqref{eq:face-join}より始点関数は
    \begin{align}
        \partial^1_1|_{\textcolor{DarkGreen}{\Delta^0_0 \times \Delta^0_0}} \colon \textcolor{DarkGreen}{\Delta^0_0 \times \Delta^0_0} &\lto \textcolor{red}{\Delta^0}_0 \sqcup \textcolor{blue}{\Delta^0}_0, \\
        (\{0\},\, \{0\}) &\lmto \textcolor{red}{\{0\}}
    \end{align}
    終点関数は
    \begin{align}
        \partial^1_0|_{\textcolor{DarkGreen}{\Delta^0_0 \times \Delta^0_0}} \colon \textcolor{DarkGreen}{\Delta^0_0 \times \Delta^0_0} &\lto \textcolor{red}{\Delta^0}_0 \sqcup \textcolor{blue}{\Delta^0}_0, \\
        (\{0\},\, \{0\}) &\lmto \textcolor{blue}{\{0\}}
    \end{align}
    となるため，$\textcolor{red}{\Delta^0}_1,\, \textcolor{blue}{\Delta^0}_1$ が縮退していることを考慮すると
    \begin{align}
        (\Delta^0 \star \Delta^0)_1
        &= \left\{ 
            \begin{tikzpicture}[baseline={([yshift=-.5ex]current bounding box.center)}]
                \path coordinate[bullet,blue,label=below:$\{0\}$] (x)
                ++(0,1) coordinate[bullet,red,label=above:$\{0\}$] (y)
                ;
                \draw[->-=.5,DarkGreen] (y) -- (x);
            \end{tikzpicture}
        \right\}     
    \end{align}
    と図示できる．
\end{myexample}

\begin{myexample}[label=ex:join-1]{join {$\Delta^1 \star \Delta^0$}}
    $\Delta^1 \star \Delta^0$ を計算してみよう．まず対象は
    \begin{align}
        (\Delta^1 \star \Delta^0)_0 &= \textcolor{red}{\Delta^1}_0 \sqcup \textcolor{blue}{\Delta^0}_0 
        = \left\{ 
            \begin{tikzpicture}[baseline={([yshift=-.5ex]current bounding box.center)}]
                \path coordinate[bullet,blue,label=below:$\{0\}$] (x)
                +(120:1) coordinate[bullet,red,label=above:$\{0\}$] (y)
                +(60:1) coordinate[bullet,red,label=above:$\{1\}$] (z)
                ;
            \end{tikzpicture}
        \right\}
    \end{align}
    である．1-射は
    \begin{align}
        (\Delta^1 \star \Delta^0)_1 &= \textcolor{red}{\Delta^1}_1 \sqcup (\textcolor{DarkGreen}{\Delta^1_0 \times \Delta^0_0}) \sqcup \textcolor{blue}{\Delta^0}_1
    \end{align}
    であるが，\eqref{eq:face-join}より始点関数は
    \begin{align}
        \partial^1_1|_{\textcolor{DarkGreen}{\Delta^1_0 \times \Delta^0_0}} \colon \textcolor{DarkGreen}{\Delta^1_0 \times \Delta^0_0} &\lto \textcolor{red}{\Delta^1}_0 \sqcup \textcolor{blue}{\Delta^0}_0, \\
        (x,\, \{0\}) &\lmto \textcolor{red}{x}
    \end{align}
    終点関数は
    \begin{align}
        \partial^1_0|_{\textcolor{DarkGreen}{\Delta^1_0 \times \Delta^0_0}} \colon \textcolor{DarkGreen}{\Delta^1_0 \times \Delta^0_0} &\lto \textcolor{red}{\Delta^1}_0 \sqcup \textcolor{blue}{\Delta^0}_0, \\
        (x,\, \{0\}) &\lmto \textcolor{blue}{\{0\}}
    \end{align}
    となるため，
    \begin{align}
        \Delta^1 \star \Delta^0
        &= \left\{ 
            \begin{tikzpicture}[baseline={([yshift=-.5ex]current bounding box.center)}]
                \path coordinate[bullet,blue,label=below:$\{0\}$] (x)
                +(120:3) coordinate[bullet,red,label=above:$\{0\}$] (y)
                +(60:3) coordinate[bullet,red,label=above:$\{1\}$] (z)
                ;
                \draw[->-=.5,red] (y) -- (z);
                \draw[->-=.5,DarkGreen] (y) -- (x);
                \draw[->-=.5,DarkGreen] (z) -- (x);
                \node[gray] at (barycentric cs:x=1,y=1,z=1) {$(\Id_{[1]},\, \{0\})$};
                \begin{scope}[on background layer]
                    \fill[gray!20] (x.center) -- (y.center) -- (z.center);
                \end{scope}
            \end{tikzpicture}
        \right\}  
    \end{align}
    と図示できる．ただし，三角形の内部は2-射 $(\Id_{[1]},\, \{0\}) \in \Delta^1_1 \times \Delta^0_0 \subset (\Delta^1 \star \Delta^0)_2$ が埋めている．
    同様に
    \begin{align}
        (\Delta^0 \star \Delta^1)_1
        &= \left\{ 
            \begin{tikzpicture}[baseline={([yshift=-.5ex]current bounding box.center)}]
                \path coordinate[bullet,red,label=below:$\{0\}$] (x)
                +(120:3) coordinate[bullet,blue,label=above:$\{0\}$] (y)
                +(60:3) coordinate[bullet,blue,label=above:$\{1\}$] (z)
                ;
                \draw[->-=.5,blue] (y) -- (z);
                \draw[-<-=.5,DarkGreen] (y) -- (x);
                \draw[-<-=.5,DarkGreen] (z) -- (x);
                \node[gray] at (barycentric cs:x=1,y=1,z=1) {$(\{0\},\, \Id_{[1]})$};
                \begin{scope}[on background layer]
                    \fill[gray!20] (x.center) -- (y.center) -- (z.center);
                \end{scope}
            \end{tikzpicture}
        \right\}
    \end{align}
    であることが分かる．
\end{myexample}

\begin{myexample}[label=ex:join-2]{join {$\Delta^2 \star \Delta^0$}}
    $\Delta^2 \star \Delta^0$ を計算してみよう．まず
    \begin{align}
        (\Delta^2 \star \Delta^0)_0 
        &= \textcolor{red}{\Delta^2}_0 \sqcup \textcolor{blue}{\Delta^0}_0 
        = \left\{ 
            \begin{tikzpicture}[baseline={([yshift=-.5ex]current bounding box.center)}]
                \path coordinate[bullet,blue,label=below:$\{0\}$] (x)
                ++(0,1) coordinate[bullet,red,label=right:$\{2\}$] (z)
                +(150:1) coordinate[bullet,red,label=above left:$\{1\}$] (y)
                +(30:1) coordinate[bullet,red,label=above right:$\{0\}$] (w)
                ;
            \end{tikzpicture}
        \right\} \\
    \end{align}
    である．次に1射は
    \begin{align}
        (\Delta^2 \star \Delta^0)_1 
        &= \textcolor{red}{\Delta^2}_1 \sqcup (\textcolor{DarkGreen}{\Delta^2_0 \times \Delta^0_0}) \sqcup \textcolor{blue}{\Delta^0}_1
    \end{align}
    であるが，\eqref{eq:face-join}より始点関数は
    \begin{align}
        \partial^1_1|_{\textcolor{DarkGreen}{\Delta^2_0 \times \Delta^0_0}} \colon \textcolor{DarkGreen}{\Delta^2_0 \times \Delta^0_0} &\lto \textcolor{red}{\Delta^2}_0 \sqcup \textcolor{blue}{\Delta^0}_0,\\
        (x,\, \Id_{\{0\}}) &\lmto \textcolor{red}{x}
    \end{align}
    となり，終点関数は
    \begin{align}
        \partial^1_0|_{\textcolor{DarkGreen}{\Delta^2_0 \times \Delta^0_0}} \colon \textcolor{DarkGreen}{\Delta^2_0 \times \Delta^0_0} &\lto \textcolor{red}{\Delta^2}_0 \sqcup \textcolor{blue}{\Delta^0}_0,\\
        (x,\, \Id_{\{0\}}) &\lmto \textcolor{blue}{\{0\}}
    \end{align}
    となる．従って図式\ref{fig:1-simp}に倣うと
    \begin{align}
        (\Delta^2 \star \Delta^0)_1 
        &= \textcolor{red}{\Delta^2}_1 \sqcup (\textcolor{DarkGreen}{\Delta^2_0 \times \Delta^0_0}) \sqcup \textcolor{blue}{\Delta^0}_1 =
        \left\{ 
            \begin{tikzpicture}[baseline={([yshift=-.5ex]current bounding box.center)}]
                \path coordinate[bullet,blue,label=below:$\{0\}$] (x)
                ++(0,1.5) coordinate[bullet,red,label=right:$\{1\}$] (y)
                +(150:1) coordinate[bullet,red,label=above left:$\{0\}$] (z)
                +(30:1) coordinate[bullet,red,label=above right:$\{2\}$] (w)
                ;
                \draw[->-=.5,DarkGreen] (y) -- (x);
                \draw[->-=.5,DarkGreen] (z) -- (x);
                \draw[->-=.5,DarkGreen] (w) -- (x);
                \draw[->-=.5,red] (z) -- (y);
                \draw[->-=.5,red] (y) -- (w);
                \draw[->-=.5,red] (z) -- (w);
            \end{tikzpicture}
        \right\} 
    \end{align}
    と図示できる．ただし，四面体の内部は3-射 $(\Id_{[2]},\, \textcolor{blue}{\{0\}}) \in \Delta^2_2 \times \Delta^0_0 \subset (\Delta^1 \star \Delta^0)_3$ が埋めている．
    同様に，$\Delta^0 \star \Delta^2$ の1-射を図示すると
    \begin{align}
        (\Delta^0 \star \Delta^2)_1 
        &= \textcolor{red}{\Delta^0}_1 \sqcup (\textcolor{DarkGreen}{\Delta^0_0 \times \Delta^2_0}) \sqcup \textcolor{blue}{\Delta^2}_1 =
        \left\{ 
            \begin{tikzpicture}[baseline={([yshift=-.5ex]current bounding box.center)}]
                \path coordinate[bullet,blue,label=below left:$\{0\}$] (x)
                ++(30:1) coordinate[bullet,blue,label=above right:$\{1\}$] (y)
                +(-30:1) coordinate[bullet,blue,label=below right:$\{2\}$] (z)
                +(0,1.5) coordinate[bullet,red,label=above:$\{0\}$] (w)
                ;
                \draw[->-=.5,blue] (x) -- (y);
                \draw[->-=.5,blue] (x) -- (z);
                \draw[->-=.5,blue] (y) -- (z);
                \draw[-<-=.5,DarkGreen] (x) -- (w);
                \draw[-<-=.5,DarkGreen] (y) -- (w);
                \draw[-<-=.5,DarkGreen] (z) -- (w);
            \end{tikzpicture}
        \right\} 
    \end{align}
    のようになる．
\end{myexample}

\begin{myexample}[label=ex:joint-n]{join $\Delta^n \star \Delta^0$}
    \exref{ex:join-0}, \exref{ex:join-1}, \exref{ex:join-2}で見てきた通り，$\Delta^n \star \Delta^0$ とは，$\Delta^n$ に1つ頂点を付け足してそれをもともとあった $\Delta^n$ の全ての頂点と繋ぎ，内部を $n+1$-単体で埋めることによって構成される．
    故に，自然同型 $\Delta^n \star \Delta^0 \cong \Delta^{n+1}$ がある．特に $\Delta^0$ の頂点は $\Delta^{n+1}$ の頂点 $\{n+1\} \in \Delta^{n+1}_0$ に埋め込まれている．
    % この自然同型を明示的に書くと以下の通り：
    % \begin{align}
    %     (\Delta^n \star \Delta^0)_m &\lto \Delta^{n+1}_m, \\
    %     \bigl( [i];\, (x,\, y) \bigr) &\lto 
    % \end{align}
\end{myexample}


\begin{mylem}[label=lem:join-infty]{{$(\infty,\, 1)$}-圏同士のjoinは{$(\infty,\, 1)$}-圏}
    \hyperref[def:infinity-1]{$(\infty,\, 1)$-圏}同士の\hyperref[def:Simp-Join]{join}は $(\infty,\, 1)$-圏である．
\end{mylem}

\begin{proof}
    ~\cite[Proposition 1.2.8.3]{lurie2008higher}
\end{proof}

\begin{mydef}[label=def:overcat-infty,breakable]{スライス {$(\infty,\, 1)$}-圏}
    \hyperref[def:infinity-1]{$(\infty,\, 1)$-圏} $\Cat{D},\, \Cat{C}$ および\hyperref[def:infinity-1]{$(\infty,\, 1)$-圏の関手} $p \in \Hom{\sSet} (\Cat{D},\, \Cat{C})$ を与える．
    $p$ に沿った $\Cat{C}$ の\textbf{スライス圏} (overcategory) 
    \begin{align}
        \bm{\Cat{C}_{/p}} \colon \OP{\Delta} &\lto \SETS
    \end{align}
    を以下で定義する：
    \begin{itemize}
        \item $\forall [n] \in \Obj{\OP{\Delta}}$ に対して，集合
        \begin{align}
            \Hom{p} (\Delta^n \star \Cat{D},\, \Cat{C})
            \coloneqq \bigl\{\, f \in \Hom{\sSet} (\Delta^n \star \Cat{D},\, \Cat{C}) \bigm| f|_{\Cat{D}} = p \,\bigr\} 
        \end{align}
        を対応付ける\footnote{$f|_{\Cat{D}}$ というのは，\hyperref[def:Simp-Join]{joinの定義}における $(\Delta^n \star \Cat{D})_k$ のdisjoint unionのうち，添字 $i=0$ が振られている成分への制限を意味する．}．
        \item $\forall \alpha \in \Hom{\OP{\Delta}} ([m],\, [n])$ に対して，写像
        \begin{align}
            \Cat{C}_{/p} (\alpha) \colon \Hom{p} (\Delta^m \star \Cat{D},\, \Cat{C}) &\lto \Hom{p} (\Delta^n \star \Cat{D},\, \Cat{C}), \\
            f &\lmto f \circ (\alpha_* \star \Id_{\Cat{D}})
        \end{align}
        を対応付ける．
    \end{itemize}
    実際，単体的集合 $\Cat{C}_{/p}$ は $(\infty,\, 1)$-圏である~\cite[\href{https://kerodon.net/tag/018F}{Tag 018F}]{kerodon}．
    
    \tcblower

    $p$ に沿った $\Cat{C}$ の\textbf{コスライス圏} (undercategory) 
    \begin{align}
        \bm{\Cat{C}_{p/}} \colon \OP{\Delta} &\lto \SETS
    \end{align}
    を以下で定義する：
    \begin{itemize}
        \item $\forall [n] \in \Obj{\OP{\Delta}}$ に対して，集合
        \begin{align}
            \Hom{p} (\Cat{D} \star \Delta^n,\, \Cat{C})
            \coloneqq \bigl\{\, f \in \Hom{\sSet} (\Cat{D} \star \Delta^n,\, \Cat{C}) \bigm| f|_{\Cat{D}} = p \,\bigr\} 
        \end{align}
        を対応付ける\footnote{$f|_{\Cat{D}}$ というのは，\hyperref[def:Simp-Join]{joinの定義}における $(\Cat{D} \star \Delta^n)_k$ のdisjoint unionのうち，添字 $i=n$ が振られている成分への制限を意味する．}．
        \item $\forall \alpha \in \Hom{\OP{\Delta}} ([m],\, [n])$ に対して，写像
        \begin{align}
            \Cat{C}_{p/} (\alpha) \colon \Hom{p} (\Cat{D} \star \Delta^m,\, \Cat{C}) &\lto \Hom{p} (\Cat{D} \star \Delta^n,\, \Cat{C}), \\
            f &\lmto f \circ (\Id_{\Cat{D}} \star \alpha_*)
        \end{align}
        を対応付ける．
    \end{itemize}
\end{mydef}

特に注目すべきは，\hyperref[def:infinity-1]{$(\infty,\, 1)$-圏の関手} $p \colon \Delta^0 \lto \Cat{C}$ をとった場合である．
このとき $X \coloneqq p_{[0]} (\{0\}) \in \Cat{C}_0$ とおいて $\bm{\Cat{C}_{/X}},\; \bm{\Cat{C}_{X/}}$ などと書く．

まず，$(\infty,\, 1)$-圏 $\Cat{C}_{/X}$ の対象 $\varphi \in (\Cat{C}_{/X})_0 = \Hom{p}(\Delta^0 \star \textcolor{blue}{\Delta^0},\, \Cat{C})$ をとる．
すると\exref{ex:join-0}および $\varphi|_{\textcolor{blue}{\Delta^0}} = p$ の条件から，$\varphi_{[1]} \colon (\Delta^0 \star \textcolor{blue}{\Delta^0})_1 \lto \Cat{C}_1$ とは図式
\begin{align}
    \varphi =  
    \begin{tikzpicture}[baseline={([yshift=-.5ex]current bounding box.center)}]
        \path coordinate[bullet,blue,label=below:$X$] (x)
        ++(0,1) coordinate[bullet,label=above:$\varphi_{[0]}|_{\Delta^0_0}(\{0\})$] (y)
        ;
        \draw[->-=.5] (y) --node[midway,right]{$\varphi_{[1]}|_{\textcolor{DarkGreen}{\Delta^0_0 \times \Delta^0_0}} (\{0\}\to \textcolor{blue}{\{0\}})$} (x);
    \end{tikzpicture}
\end{align}
である．$n \ge 2$ 射に相当する $\varphi_{[n]} \colon (\Delta^0 \star \textcolor{blue}{\Delta^0})_n \lto \Cat{C}_n$ のデータは縮退していて自明である．従って，$\varphi$ は \hyperref[def:slice-category]{$(1,\, 1)$-圏における $X$ 上のスライス圏}の対象と等価なデータを与える．

同様に，$(\infty,\, 1)$-圏 $\Cat{C}_{/X}$ の1-射 $f \in (\Cat{C}_{/X})_1 = \Hom{p}(\Delta^1 \star \textcolor{blue}{\Delta^0},\, \Cat{C})$ とは，\exref{ex:join-1}より
\begin{align}
    \label{eq:slice-1mor}
    f =     
    \begin{tikzpicture}[baseline={([yshift=-.5ex]current bounding box.center)}]
        \path coordinate[bullet,blue,label=below:$X$] (x)
        +(120:3) coordinate[bullet,label=left:$f_{[0]}|_{\Delta^1_0}(\{0\})$] (y)
        +(60:3) coordinate[bullet,label=right:$f_{[0]}|_{\Delta^1_0}(\{1\})$] (z)
        ;
        \draw[->-=.5] (y) -- (z);
        \draw[->-=.5] (y) -- (x);
        \draw[->-=.5] (z) -- (x);
         \node[gray] at (barycentric cs:x=1,y=1,z=1) {$f_{[2]}|_{\Delta^1_1 \times \Delta^0_0}(\Id_{[1]},\, \{0\})$};
                \begin{scope}[on background layer]
                    \fill[gray!20] (x.center) -- (y.center) -- (z.center);
                \end{scope}
        % \fill[red!20] (x) -- (y) -- (z) -- cycle;
    \end{tikzpicture}
\end{align}
のことである．ただし三角形の内部は2-射 $f_{[2]}|_{\Delta^1_1 \times \Delta^0_0}(\Id_{[2]},\, \textcolor{blue}{\{0\}}) \in \Cat{C}_2$ が埋めている．これは\hyperref[def:slice-category]{$(1,\, 1)$-圏における $X$ 上のスライス圏}の射のデータに対応しているが，横向きの矢印を決めるだけでは $f$ がupto \hyperref[def:infty-homotopy-morphism]{homotopy}でしか定まらないという点で異なっている．

$(\infty,\, 1)$-圏 $\Cat{C}_{/X}$ の $n$-射も同様に図示できる．

\begin{myexample}[label=def:infty-forget]{スライス圏からのforgetful functor}
    \hyperref[def:infinity-1]{$(\infty,\, 1)$-圏} $\Cat{C}$ の，$X \colon \Delta^0 \lto \Cat{C}$ に沿った\hyperref[def:overcat-infty]{スライス圏}に対して，\textbf{忘却関手} (forgetful functor)
    \begin{align}
        \text{forget} \colon \Cat{C}_{/X} \lto \Cat{C}
    \end{align}
    を次のように定義する：
    \begin{align}
        \label{eq:forget}
        \text{forget}_{[n]} \colon (\Cat{C}_{/X})_n = \Hom{X} (\Delta^n \star \Delta^0,\, \Cat{C}) &\lto \Cat{C}_n = \Hom{\sSet}(\Delta^n,\, \Cat{C}), \\
        f &\lmto f|_{\Delta^n}
    \end{align}
    $n = 0,\, 1,\, 2$ の場合，i.e. \exref{ex:join-0}, \exref{ex:join-1}, \exref{ex:join-2}の図式においては，ちょうど $X \in \Cat{C}_0$ に対応する青色の頂点（コーンポイント）を除去する操作に対応している．
    \eqref{eq:forget}の定義は\hyperref[def:infty-fib]{右ファイブレーション}である~\cite[\href{https://kerodon.net/tag/018F}{Tag 018F}]{kerodon}．
\end{myexample}


\subsection{{$(\infty,\, 1)$}-圏のlimit/colimit}

後の議論のため，先取りして $(\infty,\, 1)$-圏におけるlimit/colimitを定義しておこう．
$(\infty,\, 1)$-圏のモデルとして\hyperref[def:infinity-1]{quasi-category}を採用する場合，これは\textbf{homotopy limit/colimit}と呼ばれることもある．

\begin{mydef}[label=def:left-right-cone]{左/右錐}
    $(\infty,\, 1)$-圏 $\Cat{C}$ を与える．
    \begin{itemize}
        \item \hyperref[def:Simp-Join]{join}によって得られる $(\infty,\, 1)$-圏 $\Delta^0 \star \Cat{C}$ のことを $\Cat{C}$ の\textbf{左錐} (left cone) と呼び，$\leftCone{\Cat{C}}$ と書く．
        \item \hyperref[def:Simp-Join]{join}によって得られる $(\infty,\, 1)$-圏 $\Cat{C} \star \Delta^0$ のことを $\Cat{C}$ の\textbf{右錐} (right cone) と呼び，$\rightCone{\Cat{C}}$ と書く．
    \end{itemize}
\end{mydef}

\begin{mydef}[label=def:initial-final-infty]{{$(\infty,\, 1)$}-圏における始対象と終対象}
    \begin{itemize}
        \item \hyperref[def:infinity-1]{$(\infty,\, 1)$-圏} $\Cat{C}$ における対象 $x_0 \in \Cat{C}_0$ が\textbf{始対象} (initial object) であるとは，$\forall x \in \Cat{C}_0$ に対して\hyperref[def:Map]{射の空間} $\inftyMap{\Cat{C}} (x_0,\, x) \in \Obj{\Kan}$ が $\Delta^0$ と\hyperref[def:htpy-Kan]{ホモトピー同値}であること．
        % \hyperref[def:hcat-infty]{ホモトピー圏} $\htpy{\Cat{C}}$ における始対象\footnote{$(1,\, 1)$-圏の\textbf{始対象}とは，空の図式における\hyperref[def:colim]{余極限}のこと．}であること．
        \item \hyperref[def:infinity-1]{$(\infty,\, 1)$-圏} $\Cat{C}$ における対象 $x_0 \in \Cat{C}_0$ が\textbf{終対象} (final object) であるとは，$\forall x \in \Cat{C}_0$ に対して\hyperref[def:Map]{射の空間} $\inftyMap{\Cat{C}} (x,\, x_0) \in \Obj{\Kan}$ が $\Delta^0$ と\hyperref[def:htpy-Kan]{ホモトピー同値}であること．
        % それが\hyperref[def:hcat-infty]{ホモトピー圏} $\htpy{\Cat{C}}$ における終対象\footnote{$(1,\, 1)$-圏の\textbf{終対象}とは，空の図式における\hyperref[def:lim]{極限}のこと．}であること．
    \end{itemize}
\end{mydef}

\begin{mydef}[label=def:lim-colim-infty]{{$(\infty,\, 1)$}-圏のlimit/colimit}
    \hyperref[def:infinity-1]{$(\infty,\, 1)$-圏の関手} $D \colon \Cat{I} \lto \Cat{C}$ を与える\footnote{\hyperref[def:diagram]{$(1,\, 1)$-圏の場合}からのアナロジーで，$D$ を図式と見做す．}．
    \begin{itemize}
        \item $D$ の\textbf{limit}とは，\hyperref[def:overcat-infty]{スライス $(\infty,\, 1)$-圏} $\Cat{C}_{/D}$ における\hyperref[def:initial-final-infty]{終対象}のこと．$\bm{\lim D} \in (\Cat{C}_{/D})_0 = \Hom{D} (\leftCone{\Cat{I}},\, \Cat{C})$ と書く．
        \item $D$ の\textbf{colimit}とは，\hyperref[def:overcat-infty]{コスライス $(\infty,\, 1)$-圏} $\Cat{C}_{D/}$ における\hyperref[def:initial-final-infty]{始対象}のこと．$\bm{\colim D} \in (\Cat{C}_{D/})_0 = \Hom{D} (\rightCone{\Cat{I}},\, \Cat{C})$ と書く．
    \end{itemize}
\end{mydef}

以降で登場する\hyperref[def:lim-colim-infty]{$(\infty,\, 1)$-圏のlimit}は，$(\infty,\, 1)$-圏の成す $(\infty,\, 1)$-圏 $\CATinfty$ に値をとる図式 $D \colon \Cat{I} \lto \CATinfty$ に関するものがほとんどである．

\begin{myexample}[label=def:pullback-infty]{pullback}
    % 単体的集合の積 $S \times T \colon \OP{\Delta} \lto \SETS$ における面写像とは 
    % \begin{align}
    %     \partial^n_j \colon S_n \times T_n &\lto S_{n-1} \times T_{n-1}, \\
    %     (x,\, y) &\lmto (\partial^n_j x,\, \partial^n_j y)
    % \end{align}
    % のことであった．
    % 故に，単体的集合 $\Delta^1 \times \Delta^1$ は，$\Delta^1_0 \eqqcolon \{\bullet_0,\, \bullet_1\}$ とおくと
    % \begin{align}
    %     \Delta^1 \times \Delta^1 =
    %     \begin{tikzpicture}[baseline={([yshift=-.5ex]current bounding box.center)}]
    %         \path  coordinate[bullet,label=below left:{$(0,\,1)$}] (x)
    %         +(1,0) coordinate[bullet,label=below right:{$(1,\,1)$}] (y)
    %         +(0,1) coordinate[bullet,label=above left:{$(0,\,0)$}] (z)
    %         +(1,1) coordinate[bullet,label=above right:{$(1,\,0)$}] (w)
    %         ;
    %         \draw[->-=.5] (z) -- (x);
    %         \draw[->-=.5] (z) -- (y);
    %         \draw[->-=.5] (z) -- (w);
    %         \draw[->-=.5] (x) -- (y);
    %         \draw[->-=.5] (w) -- (y);
    %     \end{tikzpicture}
    % \end{align}
    % と図示できる．ただし，2-射以上は縮退していて見えない．
    
    　\hyperref[def:infinity-1]{$(\infty,\, 1)$-圏の関手} $D \colon (\Delta^1 \times \Delta^1) \setminus \bigl\{(0,\, 0)\bigr\} \lto \Cat{C}$ の\hyperref[def:lim-colim-infty]{limit}のことを（存在すれば）\textbf{pullback}と呼び，
    \begin{align}
        \bm{D(0,\, 1)} \bm{\times}_{\bm{D(1,\, 1)}} \bm{D(1,\, 0)} \coloneqq \lim D
    \end{align}
    と書く．
    % \exref{ex:Map1}および補題\ref{lem:n-fold-Map}より，$D \in \Hom{\sSet}(\Delta^1 \times \Delta^1,\, \Cat{C})$ は
    % \begin{align}
    %     D = \quad
    %     \begin{tikzpicture}[baseline={([yshift=-.5ex]current bounding box.center)}]
    %         \path  coordinate[bullet,label=below left:{$D_{[0]} (0,\, 1)$}] (x)
    %         +(4,0) coordinate[bullet,label=below right:{$D_{[0]} (1,\, 1)$}] (y)
    %         +(0,4) coordinate[bullet,label=above left:{$D_{[0]} (0,\, 0)$}] (z)
    %         +(4,4) coordinate[bullet,label=above right:{$D_{[0]} (1,\, 0)$}] (w)
    %         ;
    %         \begin{scope}[on background layer]
    %             \fill[red!20]  (z.center) -- (x.center) -- (y.center) -- cycle;
    %             \fill[blue!20] (z.center) -- (w.center) -- (y.center) -- cycle;
    %         \end{scope}
    %         \node[red] at (barycentric cs:z=1,x=1,y=1) {$D_{[2]}\bigl(s^1_0{}^*(\Id_{[1]}),\, s^1_1{}^* (\Id_{[1]})\bigr)$};
    %         \node[blue] at (barycentric cs:z=1,w=1,y=1) {$D_{[2]}\bigl(s^1_1{}^*(\Id_{[1]}),\, s^1_0{}^* (\Id_{[1]})\bigr)$};
    %         \draw[->-=.5] (z) -- (x);
    %         \draw[->-=.5] (z) -- (y);
    %         \draw[->-=.5] (z) -- node[midway,above]{$\sigma^0_0(x) \eqqcolon \Id_x$} (w);
    %         \draw[->-=.5] (x) -- node[midway,below]{$\sigma^0_0(y) \eqqcolon \Id_y$} (y);
    %         \draw[->-=.5] (w) -- (y);
    %     \end{tikzpicture}
    % \end{align}
    % と図示できる．従って\exref{ex:join-2}から，$D(0,\, 1) \times_{D(1,\, 1)} D(1,\, 0)$ とは次のように図示できる：
\end{myexample}

\exref{def:pullback-infty}は，今後種々の $(\infty,\, 1)$-圏の定義に用いるため，具体的な公式を紹介しておこう．
$(\infty,\, 1)$-圏 $\Cat{C}$ に対して，
\textbf{射の $(\infty,\, 1)$-圏} (arrow category) $\FUN(\Delta^1,\, \Cat{C})$ のうち\hyperref[def:isom-infty]{同型射}全体が成す\hyperref[def:fullsub-infty]{充満部分 $(\infty,\, 1)$-圏}を
\begin{align}
    \ISOM(\Cat{C}) \hookrightarrow \FUN(\Delta^1,\, \Cat{C})
\end{align}
と書く．

\begin{mylem}[label=lem:pullback-infty]{{$(\infty,\, 1)$}-圏の引き戻しの公式}
    \begin{itemize}
        \item \hyperref[def:infinity-1]{$(\infty,\, 1)$-圏} $\Cat{C},\, \Cat{C}_0,\, \Cat{C}_1$
        \item \hyperref[def:infinity-1]{$(\infty,\, 1)$-圏の関手} $F_0 \colon \Cat{C}_0 \lto \Cat{C},\; F_1 \colon \Cat{C}_1 \lto \Cat{C}$
    \end{itemize}
    を与える．さらに，$(\infty,\, 1)$-圏\footnote{~\cite[\href{https://kerodon.net/tag/0330}{Tag 0330}]{kerodon}よりこれは $(\infty,\, 1)$-圏である．} $\Cat{C}_0 \htimes_{\Cat{C}} \Cat{C}_1$ を，$(1,\, 1)$-圏 $\sSet$ における\hyperref[def:pullback-pushout]{引き戻し}
    \begin{align}
        \label{eq:htpy-pullback}
        \Cat{C}_0 \htimes_{\Cat{C}} \Cat{C}_1 \coloneqq \Cat{C}_0 \times_{\FUN (\{0\},\, \Cat{C})} \ISOM(\Cat{C}) \times_{\FUN(\{1\},\, \Cat{C})} \Cat{C}_1
    \end{align}
    として定義する~\cite[\href{https://kerodon.net/tag/032Z}{Tag 032Z}]{kerodon}．
    このとき，$\Cat{C}_0 \htimes_{\Cat{C}} \Cat{C}_1$ は $(\infty,\, 1)$-圏 $\CATinfty$ における図式 $\Cat{C}_0 \xrightarrow{F_0} \Cat{C} \xleftarrow{F_1} \Cat{C}_1$ の\hyperref[def:lim-colim-infty]{limit}を与える．
    i.e. $\sSet$ における図式\footnote{up to isomorphism で可換なのであって，可換とは限らない！}
    \begin{center}
        % https://tikzcd.yichuanshen.de/#N4Igdg9gJgpgziAXAbVABwnAlgFyxMJZABgBoBGAXVJADcBDAGwFcYkQAdDgYXp2G4BfAPrEQg0uky58hFOQrU6TVuy69+Q8ZJAZseAkQXElDFm0ScefASPLap+2UTImaZ1ZfU2hogARcABZ4ALbwwsDemoKCAdbRwvaCSjBQAObwRKAAZgBOECFIAEw0OBBICsrm7ABiiQ4geQVIZCBlFe4qFiB1YjSM9ABGMIwACtIGciC5WGnBDU2FiADMpeWIrR7dXCF8gbkhwGj5AFYikfG2ooIg-UMj406GljNzOAv5S6tt6yVVnlZdjh9odjhAzhEoldyDdkoIgA
        \begin{tikzcd}
        \Cat{C}_0 \htimes_{\Cat{C}} \Cat{C}_1 \arrow[d, "\mathrm{proj}_{\Cat{C}_0}"'] \arrow[r, "\mathrm{proj}_{\Cat{C}_1}"] & \Cat{C}_1 \arrow[d, "F_1"] \\
        \Cat{C}_0 \arrow[r, "F_0"']                                                                                          & \Cat{C}                   
        \end{tikzcd}
    \end{center}
    は\hyperref[def:pullback-infty]{$(\infty,\, 1)$-圏の引き戻し}を与える．
\end{mylem}

\begin{proof}
    
\end{proof}

\eqref{eq:htpy-pullback}をより具体的に書くと，これは $\sSet$ における2重の引き戻し
\begin{center}
    % https://tikzcd.yichuanshen.de/#N4Igdg9gJgpgziAXAbVABwnAlgFyxMJZABgBoAmAXVJADcBDAGwFcYkQAdDgYXp2G4BfAPrEQg0uky58hFAEYK1Ok1bsuAMQCqAOS4AjLAHNGACi7BiXCV1IACLr35CDxgE4BKcZJAZseAiJFeWUGFjZETg4ASQBlAHkAWXMePgFBLwkpf1kiclJiUNUIqKd04XlvbJlAlHyQmjC1SM1dVxMU4HlrUlsHVOdBds8q32kAuRICovD1AfLifoALPABbeGFgRzShQX6yoQrxZRgoI3giUAAzNwhVpABWGhwIJGIskBu7t+fXxEUVLNIhpRCAaIx6PoYIwAArjXKRNzGFajL73RD5EAvJAApolFLRKCbLgAEWhOHoAD0DMxGIwYDg9lw1vA7FBKfIKsARAAqDyUnmo27ozHYxAAFkaxXYBKJWw4ZMYFOpHH0tPpjP6LLgbI5om5wj5AqF30QAGZfkhJYDmiAQZUPmjHpbzY7hc6sX9yIJKIIgA
    \begin{tikzcd}
    \Cat{C}_0 \htimes_{\Cat{C}} \Cat{C}_1 \arrow[dd] \arrow[rr] \arrow[rd] &             \arrow[dr, phantom," " {pullback}, very near start, color=black]                                                                                                                 & \Cat{C}_1 \arrow[d, "F_1"]         \\
                                    \arrow[dr, phantom," " {pullback}, very near start, color=black]      & \ISOM(\Cat{C}) \arrow[d, "(\Id_{\Delta^\bullet} \times d^1_1{}_*)^*"] \arrow[r, "(\Id_{\Delta^\bullet} \times d^1_0{}_*)^*"] & {\FUN\bigl(\{1\},\, \Cat{C}\bigr) \cong \Cat{C}} \\
    \Cat{C}_0 \arrow[r, "F_0"']                                            & {\FUN\bigl(\{0\},\, \Cat{C}\bigr) \cong \Cat{C}}                                                                                          &                                   
    \end{tikzcd}
\end{center}
のことである．
命題\ref{prop:SimpSet-basic}-(2) の証明より，$(1,\, 1)$-圏 $\sSet$ における引き戻しとは $n$-単体毎に\hyperref[def:pullback-pushout]{$\SETS$ における引き戻し}を取ることで構成されるから，この引き戻しは具体的に計算できる．

\begin{mylem}[label=lem:pullback-iso]{{$(\infty,\, 1)$-圏}の引き戻しとアイソファイブレーション}
    \begin{itemize}
        \item \hyperref[def:infinity-1]{$(\infty,\, 1)$-圏} $\Cat{C},\, \Cat{C}_0,\, \Cat{C}_1$
        \item \hyperref[def:infinity-1]{$(\infty,\, 1)$-圏の関手} $F_0 \colon \Cat{C}_0 \lto \Cat{C},\; F_1 \colon \Cat{C}_1 \lto \Cat{C}$
    \end{itemize}
    を与える．このとき $F_0,\, F_1$ の少なくともどちらか一方が\hyperref[def:isofibration]{アイソファイブレーション}ならば，
    \hyperref[def:equiv-infty]{$(\infty,\, 1)$-圏同値}
    \begin{align}
        \Cat{C}_0 \times_{\Cat{C}} \Cat{C}_1 \overset{\simeq}{\hookrightarrow} \Cat{C}_0 \htimes_{\Cat{C}} \Cat{C}_1
    \end{align}
    がある．ただし，$\Cat{C}_0 \times_{\Cat{C}} \Cat{C}_1$ は $(1,\, 1)$-圏 $\sSet$ における引き戻しである．
\end{mylem}

\begin{proof}
    ~\cite[\href{https://kerodon.net/tag/033T}{Tag 033T}]{kerodon}
\end{proof}

補題\ref{lem:pullback-infty}と補題\ref{lem:pullback-iso}を併せると，$(\infty,\, 1)$-圏 $\CATinfty$ の図式 $\Cat{C}_0 \xrightarrow{F_0} \Cat{C} \xleftarrow{F_1} \Cat{C}_1$ において $F_0,\, F_1$ のどちらか一方が\hyperref[def:isofibration]{アイソファイブレーション}ならば，
その図式の\hyperref[def:pullback-infty]{$(\infty,\, 1)$-圏の引き戻し}は\textbf{単に $(1,\, 1)$-圏 $\sSet$ における引き戻しとして（up to equivalenceで）計算できることが分かる！}
この事実は実用上極めて便利である．
実際，後に $F_0,\, F_1$ のどちらか一方が\hyperref[def:infty-fib]{右ファイブレーション}である場合が頻発するが，その際の計算では常に補題\ref{lem:pullback-iso}を用いる．



\subsection{Unstraightening construction}

\begin{mytheo}[label=thm:unstraightening]{unstraightening construction}
    \hyperref[def:equiv-infty]{$(\infty,\, 1)$-圏同値}
    \begin{align}
        \inftyPsh{\Cat{B}} \xrightarrow{\simeq} \RFIB_{\Cat{B}}
    \end{align}
    が存在する．
\end{mytheo}

\begin{proof}
    \hyperref[def:equiv-infty]{$(\infty,\, 1)$-圏同値}
    \begin{align}
        \Un\colon \inftyPsh{\Cat{B}} \lto \RFIB_{\Cat{B}}
    \end{align}
    は，次のようにして構成される (\textbf{unstraightening construction})：
    \begin{description}
        \item[\textbf{対象}] 
        $F \in \inftyPsh{\Cat{B}}_0$ に対して，\hyperref[def:pullback-infty]{$(\infty,\, 1)$-圏の引き戻し}
        \begin{center}
            % https://tikzcd.yichuanshen.de/#N4Igdg9gJgpgziAXAbVABwnAlgFyxMJZABgBpiBdUkANwEMAbAVxiRAB12BVMACgDEAlCAC+pdJlz5CKAIzkqtRizacAygAUAggGEAomoD6wAPQAqEaPEgM2PASLzZi+s1aIO7TboNWJd6SIyZ2pXFQ9OHTocYAAhSxFFGCgAc3giUAAzACcIAFskAGZqHAgkACZqBjoAIxgGDUl7GRBsrBSACxwQUOV3EH4-EBz8pHkQUoret1V2HBgADxjMiGy0nATrEYLEMgmyxEKxLNydvcnEWUSRIA
            \begin{tikzcd}
            \textcolor{red}{\Un(F)} \arrow[dr, phantom," " {pullback}, very near start, color=black] \arrow[red,d] \arrow[r] & \SPACES_{/*} \arrow[d, "\text{forget}"] \\
            \Cat{B} \arrow[r, "F"']    & \SPACES                                
            \end{tikzcd}
        \end{center}
        により得られる\hyperref[def:infty-fib]{右ファイブレーション} $\textcolor{red}{\Un(F) \lto}\, \Cat{B} \in (\RFIB_{\Cat{B}}){}_0$ を対応付ける．
        \item[\textbf{$\bm{n}$-射}] 
    \end{description}
    逆向きの\hyperref[def:equiv-infty]{$(\infty,\, 1)$-圏同値} (\textbf{straightning construction})
    \begin{align}
        \St\colon \RFIB_{\Cat{B}} \lto \inftyPsh{\Cat{B}}
    \end{align}
    は難しい．詳細は~\cite[Proposition 2.2.3.11]{lurie2008higher}を参照．
\end{proof}

\subsection{{$(\infty,\, 1)$}-圏における米田埋め込み}

\begin{mydef}[label=def:Tw]{twisted arrow category}
    \hyperref[def:infinity-1]{$(\infty,\, 1)$-圏} $\Cat{C}$ を与える．
    このとき，$\Cat{C}$ の\textbf{twisted arrow category}と呼ばれる $(\infty,\, 1)$-圏を以下で定義する：
    \begin{align}
        \Tw(\Cat{C}) \colon \OP{\Delta} &\lto \SETS, \\
        [n] &\lmto \Hom{\sSet} \bigl( \OP{(\Delta^n)} \star \Delta^n,\, \Cat{C} \bigr), \\
        \bigl( [m] \xrightarrow{\alpha} [n] \bigr) &\lmto \bigl( f \mapsto f \circ (\alpha_* \times \alpha_*) \bigr) 
    \end{align}
\end{mydef}

\hyperref[def:infinity-1]{$(\infty,\, 1)$-圏の関手} 
\begin{align}
    \mathrm{pr} \colon \Tw(\Cat{C}) \lto \OP{\Cat{C}} \times \Cat{C}
\end{align}
を以下で定義すると，これは\hyperref[def:infty-fib]{左ファイブレーション}になる~\cite[\href{https://kerodon.net/tag/03JQ}{Tag 03JQ}]{kerodon}：
\begin{align}
    \mathrm{pr}_{[n]} \colon \Tw(\Cat{C})_n &\lto \OP{\Cat{C}}_n \times \Cat{C}_n = \Hom{\sSet}(\Delta^n,\, \OP{\Cat{C}}) \times \Hom{\sSet}(\Delta^n,\, \Cat{C}) \\
    f &\lmto \bigl( f|_{\OP{(\Delta^n)}},\, f|_{\Delta^n} \bigr) 
\end{align}
\hyperref[def:infinity-1]{$(\infty,\, 1)$-圏}におけるHom関手
\begin{align}
    \label{def:infty-Hom}
    \inftyMap{\Cat{C}} \colon \OP{\Cat{C}} \times \Cat{C} \lto \SPACES
\end{align}
の自然な構成は，\hyperref[thm:unstraightening]{straightning construction}を用いた
\begin{center}
    % https://tikzcd.yichuanshen.de/#N4Igdg9gJgpgziAXAbVABwnAlgFyxMJZABgBpiBdUkANwEMAbAVxiRAB12BVMACgDEAlCAC+pdJlz5CKAIzkqtRizacAygAUAggGEAomoD6wAPQAqEaPEgM2PASLzZi+s1aIO7TboNWJd6SIyZ2pXFQ9OHTocYAAhSxFFGCgAc3giUAAzACcIAFskAGZqHAgkACZqBjoAIxgGDUl7GRBsrBSACxwQUOV3EH4-EBz8pHkQUoret1V2HBgADxjMiGy0nATrEYLEMgmyxEKxLNydvcnEWUSRIA
    \begin{tikzcd}[column sep=large]
    \Tw(\Cat{C}) \arrow[dr, phantom," " {pullback}, very near start, color=black] \arrow[d, "\mathrm{pr}"'] \arrow[r] & \SPACES_{*/} \arrow[d, "\text{forget}"] \\
    \OP{\Cat{C}} \times \Cat{C} \arrow[r, red, "\inftyMap{\Cat{C}} \coloneqq \St(\mathrm{pr})"']    & \SPACES                                
    \end{tikzcd}
\end{center}
である~\cite[I.26., p.19]{Hebestreit2020K-theory}．

\begin{mydef}[label=def:Yoneda-infty]{{$(\infty,\, 1)$}-圏の米田埋め込み (informal)}
    $\Cat{C}$ を\hyperref[def:infinity-1]{$(\infty,\, 1)$-圏}とする．$(\infty,\, 1)$-圏の\textbf{米田埋め込み} (Yoneda embedding)
    \begin{align}
        \yo \colon \Cat{C} \lto \inftyPsh{\Cat{C}}
    \end{align}
    とは，$\Cat{C}$ から\hyperref[def:infinity-presheaf]{$(\infty,\, 1)$-前層の成す $(\infty,\, 1)$-圏} $\inftyPsh{\Cat{C}}$ への\hyperref[def:infinity-1]{$(\infty,\, 1)$-圏の関手}であって，
    対象 $x \in \Cat{C}_0$ に対して
    \begin{align}
        \yo_{[0]}(x)_{[0]} \colon \OP{\Cat{C}}_0 &\lto \SPACES_0, \\
        y &\lmto \inftyMap{\Cat{C}} (y,\, x)
    \end{align}
    を充たす\footnote{$(\infty,\, 1)$-圏 $\Cat{C}$ において，対象 $x,\, y \in \Cat{C}_0$ の間の\hyperref[def:Map]{射の空間} $\inftyMap{\Cat{C}}(x,\, y)$ はKan複体を成すのだった~\cite[\href{https://kerodon.net/tag/01JC}{Tag 01JC}]{kerodon}．}ような $(\infty,\, 1)$-圏の関手 $\yo_{[0]}(x) \in \inftyPsh{\Cat{C}}_0 = \Hom{\sSet}(\OP{\Cat{C}},\, \SPACES)$ を対応付けるもののこと\footnote{厳密な構成については~\cite[\href{https://kerodon.net/tag/03NF}{Tag 03NF}]{kerodon}を参照．}．
\end{mydef}
% \subsection{$\mathcal{B}\mathrm{sc}$ の構造}



\subsection{層状化空間の接構造}

\begin{mydef}[label=def:Enter]{enter-path category}
    \hyperref[def:c-smooth]{conically smoothな層状化空間} $X \in \Obj{\Strat}$ の\textbf{enter-path $(\infty,\, 1)$-category}とは，
    $(\infty,\, 1)$-圏 \hyperref[def:Strat-infty]{$\BSC$} の\hyperref[def:overcat-infty]{スライス圏}
    \begin{align}
        \ENTR{X} \coloneqq \BSC_{/X}
    \end{align}
    のこと．
\end{mydef}

\begin{mydef}[label=def:tangent-classifier]{tangent classifier}
    $\iota \colon \BSC \hookrightarrow \SNGLR$ を包含とする．
    \textbf{tangent classifier}とは，\hyperref[def:infinity-1]{$(\infty,\, 1)$-圏の関手}
    \begin{align}
        \bm{\tau} \colon \SNGLR \xrightarrow{\yo} \inftyPsh{\SNGLR} \xrightarrow{\iota^*} \inftyPsh{\BSC}
    \end{align}
    のこと．
\end{mydef}

定義\ref{def:Yoneda-infty}より，
\hyperref[def:tangent-classifier]{tangent classifier}は\hyperref[def:c-smooth]{conically smoothな層状化空間} $X \in \SNGLR_0$ に対して $(\infty,\, 1)$-圏の表現可能前層
\begin{align}
    \label{eq:representable-infty}
    \tau_{[0]}(X) = \inftyMap{\BSC}(\, \mhyphen\, ,\, X) \in \inftyPsh{\BSC}_0
\end{align}
を対応付ける．

\begin{marker}
    定理\ref{thm:unstraightening}により，\hyperref[def:tangent-classifier]{tangent classifier} $\tau$ のことを
    \begin{align}
        \tau \colon \SNGLR \xrightarrow{\yo} \inftyPsh{\SNGLR} \xrightarrow{\iota^*} \inftyPsh{\BSC} \xrightarrow{\simeq} \RFIB_{\BSC}
    \end{align}
    と見做すこともできる．このとき，\hyperref[def:RFIB]{$\RFIB_{\BSC}$} の構成および $(\infty,\, 1)$-前層\eqref{eq:representable-infty}に対する定理\ref{thm:unstraightening}の具体的構成から，
    \hyperref[def:c-smooth]{conically smoothな層状化空間} $X \in \SNGLR_0$ に対して定まる $(\infty,\, 1)$-圏の\hyperref[def:infty-fib]{右ファイブレーション} $\tau_{[0]}(X) \in (\RFIB_{\BSC})_0$ とは
    忘却関手
    \begin{align}
        \tau_{[0]}(X) \colon \ENTR{X} \lto \BSC
    \end{align}
    のことである．この忘却関手を以下では $\bm{\tau_X} \coloneqq \tau_{[0]}(X)$ と書く．
\end{marker}

\begin{mydef}[label=def:B-mfld,breakable]{{$\Cat{B}$}-多様体}
    \begin{itemize}
        \item \textbf{$\bm{(\Cat{B},\, f)}$ 構造}\footnote{~\cite[Definition 1.1.6]{AFT2014stratified}では\textbf{$(\infty,\, 1)$-category of basics}と呼ばれている．}とは，$(\infty,\, 1)$-圏の\hyperref[def:infty-fib]{右ファイブレーション} $(\Cat{B} \xrightarrow{f} \BSC) \in (\RFIB_{\BSC})_0$ のこと．
        \item $(\Cat{B},\, f)$ 構造 $\Cat{B} \xrightarrow{f} \BSC$ を1つ固定する．このとき，\textbf{$\bm{\Cat{B}}$-多様体} ($\Cat{B}$-manifold) の成す $(\infty,\, 1)$-圏 $\MFLDB{\Cat{B}}$ とは，
        \hyperref[def:pullback-infty]{$(\infty,\, 1)$-圏の引き戻し}
        \begin{center}
            % https://tikzcd.yichuanshen.de/#N4Igdg9gJgpgziAXAbVABwnAlgFyxMJZABgBpiBdUkANwEMAbAVxiRAB12BZAMQBkAIgCFgnAMJ0cwIQF8ZIGaXSZc+QigCM5KrUYs2ACk4AlHgEkhAfVHshAZTEyAlNfGTpMgAScAHgCcsAHMACxw6Pz8IAHdgNC9Oe0cFJRAMbDwCIi0NHXpmVkQOdlMLV1sHeUVldLUiMhzqPP1CzjsAOQBxPmMFHRgoQPgiUAAzSIBbJABmahwIJAAmRr0CorCmEGoGOgAjGAYABRUM9RAAkJxk0YmkMhA56aqQMYhJxC17+cQl3Xy2TnGkmCfnGwBGED8gxwlRSLzedwe7xkFBkQA
            \begin{tikzcd}
            \MFLDB{\Cat{B}} \arrow[dr, phantom," " {pullback}, very near start, color=black] \arrow[d] \arrow[r] & (\RFIB_{\BSC})_{/f} \arrow[d, "\mathrm{forget}"] \\
            \SNGLR \arrow[r, "\tau"']           & \RFIB_{\BSC}                                                              
            \end{tikzcd}
        \end{center}
        のこと
        \footnote{実際には\exref{def:infty-forget}より $(\RFIB_{\BSC})_{/f} \lto \RFIB_{\BSC}$ が\hyperref[def:infty-fib]{右ファイブレーション}（したがって図\ref{fig:infty-fib}より\hyperref[def:isofibration]{アイソファイブレーション}）なので，補題\ref{lem:pullback-infty}, \ref{lem:pullback-iso}からこの $(\infty,\, 1)$-圏の引き戻しは $(1,\, 1)$-圏 $\sSet$ の引き戻しとして計算できる．}．
        特に，$\MFLDB{\Cat{B}}_0$ の元は以下の2つのデータから成り，\textbf{$\bm{\Cat{B}}$-多様体}と呼ばれる：
        \begin{itemize}
            \item \hyperref[def:c-smooth]{conically smoothな層状化空間} $X \in \SNGLR_0$
            \item \hyperref[def:infinity-1]{$(\infty,\, 1)$-圏の関手} $g \colon \ENTR{X} \lto \Cat{B}$
        \end{itemize}
        これらは以下の条件を充たさねばならない：
        \begin{description}
            \item[\textbf{(lift of tangent classifier)}]　
            
            $\sSet$ における以下の図式は可換である：
            \begin{center}
                % https://tikzcd.yichuanshen.de/#N4Igdg9gJgpgziAXAbVABwnAlgFyxMJZABgBoBGAXVJADcBDAGwFcYkQAdDgUQDkAVAErAAGgF8QY0uky58hFOQrU6TVuy4AhAMoBhSdJAZseAkSXEVDFm0ScOu+jmCaJYlTCgBzeEVAAzACcIAFskMhAcCCQAJhprdTsvAwDgsMQIqKQlVRsNDhx6ZgB9ERAaRnoAIxhGAAVZUwUQQKwvAAscFJAg0NiaLMQchNseyUoxIA
                \begin{tikzcd}
                                                            & \Cat{B} \arrow[d, "f"] \\
                \ENTR{X} \arrow[ru, "g"] \arrow[r, "\tau_X"'] & \BSC                  
                \end{tikzcd}
            \end{center}
        \end{description}
    \end{itemize}
\end{mydef}

% \subsection{$C^\infty$-多様体の接構造との比較}
以降の小節で定義\ref{def:B-mfld}を解読しよう．
特に，定理\ref{thm:tangentclassifier}により，\hyperref[def:tangent-classifier]{tangent classifier} $\tau_X$ は\hyperref[def:c-smooth]{conically smoothな層状化空間} $X \in \SNGLR_0$ の上に「stratum-wise」に（通常の意味の）接束をアサインしていることが分かる．
ここで言う「stratum-wise」とは，$X$ が元来持っている\hyperref[def:stratified-space]{層状化}のことではなく，補題\ref{lem:strat-basic}による $X$ の\hyperref[thm:basic-easy]{$[\BSC]$}-層状化によるものである．
補題\ref{lem:XU-mfd}により，\hyperref[thm:basic-easy]{$[\BSC]$}-層状化による $X$ の $[U] \in [\BSC]$-層はただの $C^\infty$-多様体であることが分かり，その上の\hyperref[prop:tangentbundle]{接束}を考えることができるのである．

\subsection{{$\mathcal{B}\mathrm{sc}$} におけるHom}

しばらくの間，\hyperref[def:Snglr]{basic} $U \coloneqq \bigl( (\mathbb{R}^n \to [0]) \times \Cone{Z \to P},\, \mathcal{A}_Z \bigr) \in \Obj{\Bsc}$ を1つ固定する．
これまでと同様に，\hyperref[def:stratified-space]{層状化空間}$\Cone{Z \xrightarrow{s} P} = \bigl(\Cone{Z} \xrightarrow{\Cone{s}} \Cone{P}\bigr)$ の\hyperref[def:str-cone]{コーンポイント} $[0,\, *] \in \Cone{Z} = \frac{\mathbb{R}_{\ge 0} \times Z}{\{0\} \times Z}$ のことを $\mathrm{pt}$ と書き，$\mathbb{R}^n \times \{\mathrm{pt}\} \subset \mathbb{R}^n \times \Cone{Z}$ のことを $\mathbb{R}^n$ と略記することにする．
さらに，点 $(0,\, \mathrm{pt}) \in \mathbb{R}^n \times \Cone{Z}$ のことを $\bm{0}$ と略記し，$U$ の\textbf{原点} (origin) と呼ぶことにする．

\begin{marker}
    以下では，混乱が生じにくい場合は層状化空間 $(X \xrightarrow{s} P)$ の $s$ と $P$ を省略する．さらに，\hyperref[def:Sm-induction]{conically smooth atlas}を明示しない．
\end{marker}

いま，\hyperref[def:simplicial-top]{特異単体}の変種として，\hyperref[def:c-smooth]{conically smoothな層状化空間} $X \in \Obj{\Strat}$ の\textbf{滑らかな特異単体} (smooth singular simplicial set) と呼ばれる単体的集合を
\begin{align}
    \Sing^{\text{sm}}(X) \colon \OP{\Delta} &\lto \SETS, \\
    [n] &\lmto \Hom{\Strat}(\textcolor{red}{\Delta^n_e},\, X), \\
    \bigl([m] \xrightarrow{\alpha} [n]\bigr) &\lmto \bigl( \Sing^{\text{sm}}(X)_m \xrightarrow{\Delta_e(\alpha)^*} \Sing^{\text{sm}}(X)_n \bigr)
\end{align}
と定義する\footnote{滑らかな特異単体と言ったときは\hyperref[def:simplicial-top]{幾何学的 $n$-単体}を用いることが多く，このように定義することは稀だと思う．}．これはKan複体になる．
ここで\eqref{eq:auto-basic}を思い出して，写像
\begin{align}
    \tilde{\gamma} \colon \mathbb{R}_{\ge 0} \times \mathbb{R}^n &\lto \Hom{\Bsc} (U,\, U), \\
    (t,\, p) &\lmto \Bigl( \bigl( v,\, [s,\, z] \bigr) \mapsto \bigl( tv+p,\, [ts,\, z] \bigr) \Bigr) \eqqcolon \tilde{\gamma}_{t,\, p}
\end{align}
を考える．これをKan複体の間の射（i.e. 自然変換）
\begin{align}
    \label{eq:gamma-Kan}
    \gamma \colon \Sing^{\text{sm}} (\mathbb{R}_{\ge 0} \times \mathbb{R}^n) &\lto \Hom{\BSC} (U,\, U)
\end{align}
へと格上げすることができる．
実際，$\forall m \ge 0$ に対して，$m$-単体はそれぞれ
\begin{align}
    \Sing^{\text{sm}} (\mathbb{R}_{\ge 0} \times \mathbb{R}^n)_m &= \Hom{\Strat} (\Delta_e^m,\, \mathbb{R}_{\ge 0} \times \mathbb{R}^n), \\
    \Hom{\BSC} (U,\, U)_m &= \bigl\{\, f \in \Hom{\Snglr}(\Delta^m_e \times U,\, \Delta^m_e \times U) \bigm| \mathrm{proj}_{\Delta^m_e} \circ f = \mathrm{proj}_{\Delta^m_e} \,\bigr\}
\end{align}
であったから，
\begin{align}
    \gamma_{[m]} \colon \Sing^{\text{sm}} (\mathbb{R}_{\ge 0} \times \mathbb{R}^n)_m &\lto \Hom{\BSC} (U,\, U)_m, \\
    \Bigl( x \mapsto \bigl( t(x),\, p(x) \bigr)  \Bigr)  &\lmto \Bigl( \bigl(x,\, u\bigr) \mapsto \bigl(x,\, \tilde{\gamma}_{t(x),\, p(x)}(u)\bigr) \Bigr) \eqqcolon \gamma_{[m]}{}_{t,\, p}
\end{align}
と定義すればよい．

さらに，勝手な\hyperref[def:Snglr]{conically smoothな層状化空間} $Z \in \Obj{\Snglr}$ に対して，その\textbf{自己同相群} (automorphism group) 
\begin{align}
    \Aut(Z) \colon \OP{\Delta} \lto \SETS
\end{align}
を，\hyperref[lem:Kan-enriched]{Kan複体 $\Hom{\SNGLR} (Z,\, Z)$} の\hyperref[def:SimpSet]{部分Kan複体}として次のように定義する：
\begin{align}
    \label{def:Kan-Aut}
    \Aut(Z)_m \coloneqq \bigl\{\, f \in \Hom{\SNGLR} (Z,\,Z)_m \bigm| \forall x \in \Delta^m_e,\; \exists (\mathrm{proj}_{Z} \circ f|_{\{x\} \times Z})^{-1} \in \Hom{\Snglr}(Z,\, Z) \,\bigr\} 
\end{align}
なお，$\forall m \ge 0$ に対し集合 $\Aut(Z)_m$ は $(1,\, 1)$-圏 $\Snglr$ における1-射の合成に関して群になるため，Kan複体 $\Aut(Z)$ は\textbf{群的なKan複体} (group-like Kan complex)\footnote{Kan複体であって，$\OP{\Delta} \lto \GRP$ でもあるもの．従って，homotopy hypothesisより位相群と見做すことができる．}と見做すことができる．\label{def:Kan-group}

\begin{mydef}[label=def:stratified-GL]{層状化された一般線形群}
    勝手な\hyperref[def:Snglr]{basic} $U \coloneqq \mathbb{R}^n \times \Cone{Z} \in \Obj{\Bsc}$ を与える．

    \hyperref[lem:Kan-enriched]{Kan複体 $\Hom{\SNGLR} (U,\, U)$} の部分Kan複体
    \footnote{より正確には，この時点では\hyperref[def:SimpSet]{部分単体的集合} $\LGL(U) \subset \Hom{\SNGLR}(U,\, U)$ がKan複体であることは非自明である．補題\ref{lem:Hom-sub-1}により $\LGL(U)$ が群的な単体的集合（単体的群）こと，i.e. $\LGL(U) \colon \OP{\Delta} \lto \GRP$ であることが分かり，故に\cite[\href{https://kerodon.net/tag/00MG}{Tag 00MG}]{kerodon}からKan複体であることが分かるのである．}
    \begin{align}
        \LGL (U) \colon \OP{\Delta} \lto \SETS
    \end{align}
    を次のように定義し，\textbf{層状化された一般線形群}と呼ぶ：
    \small
    \begin{align}
        \bm{\LGL(U)_m} \coloneqq \bigl\{\, T \in \Hom{\SNGLR} (U,\, U)_m \bigm| \forall (t,\, p) \in \Sing^{\text{sm}}(\mathbb{R}_{\ge 0} \times \mathbb{R}^n)_m,\; T \circ \gamma_{[m]}{}_{t,\, p} = \gamma_{[m]}{}_{t,\, (Tp)_{\mathbb{R}^n}} \circ T  \,\bigr\} 
    \end{align}
    \normalsize
    ただし，以下の合成写像を $(Tp)_{\mathbb{R}^n} \colon \Delta^m_e \lto \mathbb{R}^n$ と略記した：
    \scriptsize
    \begin{align}
        &\Delta^m_e & &\xrightarrow{\mathrm{diag}} & &\Delta^m_e \times \Delta^m_e & &\xrightarrow{\Id \times p} & &\Delta^m_e \times \mathbb{R}^n & &\xrightarrow{T|_{\Delta^m_e \times \mathbb{R}^n \times \{\mathrm{pt}\}}}  & &\Delta^m_e \times \mathbb{R}^n \times \Cone{Z} & &\xrightarrow{\mathrm{proj}_{\mathbb{R}^n}} \mathbb{R}^n, \\
        &x & &\lmto & &(x,\, x) & &\lmto & &\bigl( x,\, p(x) \bigr) & &\lmto & &T \bigl( x,\, p(x),\, \mathrm{pt} \bigr) & &\lmto \mathrm{proj}_{\mathbb{R}^n} \circ T\bigl( x,\, p(x),\, \mathrm{pt} \bigr)
    \end{align}
    \normalsize
    
    \tcblower

    \hyperref[lem:Kan-enriched]{Kan複体 $\Hom{\SNGLR} (U,\, U)$} の部分Kan複体
    \begin{align}
        \LO (U) \colon \OP{\Delta} \lto \SETS
    \end{align}
    を次のように定義し，\textbf{層状化された直交群}と呼ぶ：
    \begin{align}
        \LO (U) \coloneqq \Sing^{\text{sm}}\bigl(\LO(n)\bigr) \times \Aut(Z)
    \end{align}
\end{mydef}

\begin{myexample}[label=ex:ordinary-GL]{通常の一般線形群}
    $Z = \emptyset$ の場合を考える．このとき $U = \mathbb{R}^n$ なので，Kan複体の射\eqref{eq:gamma-Kan}とは単に
    \begin{align}
        \gamma_{[m]} \colon \Bigl( x \mapsto \bigl( t(x),\, p(x) \bigr)  \Bigr)  &\lmto \Bigl( \bigl(x,\, u\bigr) \mapsto \bigl(x,\, t(x)u + p(x)\bigr) \Bigr)
    \end{align}
    のことであり，$T \in \LGL(U)_m$ とは $\forall (t,\, p) \in \Sing^{\text{sm}}(\mathbb{R}_{\ge 0} \times \mathbb{R}^n)$ および
    $\forall (x,\, u) \in \Delta^m_e \times U = \Delta^m_e \times \mathbb{R}^n$ に対して
    \begin{align}
        T \circ \gamma_{[m]}{}_{t,\, p}(x,\, u)
        &= T \bigl(x,\, t(x)u + p(x)\bigr) \\
        &= \gamma_{[m]}{}_{t,\, \mathrm{proj}_{\mathbb{R}^n} \circ T|_{\mathbb{R}^n} (p)} \circ T (x,\, u) \\
        &= \Bigl(x,\, t(x)\, \mathrm{proj}_{\mathbb{R}^n} \circ T(x,\, u) + \mathrm{proj}_{\mathbb{R}^n} \circ T\bigl(x,\, p(x)\bigr)  \Bigr) 
    \end{align}
    が成り立つことを意味する．i.e. $\forall x \in \Delta^m_e$ に対して，$T_x \coloneqq \mathrm{proj}_{\mathbb{R}^n} \circ T|_{\{x\} \times \mathbb{R}^n} \colon \mathbb{R}^n \lto \mathbb{R}^n$ とおくと
    \begin{align}
        T_x \bigl( t(x) u + p(x) \bigr) = t(x)\, T_x u + T_x p(x)
    \end{align}
    が成り立つため，$T_x$ は線型写像である．
    \hyperref[def:Snglr]{$(1,\, 1)$-圏 $\Snglr$ の定義}より $T$ は開埋め込みであり，単射である．次元の等しい有限次元ベクトル空間の間の単射線型写像 $T_x$ は全単射でもあり，$T_x \in \LGL(n,\, \mathbb{R})$ だと分かった．
    この事実は，Kan複体の同型 $\LGL(\mathbb{R}^n) \cong \Sing^{\text{sm}}\bigl(\LGL(n,\, \mathbb{R})\bigr)$ を意味する．
\end{myexample}

さて，部分Kan複体 $\Aut^0 (U) \subset \Hom{\BSC}^0 (U,\, U) \subset \Hom{\SNGLR} (U,\, U)$ を次のように定義しよう：
\begin{align}
    \Hom{\BSC}^0 (U,\, U)_m &\coloneqq \bigl\{\, f \in \Hom{\SNGLR} (U,\, U)_m \bigm| \forall x \in \Delta^m_e,\; f(x,\, \bm{0}) = (x,\, \bm{0}) \,\bigr\}, \\
    \Aut^0 (U)_m &\coloneqq \Aut (U)_m \cap \Hom{\BSC}^0 (U,\, U)_m
\end{align}

\begin{mylem}[label=lem:Hom-sub-1]{層状化された一般線形群は原点を保つ自己同相群}
    \begin{align}
        \LGL(U) \subset \Aut^0(U) \subset \Hom{\BSC}^0(U,\, U)
    \end{align}
    さらに，$\forall m \ge 0,\; \forall S,\, T \in \LGL(U)_m$ に対して $S \circ T \in \LGL(U)$ が成り立つ．i.e. $\LGL(U)$ は群的なKan複体である．
\end{mylem}

\begin{proof}
    $\forall m \ge 0$ および $\forall T \in \LGL(U)_m$ を1つ固定する．
    
    \begin{description}
        \item[$\bm{T} \bm{\in} \bm{\Hom{\BSC}^0(U,\, U)_m}$]　
        
        $\forall x \in \Delta^m_e$ および $\forall t \in \mathbb{R}_{\ge 0}$ に対して $\gamma_{[m]}{}_{t,\, 0} (x,\, \bm{0}) = (x,\, \bm{0})$ が成り立つ\footnote{定数写像 $\Delta^m_e \lto \mathbb{R}_{\ge 0},\, x \lmto  t$ のことを $t$ と略記している．}ので，
        \begin{align}
            \bigl( x,\, \mathrm{proj}_{U} \circ T(x,\, \bm{0}) \bigr)
            &=T(x,\, \bm{0}) \\
            &= T \circ \gamma_{[m]}{}_{t,\, 0} (x,\, \bm{0}) \\
            &= \gamma_{[m]}{}_{t,\, (T0)_{\mathbb{R}^n}} \circ T (x,\, \bm{0}) \\
            &= \gamma_{[m]}{}_{t,\, \mathrm{proj}_{\mathbb{R}^n} \circ T(x,\, \bm{0})}\bigl( x,\, \mathrm{proj}_{U} \circ T(x,\, \bm{0})\bigr) \label{eq:save-origin}
        \end{align}
        が言える\footnote{3つ目の等号で登場する $0$ とは，定数写像 $\Delta^m_e \lto \mathbb{R}^n,\, x \lmto 0$ のことである．}．
        ここで $\mathrm{proj}_{U} \circ T(x,\, \bm{0}) = \bigl( \mathrm{proj}_{\mathbb{R}^n} \circ T(x,\, \bm{0}),\, \mathrm{proj}_{\Cone{Z}} \circ T(x,\, \bm{0}) \bigr) \in U$ であるから，$v_0 \coloneqq \mathrm{proj}_{\mathbb{R}^n} \circ T(x,\, \bm{0}),\, [s_0,\, z_0] \coloneqq \mathrm{proj}_{\Cone{Z}} \circ T(x,\, \bm{0})$ とおくと\eqref{eq:save-origin}は
        \begin{align}
            \bigl(v_0,\, [s_0,\, z_0] \bigr)
            = \bigl( t v_0 + v_0,\, [ts_0,\, z_0] \bigr)
        \end{align}
        を意味する．いま $t \in \mathbb{R}_{\ge 0}$ は任意であったから，$v_0 = 0,\, s_0 = 0$ でなくてはいけない．i.e. $T(x,\, 0) = \bigl(x,\, v_0,\, [s_0,\, z_0]\bigr) = (x,\, 0)$ が示された．

        \item[$\bm{T} \bm{\in} \bm{\Aut^0(U,\, U)_m}$]　
        
        $\forall x \in \Delta^m_e$ を1つ固定する．
        \hyperref[def:Snglr]{$(1,\, 1)$-圏 $\Snglr$ の定義}より $T \in \Hom{\Snglr}(\Delta^m_e \times U,\, \Delta^m_e \times U)$ は開埋め込みであるから単射であり，従って $T_x \coloneqq \mathrm{proj}_{U} \circ T|_{\{x\} \times U} \in \Hom{\Snglr}(U,\, U)$ もまた開埋め込みであって単射である．
        故に $T_x$ が全射であること，i.e. $T_x (U) = U$ であることを示せば十分である．
        実際，原点の勝手な開近傍 $\bm{0} \in O \subset U$ を $\tilde{\gamma}_{t,\, 0}$ によって拡大することで $U$ の開基を作ることができ，かつ $T \circ \gamma_{t,\, 0} = \gamma_{t,\, 0} \circ T$ が成り立つことが先に示されているので $T_x$ は全射である~\cite[Lemma 4.3.4.]{AFT2014stratified}．

        \item[\textbf{合成について閉じていること}]　
        
        $\forall (x,\, p) \in \Delta^m_e \times \mathbb{R}^n$ に対して $(x,\, p,\, \mathrm{pt}) = \gamma_{[m]}{}_{1,\, p}(x,\, \bm{0})$ が成り立つ\footnote{定数写像 $\Delta^m_e \lto \mathbb{R}^n,\; x \lmto p$ のことを $p$ と略記している．}から，
        \begin{align}
            T(x,\, p,\, \mathrm{pt})
            &= T \circ \gamma_{[m]}{}_{1,\, p}(x,\, \bm{0}) \\
            &= \gamma_{[m]}{}_{1,\, (Tp)_{\mathbb{R}^n}} \circ T(x,\, \bm{0}) \\
            &= \gamma_{[m]}{}_{1,\, (Tp)_{\mathbb{R}^n}}(x,\, \bm{0}) \in \{x\} \times \mathbb{R}^n \times \{\mathrm{pt}\}
        \end{align}
        が言える，i.e. $\LGL(U)_{m}$ の任意の元はコーンポイントを保存することに注意する．

        　ここでもう一つの $\forall S \in \LGL(U)_m$ を固定する．
        $\forall (t,\, p) \in \Sing^{\text{sm}} (\mathbb{R}_{\ge 0} \times \mathbb{R}^n)_m$ および $\forall x \in \Delta^m_e$ に対して
        \begin{align}
            \bigl((S \circ T)p\bigr)_{\mathbb{R}^n}(x)
            &= \mathrm{proj}_{\mathbb{R}^n} \circ (S \circ T)\bigl( x,\, p(x),\, \mathrm{pt} \bigr) \\
            &= \mathrm{proj}_{\mathbb{R}^n} \circ S \Bigl(x,\, \mathrm{proj}_{\mathbb{R}^n} \circ T\bigl( x,\, p(x),\, \mathrm{pt} \bigr),\, \mathrm{proj}_{\Cone{Z}} \circ T\bigl( x,\, p(x),\, \mathrm{pt} \bigr) \Bigr) \\
            &= \mathrm{proj}_{\mathbb{R}^n} \circ S \Bigl(x,\, (Tp)_{\mathbb{R}^n}(x),\, \mathrm{pt} \Bigr) \\
            &= \bigl( S(Tp)_{\mathbb{R}^n} \bigr)_{\mathbb{R}^n}(x)
        \end{align}
        と計算できるため，
        \begin{align}
            (S \circ T) \circ \gamma_{[m]}{}_{t,\, p}
            &= S \circ \gamma_{[m]}{}_{t,\, (Tp)_{\mathbb{R}^n}} \circ T \\
            &= \gamma_{[m]}{}_{t,\, \bigl(S(Tp)_{\mathbb{R}^n}\bigr)_{\mathbb{R}^n}} \circ S \circ T \\
            &= \gamma_{[m]}{}_{t,\, \bigl((S\circ T)p\bigr)_{\mathbb{R}^n}} \circ (S \circ T)
        \end{align}
        i.e. $S \circ T \in \LGL(U)_m$ が言えた．
    \end{description}
\end{proof}

\begin{mylem}[label=lem:Hom-sub-2]{層状化された直交群は一般線形群の部分Kan複体}
    \begin{align}
        \LO(U) \subset \LGL(U)
    \end{align}
\end{mylem}

\begin{proof}
    $\forall m \ge 0$ および $\forall T \coloneqq (A,\, f) \in \LO(U)_m = \Sing^{\text{sm}}\bigl( \LO(n) \bigr)_m \times \Aut(Z)_m$ を1つ固定する．
    まず，$\forall \bigl(x,\, (v,\, [s,\, z]) \bigr) \in \Delta^m_e \times U = \Delta^m_e \times \mathbb{R}^n \times \Cone{Z}$ に対して
    \begin{align}
        T\bigl(x,\, (v,\, [s,\, z]) \bigr)
        &\coloneqq \Bigl( x,\, \bigl(A(x)v,\, [s,\, \mathrm{proj}_Z \circ f(x,\, z)]\bigr) \Bigr) 
    \end{align}
    と定めることで $T \in \Hom{\BSC}^0 (U,\, U)_m$ と見做せる．このとき，$\forall (t,\, p) \in \Sing^{\text{sm}} (\mathbb{R}_{\ge 0} \times \mathbb{R}^n)_m$ に対して
    \begin{align}
        (Tp)_{\mathbb{R}^n} (x)
        &= \mathrm{proj}_{\mathbb{R}^n} \circ T \Bigl(x,\, \bigl( p(x),\, [0,\, *] \bigr)  \Bigr) \\
        &= \mathrm{proj}_{\mathbb{R}^n}\Bigl( x,\, \bigl(A(x)p(x),\, [0,\, *]\bigr) \Bigr) \\
        &= A(x) p(x)
    \end{align}
    および
    \begin{align}
        T \circ \gamma_{[m]}{}_{t,\, p}\bigl(x,\, (v,\, [s,\, z]) \bigr)
        &= T \Bigl(x,\, \bigl( t(x)v+p(x),\, [t(x)s,\, z] \bigr)\Bigr) \\
        &= \Bigl(x,\, \bigl(t(x) A(x)v + A(x)p(x),\, [t(x)s,\, \mathrm{proj}_Z \circ f(x,\, z)]\bigr) \Bigr) \\
        &= \Bigl(x,\, \tilde{\gamma}_{t(x),\, A(x)p(x)}\bigl( A(x)v,\, [s,\, \mathrm{proj}_Z \circ f(x,\, z)] \bigr) \Bigr) \\
        &= \gamma_{[m]}{}_{t,\, (Tp)_{\mathbb{R}^n}} \circ T \bigl(x,\, (v,\, [s,\, z]) \bigr)
    \end{align}
    が成り立つため $T \in \LGL(U)_m$ が言えた．
\end{proof}

\begin{myprop}[label=prop:Hom-htpy-equiv]{$\LO(U)$ へのレトラクション}
    Kan複体の包含写像
    \begin{align}
        \LO(U) \overset{i_1}{\hookrightarrow} \LGL(U) \overset{i_2}{\hookrightarrow} \Hom{\BSC}^0 (U,\, U) \overset{i_3}{\hookrightarrow} \Hom{\BSC}(U,\, U)
    \end{align}
    は全てKan複体のホモトピー同値である．
\end{myprop}

\begin{proof}
    $\forall A \in \LGL(n,\, \mathbb{R}^n)$ を1つとる．$A$ を列ベクトルに分割し，それを $A = [a_1,\, \dots,\, a_n]$ と書こう．
    このとき，正規直交基底 $e_1,\, \dots,\, e_n \in \mathbb{R}^n$ と上三角行列 $R_A \in \Mat{n}{\mathbb{R}}$ を次のようにして得る（Gram-Schmidtの正規直交化）：
    \begin{align}
        e_i &\coloneqq 
        \begin{cases}
            \frac{a_1}{\norm{a_1}}, &i=1 \\
            \frac{a_i - \sum_{k=1}^{i-1} \braket{e_k}{a_i} e_k}{\norm{a_i - \sum_{k=1}^{i-1} \braket{e_k}{a_i} e_k}}, &1 < i \le n
        \end{cases}
        \\
        R_A{}_{ij} &\coloneqq 
        \begin{cases}
            \norm{a_1}, &i=j=1 \\
            \norm{a_i - \sum_{k=1}^{i-1} \braket{e_k}{a_i} e_k}, &1 < i=j \le n \\
            \braket{e_i}{a_j}, &i < j \\
            0, &i > j
        \end{cases}
    \end{align}
    ただし $\norm{\cdot}$ は $\mathbb{R}^n$ の標準的なノルムである．
    ここで $Q_A \coloneqq [e_1,\, \dots,\, e_n]$ とおくと $A = Q_A R_A$ が成り立つ．
    連続写像
    \begin{align}
        \label{def:GrSm}
        \mathrm{GrSm} \colon \LGL(n,\, \mathbb{R}) &\lto \LO(n), \\
        A &\lmto Q_A
    \end{align}
    は，包含写像 $i_0 \colon \LO(n) \hookrightarrow \LGL(n)$ について $\mathrm{GrSm} \circ i_0 = \Id_{\LO(m)}$ を充たす．

    \begin{description}
        \item[\textbf{$i_1$ のホモトピー逆}]　
        
        % Gram-Schmidtの正規直交化を行う．
        % 具体的には，まずKan複体の射
        % \begin{align}
        %     r_1 \colon \LGL(U) \lto \LO(U)
        % \end{align}
        % を次のように定める：
        $\forall T \in \LGL(U)_m$ および $\forall \bigl( x,\, (v,\, [0,\, *]) \bigr) \in \Delta^m_e \times \mathbb{R}^n \times \Cone{Z}$ をとる．
        補題\ref{lem:Hom-sub-1}の証明より $T$ はコーンポイントを保存するから，\exref{ex:ordinary-GL}と同様にして $\forall (t,\, p) \in \Sing^{\text{sm}} (\mathbb{R}_{\ge 0},\, \mathbb{R}^n)_m$ に対して
        \begin{align}
            &\textcolor{red}{T} \Bigl( \textcolor{red}{x},\, \bigl( t(x)v + p(x),\, \textcolor{red}{[0,\, *]} \bigr) \Bigr)  \\
            &= T \circ \gamma_{[m]}{}_{t,\, p} \bigl( x,\, (v,\, [0,\, *]) \bigr) \\
            &= \gamma_{[m]}{}_{t,\, (Tp)_{\mathbb{R}^n}} \circ T\bigl( x,\, v,\, [0,\, *] \bigr) \\
            &= \biggl( x,\, \Bigl( t(x)\, \textcolor{red}{\mathrm{proj}_{\mathbb{R}^n} \circ T} \bigl( \textcolor{red}{x},\, v,\, \textcolor{red}{[0,\, *]} \bigr) + \textcolor{red}{\mathrm{proj}_{\mathbb{R}^n} \circ T} \bigl( \textcolor{red}{x},\, p(x),\, \textcolor{red}{[0,\, *]} \bigr),\, [0,\, *]  \Bigr)  \biggr) 
        \end{align}
        と計算できる．i.e. $T_x \coloneqq \mathrm{proj}_{\mathbb{R}^n} \circ T|_{\{x\} \times \mathbb{R}^n \times \{[0,\, *]\}} \colon \mathbb{R}^n \lto \mathbb{R}^n$ は線型写像であり，$T$ が開埋め込みであることから $T_x \in \LGL(n,\, \mathbb{R})$ が言える．
        故に連続写像\eqref{def:GrSm}を $T_x$ に適用することができて，$\mathrm{GrSm}(T_x) \in \LO(n)$ を得る．
        
        　一方，$\forall \bigl( x,\, (v,\, [s,\, z]) \bigr) \in \Delta^m_e \times \mathbb{R}^n \times \Cone{Z}$ に対して
        $\bigl( x,\, (v,\, [s,\, z]) \bigr) = \gamma_{[m]}{}_{s,\, v} \bigl( x,\, (0,\, [1,\, z]) \bigr)$ が成り立つ\footnote{$(s,\, v) \in \Sing^{\text{sm}} (\mathbb{R}_{\ge 0} \times \mathbb{R}^n)_m$ とは，定数写像 $\Delta^m_e \lto \mathbb{R}_{\ge 0} \times \mathbb{R}^n,\; x \lmto (s,\, v)$ のことである．}ので，
        \begin{align}
            &T \bigl( x,\, (v,\, [s,\, z]) \bigr) \\
            &= T \circ \gamma_{[m]}{}_{s,\, v} \bigl( x,\, (0,\, [1,\, z]) \bigr) \\
            &= \gamma_{[m]}{}_{s,\, (Tv)_{\mathbb{R}^n}} \circ T \bigl( x,\, (0,\, [1,\, z]) \bigr)
        \end{align}
        が言える．故に，
        \begin{align}
            \mathrm{proj}_{U} \circ T \bigl( x,\, (v,\, [s,\, z]) \bigr) \eqqcolon \Bigl(T^{\mathbb{R}^n}_{x,\, v,\, [s,\, z]},\, \bigl[ T^{\mathbb{R}_{\ge 0}}_{x,\, v,\, [s,\, z]},\, T^Z_{x,\, v,\, [s,\, z]} \bigr] \Bigr)
        \end{align}
        とおくと
        \begin{align}
            \bigl[T^{\mathbb{R}_{\ge 0}}_{x,\, v,\, [s,\, z]},\, T^Z_{x,\, v,\, [s,\, z]}\bigr] &= \bigl[ s\, T^{\mathbb{R}_{\ge 0}}_{x,\, 0,\, [1,\, z]},\, T^Z_{x,\, 0,\, [1,\, z]}\bigr]
        \end{align}
        が成り立つ．

        　以上の考察から，写像
        \begin{align}
            r_1{}_{[m]} \colon \LGL(U)_m &\lto \LO(U)_m, \\
            T &\lmto \Bigl( (x,\, v,\, [s,\, z]) \mapsto \bigl( x,\, \mathrm{GrSm}(T_x)(v),\, [s,\, T^Z_{x,\, 0,\, [1,\, z]}] \bigr)  \Bigr) 
        \end{align}
        が定まり，Kan複体の射
        \begin{align}
            r_1 \colon \LGL(U) \lto \LO(U)
        \end{align}
        を構成する．構成より明らかに $r_1 \circ i_1 = \Id_{\LO(U)}$ が成り立つ．その上，Kan複体の射
        \begin{align}
            H \in \FUN \bigl(\LGL(U),\, \LGL(U)\bigr)_1 \cong \Hom{\Kan} \bigl(\Sing^{\text{sm}}([0,\, 1]) \times \LGL(U),\, \LGL(U)\bigr)
        \end{align}
        の $m$-射の対応 
        $H_m \colon \Sing^{\text{sm}}([0,\, 1])_m \times \LGL(U)_m \lto \LGL(U)_m$ 
        を次のように定義すると，$H$ が $\Id_{\LGL(U)}$ と $i_1 \circ r_1$ を繋ぐ\hyperref[def:infty-homotopy-morphism]{ホモトピー}となる．i.e. $r_1$ が $i_1$ のホモトピー逆であることが示された：
        \footnotesize
        \begin{align}
            % \Sing^{\text{sm}}([0,\, 1]) \times \LGL(U) &\lto \LGL(U), \\
            (t,\, T) &\lmto \biggl( (x,\, v,\, [s,\, z]) \mapsto \Bigl( x,\, \bigl(1-t(x)\bigr)\, T^{\mathbb{R}^n}_{x,\, v,\, [s,\, z]} + t(x) \mathrm{GrSm}(T_x)(v),\, \Bigl[ \bigl(1-t(x)\bigr) T^{\mathbb{R}_{\ge 0}}_{x,\, 0,\, [1,\, z]} + t(x)s,\, T^Z_{x,\, 0,\, [1,\, z]}\Bigr] \Bigr)  \biggr) 
        \end{align}
        \normalsize

        \item[\textbf{$i_2$ のホモトピー逆}]　
        
        $\forall f \in \Hom{\BSC}^0(U,\, U)_m$ および $\forall x \in \Delta^m_e$ に対して，$f_x \coloneqq \mathrm{proj}_U \circ f|_{\{x\} \times U} \in \Hom{\Snglr}^0 (U,\, U)$ とおこう．
        このとき $(\infty,\, 1)$-圏 $\BSC$ の定義において登場する射は全て\hyperref[def:c-smooth-map]{conically smooth}なので，$f_x$ の\hyperref[def:c-smooth-along]{微分}
        \begin{align}
            Df_x \colon T\mathbb{R}^n \times \Cone{Z} &\lto T\mathbb{R}^n \times \Cone{Z}
        \end{align}
        が存在する．具体的には，
        \small
        \begin{align}
            &Df_x \bigl((v,\, p),\, [t,\, z]\bigr) \\
            &= \lim_{s \to +0} \gamma^{-1} \circ f_x{}_{\Delta} \circ \gamma \bigl(s;\, (v,\, p),\, [t,\, z]\bigr) \\
            &= \lim_{s \to +0} \gamma^{-1} \circ f_x{}_{\Delta} \bigl(s;\, (sv+p,\, p),\, [st,\, z]\bigr) \\
            &= \lim_{s \to +0} \gamma^{-1} \bigl(s;\, f_x|_{\mathbb{R}^n}(sv+p),\, f_x(p,\, [st,\, z])\bigr) \\
            &= \left( \Bigl(\lim_{s \to +0} \frac{f_x|_{\mathbb{R}^n}(sv+p) - \mathrm{proj}_{\mathbb{R}^n} \circ f_x(p,\, [st,\, z])}{s},\, f_x|_{\mathbb{R}^n}(p) \Bigr),\, \mathrm{proj}_{\Cone{Z}} \Bigl(\lim_{s \to +0}\tilde{\gamma}_{\frac{1}{s},\, -\frac{\mathrm{proj}_{\mathbb{R}^n} \circ f_x(p,\, [st,\, z])}{s}} \circ f_x(p,\, [st,\, z])\Bigr) \right) \\
            &= \left( \Bigl(\lim_{s \to +0} \frac{f_x|_{\mathbb{R}^n}(sv+p) - \mathrm{proj}_{\mathbb{R}^n} \circ f_x(p,\, [st,\, z])}{s},\, f_x|_{\mathbb{R}^n}(p) \Bigr),\, \mathrm{proj}_{\Cone{Z}} \Bigl(\lim_{s \to +0}\tilde{\gamma}_{\frac{1}{s},\, 0} \circ f_x(p,\, [st,\, z])\Bigr) \right)
        \end{align}
        \normalsize
        である\footnote{$\tilde{\gamma}$ の定義から，$\forall (t,\, p) \in \mathbb{R}_{\ge 0} \times \mathbb{R}^n$ に対して $\mathrm{proj}_{\Cone{Z}} \circ \tilde{\gamma}_{t,\, p} = \mathrm{proj}_{\Cone{Z}} \circ \tilde{\gamma}_{t,\, 0}$ が成り立つ．}．
        さらに，$T \mathbb{R}^n \approx \mathbb{R}^n \times \mathbb{R}^n$ の底空間\footnote{第二成分である．}の勝手な元 $p \in \mathbb{R}^n$ に対して
        % $T(\Delta^m_e \times \mathbb{R}^n) \approx (\Delta^m_e \times \mathbb{R}^n) \times (\Delta^m_e \times \mathbb{R}^n)$ であり，底空間の勝手な元 $(x,\, p) \in \Delta^m_e \times \mathbb{R}^n$ に対して
        \begin{align}
            D_{p}f_x &\coloneqq \mathrm{proj}_{\mathbb{R}^n \times \{f_x|_{\mathbb{R}^n} (p)\} \times \Cone{Z}} \circ Df_x|_{\mathbb{R}^n \times \{ p \} \times \Cone{Z}} \in \Hom{\Snglr} (U,\, U), \\
            D_{p}{}_{[m]}f &\coloneqq \Bigl( \bigl(x,\, u\bigr) \lmto \bigl(x,\, D_{p}f_x(u)\bigr) \Bigr) \in \Hom{\BSC} (U,\, U)_m
        \end{align}
        とおこう．特に，$D_{0}{}_{[m]} f \in \Hom{\BSC}^0(U,\, U)_m$ であることに注意する．
        あからさまには
        \small
        \begin{align}
            &D_0 f_x (v,\, [t,\, z]) \\
            &= \left( \lim_{s \to +0} \frac{f_x|_{\mathbb{R}^n}(sv) - \mathrm{proj}_{\mathbb{R}^n} \circ f_x(0,\, [st,\, z])}{s},\, \mathrm{proj}_{\Cone{Z}} \Bigl(\lim_{s \to +0}\tilde{\gamma}_{\frac{1}{s},\, 0} \circ f_x(0,\, [st,\, z])\Bigr) \right) \\
            &= \left( \lim_{s \to +0} \frac{f_x|_{\mathbb{R}^n}(sv) - \mathrm{proj}_{\mathbb{R}^n} \circ f_x(\bm{0})}{s},\, \mathrm{proj}_{\Cone{Z}} \Bigl(\lim_{s \to +0}\tilde{\gamma}_{\frac{1}{s},\, 0} \circ f_x(0,\, [st,\, z])\Bigr) \right) \\
            &= \left( \lim_{s \to +0} \frac{f_x|_{\mathbb{R}^n}(sv)}{s},\, \mathrm{proj}_{\Cone{Z}} \Bigl(\lim_{s \to +0}\tilde{\gamma}_{\frac{1}{s},\, 0} \circ f_x(0,\, [st,\, z])\Bigr) \right) \\
            &= \lim_{s \to +0} \tilde{\gamma}_{\frac{1}{s},\, 0} \circ f_x{}_{\Delta} \circ \tilde{\gamma}_{s,\, 0} (v,\, [t,\, z])
        \end{align}
        \normalsize
        と書ける．

        　以上の準備の下，まず $\forall m \ge 0$ および $\forall f \in \Hom{\BSC}^0(U,\, U)_m$ に対して $D_{0}{}_{[m]} f \in \LGL(U)_m$ であること，i.e. Kan複体の射
        \begin{align}
            D_{0} \colon \Hom{\BSC}^0(U,\, U) \lto \LGL(U)
        \end{align}
        が定まることを示そう．
        実際，$\forall (t,\, p) \in \Sing^{\text{sm}} (\mathbb{R}_{\ge 0} \times \mathbb{R}^n)_m,\, \forall (x,\, u) \in \Delta^m_e \times U$ に対して
        \begin{align}
            (D_{0}{}_{[m]} f(p))_{\mathbb{R}^n} (x) 
            &= \mathrm{proj}_{\mathbb{R}^n} \circ D_{0}{}_{[m]} f \bigl( x,\, p(x),\, \mathrm{pt} \bigr) \\
            &= \mathrm{proj}_{\mathbb{R}^n} \circ D_0 f_x \bigl( p(x),\, \mathrm{pt} \bigr) \\
            &= \lim_{s \to +0} \frac{f_x|_{\mathbb{R}^n} \bigl(sp(x)\bigr)}{s}
        \end{align}
        であるから，
        \begin{align}
            &D_{0}{}_{[m]} f \circ \gamma_{[m]}{}_{t,\, p} (x,\, u) \\
            &= \Bigl( x,\, D_{0}f_x \circ \tilde{\gamma}_{t(x),\, p(x)}(u) \Bigr) \\
            &= \Bigl( x,\, \bigl(\lim_{s \to +0} \tilde{\gamma}_{\frac{1}{s},\, 0} \circ f_x{}_\Delta \circ \tilde{\gamma}_{s,\, 0} \bigr) \circ \tilde{\gamma}_{t(x),\, p(x)}(u) \Bigr) \\
            &= \Bigl( x,\, \lim_{s \to +0} \tilde{\gamma}_{\frac{1}{s},\, 0} \circ f_x{}_\Delta \circ \tilde{\gamma}_{st(x),\,sp(x)} (u) \Bigr) \\
            &= \Bigl( x,\, \bigl( \lim_{s \to +0} \bigl(\tilde{\gamma}_{1,\, \frac{f_x|_{\mathbb{R}^n} (sp(x))}{s}}\circ \tilde{\gamma}_{1,\, - \frac{f_x|_{\mathbb{R}^n} (sp(x))}{s}}\bigr) \circ \tilde{\gamma}_{\frac{1}{s},\, 0} \circ f_x{}_{\Delta} \circ \tilde{\gamma}_{st(x),\, sp(x)} \bigr)(u) \Bigr)  \\
            &= \Bigl( x,\, \bigl( \lim_{r \to +0} \tilde{\gamma}_{1,\, \frac{f_x|_{\mathbb{R}^n} (rp(x))}{r}} \bigr) \circ \bigl(\lim_{s \to +0} \tilde{\gamma}_{1,\, - \frac{f_x|_{\mathbb{R}^n} (sp(x))}{s}} \circ \tilde{\gamma}_{\frac{1}{s},\, 0} \circ f_x{}_{\Delta} \circ \tilde{\gamma}_{st(x),\, sp(x)} \bigr)(u) \Bigr)  \\
            &= \Bigl( x,\, \bigl( \lim_{r \to +0} \tilde{\gamma}_{1,\, \frac{f_x|_{\mathbb{R}^n} (rp(x))}{r}} \bigr) \circ \bigl(\lim_{s \to +0} \lim_{a \to +0} \tilde{\gamma}_{1,\, - \frac{f_x|_{\mathbb{R}^n} (ap(x))}{s}} \circ \tilde{\gamma}_{\frac{1}{s},\, 0} \circ f_x{}_{\Delta} \circ \tilde{\gamma}_{st(x),\, ap(x)} \bigr)(u) \Bigr)  \\
            &= \Bigl( x,\, \bigl( \lim_{r \to +0} \tilde{\gamma}_{1,\, \frac{f_x|_{\mathbb{R}^n} (rp(x))}{r}} \bigr) \circ \bigl(\lim_{s \to +0} \tilde{\gamma}_{\frac{1}{s},\, 0} \circ \bigl( \lim_{a \to +0}\tilde{\gamma}_{1,\, - f_x|_{\mathbb{R}^n} (ap(x))} \circ f_x{}_{\Delta} \circ \tilde{\gamma}_{st(x),\, ap(x)}\bigr) \bigr)(u) \Bigr)  \\
            &= \Bigl( x,\, \bigl( \lim_{r \to +0} \tilde{\gamma}_{1,\, \frac{f_x|_{\mathbb{R}^n} (rp(x))}{r}} \bigr) \circ \bigl(\lim_{s \to +0} \tilde{\gamma}_{\frac{1}{s},\, 0} \circ (\tilde{\gamma}_{1,\, 0} \circ f_x{}_{\Delta} \circ \tilde{\gamma}_{t(x) s,\, 0}) \bigr)(u) \Bigr)  \\
            &= \Bigl( x,\, \bigl( \lim_{r \to +0} \tilde{\gamma}_{1,\, \frac{f_x|_{\mathbb{R}^n} (rp(x))}{r}} \bigr) \circ \bigl(\lim_{s \to +0} \tilde{\gamma}_{\frac{1}{s},\, 0} \circ f_x{}_{\Delta} \circ \tilde{\gamma}_{t(x) s,\, 0} \bigr)(u) \Bigr)  \\
            &= \Bigl( x,\, \bigl( \lim_{r \to +0} \tilde{\gamma}_{1,\, \frac{f_x|_{\mathbb{R}^n} (rp(x))}{r}} \bigr) \circ \bigl(\lim_{s' \to +0} \tilde{\gamma}_{\frac{t(x)}{s'},\, 0} \circ f_x{}_{\Delta} \circ \tilde{\gamma}_{s',\, 0} \bigr)(u) \Bigr)  \\
            &= \Bigl( x,\, \bigl( \lim_{r \to +0} \tilde{\gamma}_{1,\, \frac{f_x|_{\mathbb{R}^n} (rp(x))}{r}} \circ \gamma_{t(x),\, 0}\bigr) \circ \bigl(\lim_{s' \to +0} \tilde{\gamma}_{\frac{1}{s'},\, 0} \circ f_x{}_{\Delta} \circ \tilde{\gamma}_{s',\, 0} \bigr)(u) \Bigr)  \\
            &= \Bigl( x,\, \bigl(\lim_{r \to +0} \tilde{\gamma}_{t(x),\, \frac{f_x|_{\mathbb{R}^n} (rp(x))}{r}} \bigr) \circ \bigl(\lim_{s' \to +0} \tilde{\gamma}_{\frac{1}{s'},\, 0} \circ f_x{}_{\Delta} \circ \tilde{\gamma}_{s',\, 0} \bigr)(u) \Bigr)  \\
            &= \Bigl( x,\, \tilde{\gamma}_{t(x),\, \lim_{r \to +0} \frac{f_x|_{\mathbb{R}^n} (rp(x))}{r}} \circ D_{0}f_x(u) \Bigr)  \\
            &= \gamma_{[m]}{}_{t,\, (D_{0}{}_{[m]} f(p))_{\mathbb{R}^n}} \circ D_{0}{}_{[m]} f (x,\, u)
        \end{align}
        と計算できる\footnote{$\tilde{\gamma}_{t_2,\, p_2} \circ \tilde{\gamma}_{t_1,\, p_1} \bigl(v,\, [s,\, z]\bigr) = \tilde{\gamma}_{t_2,\, p_2}\bigl( t_1 v + p_1,\, [t_1 s,\, z] \bigr) = \bigl( t_2t_1 v + t_2 p_1 + p_2,\, [t_2t_1 s,\, z] \bigr) = \tilde{\gamma}_{t_2t_1,\, t_2 p_1 + p_2} (v,\, [s,\, z])$ である．よって $\tilde{\gamma}^{-1}_{t,\, p} = \tilde{\gamma}_{\frac{1}{t},\, -\frac{p}{t}}$ が成り立つ．}．

        　さらに，\hyperref[def:stratified-GL]{$\LGL(U)$ の定義}から $\forall T \in \LGL(U)_m$ および $\forall t > 0$ に対して
        \begin{align}
            \gamma_{[m]}{}_{\frac{1}{t},\, 0} \circ T \circ \gamma_{[m]}{}_{t,\, 0} = T
        \end{align}
        が成り立つので，\hyperref[def:c-smooth-along]{微分}の一意性より
        \begin{align}
            D_0{}_{[m]} T (x,\, u)
            &= \Bigl( x,\, \lim_{s \to +0} \tilde{\gamma}_{\frac{1}{s},\, 0} \circ T_x{}_{\Delta} \circ \tilde{\gamma}_{s,\, 0}(u) \Bigr) \\
            &=  \lim_{s \to +0} \tilde{\gamma}_{[m]}{}_{\frac{1}{s},\, 0} \circ T \circ \gamma_{[m]}{}_{s,\, 0} (x,\, u) \\
            &= T(x,\, u)
        \end{align}
        だと分かる．i.e. $D_0 \circ i_2 = \Id_{\LGL(U)}$ である．
        
        　その上，
        \begin{align}
            \Sing^{\text{sm}} ([0,\, 1]) \times \Hom{\BSC}^0 (U,\, U) &\lto \Hom{\BSC}^0 (U,\, U), \\
            (s,\, f) &\lmto \gamma_{\frac{1}{s},\, 0} \circ f \circ \gamma_{s,\, 0}
        \end{align}
        がちょうど $i_2 \circ D_0$ と $\Id_{\Hom{\BSC}^0 (U,\, U)}$ を繋ぐホモトピーになっており，$i_2$ のホモトピー逆が $D_0$ であることが示された．

        \item[\textbf{$i_3$ のホモトピー逆}]　
        
        ホモトピー逆は
        \begin{align}
            r_3{}_{[m]} \colon \Hom{\BSC}(U,\, U)_m &\lto \Hom{\BSC}^0 (U,\, U)_m, \\
            f &\lmto \gamma_{[m]}{}_{1,\, - \mathrm{proj}_{\mathbb{R}^n} \circ f|_{\Delta^m_e \times \{\bm{0}\}}} \circ f
        \end{align}
        である．
    \end{description}
\end{proof}

\subsection{$\mathcal{B}\mathrm{sc}$ の構造}

\begin{mylem}[label=lem:basic-conservative]{}
    2つの\hyperref[def:Bsc-induction]{basic} $U \coloneqq \mathbb{R}^n \times \Cone{Z},\, V \coloneqq \mathbb{R}^m \times \Cone{W} \in \Obj{\Bsc}$ であって，$\Hom{\Bsc}(U,\, V) \neq \emptyset$ を充たすものをとる．
    このとき，以下の条件は全て同値である：
    \begin{enumerate}
        \item $\depth(U) = \depth(V)$
        \item $f \in \Hom{\Bsc}(U,\, V)$ であって，$f(U) \cap \bigl(\mathbb{R}^m \times \{*\}\bigr) \neq \emptyset$ を充たすものが存在する．
        \item \hyperref[lem:Kan-enriched]{Kan複体 $\Hom{\BSC}(U,\, V)$} の\hyperref[def:fullsub-infty]{部分Kan複体}
        \begin{align}
            \mathcal{I}\mathrm{so}_{\BSC} (U,\, V) \hookrightarrow \Hom{\BSC}(U,\, V)
        \end{align}
        を次のように定義すると，その包含写像はKan複体の弱ホモトピー同値である：
        \small
        \begin{align}
            \bm{\mathcal{I}\mathrm{so}}_{\BSC} \bm{(U,\, V)}_{\bm{m}} \coloneqq \bigl\{\, f \in \Hom{\BSC}(U,\, V)_m \bigm| \forall x \in \Delta^m_e,\, \exists (\mathrm{proj}_V \circ f|_{\{x\} \times U})^{-1} \in \Hom{\Snglr}(V,\, U) \,\bigr\} 
        \end{align}
        \normalsize
        \item $(1,\, 1)$-圏 $\Bsc$ において $U$ と $V$ は同型である．
        \item $n=m$ かつ，\hyperref[def:c-smooth]{$(1,\, 1)$-圏 $\Strat$} において $Z$ と $W$ は同型である．
    \end{enumerate}
\end{mylem}

\begin{proof}
    \begin{description}
        \item[\textbf{(5)$\bm{\Longrightarrow}$(4)}]　
        
        明らか．

        \item[\textbf{(4)$\bm{\Longrightarrow}$(3)}]　
        
        命題\ref{prop:Hom-htpy-equiv}と同様の議論により従う．

        \item[\textbf{(3)$\bm{\Longrightarrow}$(2)}]　
        
        \hyperref[def:Bsc-induction]{$(1,\, 1)$-圏 $\Bsc$ における射の定義}より明らか．

        \item[\textbf{(2)$\bm{\Longrightarrow}$(1)}]　
        
        開埋め込み $\mathbb{R}^n \times \Cone{Z} \lto \mathbb{R}^m \times \Cone{W}$ の存在より $\dim(Z) = \dim(W)$ が言える．

        \item[\textbf{(1)$\bm{\Longrightarrow}$(5)}]　
        
        $f \in \Hom{\Bsc}(U,\, V)$ を1つとる．すると，\hyperref[def:Bsc-induction]{conically smoothnessの定義}よりconically smoothな層状化開埋め込み
        $\tilde{f} \in \Hom{\Snglr} (\mathbb{R}^n \times \mathbb{R}_{\ge 0} \times Z,\, \mathbb{R}^n \times \mathbb{R}_{\ge 0} \times W)$ が存在する．
        $\tilde{f}$ を制限することで，conically smoothな層状化開埋め込み $\bar{f} \in \Hom{\Snglr}\bigl(\{0\} \times \{0\} \times Z,\, \{f(0)\} \times \{0\} \times W\bigr)$ を得る．
        このとき $Z$ がコンパクトなので，$\bar{f}(Z) \subset W$ もまたコンパクトである．
        \hyperref[def:StTop]{$(1,\, 1)$-圏 $\StTOP$ の定義}より $W$ はHausdorff空間なので，$W$ のコンパクト部分空間である $\bar{f}(Z)$ は閉部分集合である．一方 $\bar{f}$ が開埋め込みであることから，$\bar{f}(Z)$ は $W$ の開部分集合でもある．
        $W$ の開かつ閉な部分集合は，$W$ の連結成分の直和しかない\footnote{位相空間の連結性の定義より，各連結成分における開かつ閉な部分集合は空集合か自分自身のちょうど2つしかない．}．然るにもし $\bar{f}(Z) \subsetneq W$ だとすると，空でない部分集合 $\Cone{\bar{f}(Z)} \subsetneq \Cone{W}$ が開かつ閉となり，$\Cone{W}$ の連結性に矛盾する．よって背理法から $\bar{f}(Z) = W$ であること，i.e. $\bar{f} \colon Z \lto W$ が $(1,\, 1)$-圏 $\Snglr$ における同型射であることが示された．
    \end{description}
    
\end{proof}


$(1,\, 1)$-圏 $\Bsc$ における対象の同型類が成す集合を
\begin{align}
    \Obj{[\BSC]} \coloneqq \bigl\{\, [U] \bigm| U \in \Obj{\Bsc} \,\bigr\}
\end{align}
と書く．この集合の元を\textbf{特異性タイプ} (singularity type) と呼ぼう．\label{def:singularity-type}
集合 $\Obj{[\BSC]}$ の上には二項関係
\begin{align}
    \label{eq:relation-Bsc}
    \le \;\coloneqq \Bigl\{\, \bigl([U],\, [V]\bigr) \in \Obj{[\BSC]}^{\times 2} \Bigm| \Hom{\Bsc} (U,\, V) \neq \emptyset \,\Bigr\} 
\end{align}
が定まる．
定理\ref{thm:basic-easy}は，特異性タイプが開埋め込みについて半順序集合を成し，さらにそれが\hyperref[def:Bsc-induction]{basic}の\hyperref[def:dim-depth]{深さ}によって特徴付けられることを主張する．
% この定理の直感は次の絵が物語っている：

\begin{mytheo}[label=thm:basic-easy]{Basics are easy}
    \begin{enumerate}
        \item $\forall U = \mathbb{R}^n \times \Cone{Z} \in \Obj{\Bsc}$ に対して，Kan複体の包含
        \begin{align}
            \LO(\mathbb{R}^n) \times \Aut (Z) \hookrightarrow \Hom{\BSC}(U,\, U)
        \end{align}
        はKan複体のホモトピー同値である\footnote{左辺を $\Sing^{\text{sm}}\bigl( \LO(n) \bigr) \times \Aut(Z)$ としても等価な主張である．}．
        \item $\forall f \in \Hom{\Bsc}(U,\, V)$ に対して，以下のいずれかちょうど1つが真である：
        \begin{enumerate}
            \item $f$ は $(\infty,\, 1)$-圏 $\BSC$ における\hyperref[def:isom-infty]{同型射}である．
            \item $\depth (U) < \depth (V)$
        \end{enumerate}
        \item 集合 $\Obj{[\BSC]}$ 上の二項関係\eqref{eq:relation-Bsc}は半順序である．半順序集合 $\bigl(\Obj{[\BSC]},\, \le \bigr)$ を $(1,\, 1)$-圏と見做し，それを $[\BSC]$ と書く．
        \item \hyperref[def:infinity-1]{$(\infty,\, 1)$-圏の関手}
        \begin{align}
            [\,\mhyphen\,] \colon \BSC \lto \Ner\bigl([\BSC]\bigr)
        \end{align}
        はconservativeである．i.e. $(U \xrightarrow{f} V) \in \BSC_1$ が $(\infty,\, 1)$-圏 $\BSC$ における\hyperref[def:isom-infty]{同型射}であるためには，$U,\, V$ が $(1,\, 1)$-圏 $\Bsc$ において同型であることが必要十分である．
        \item 写像
        \begin{align}
            \depth \colon \Obj{[\BSC]} \lto \mathbb{Z}_{\ge -1}
        \end{align}
        は順序を保つ．
    \end{enumerate}
\end{mytheo}

\begin{proof}
    \begin{enumerate}
        \item 命題\ref{prop:Hom-htpy-equiv}
        \item 補題\ref{lem:basic-conservative}の (1), (3), (4) の同値性より明らか．
        \item 反射律と推移律は自明．対称律は補題\ref{lem:basic-conservative}の (3) と (4) の同値性から従う．
        \item 補題\ref{lem:basic-conservative}より従う．
        \item 補題\ref{lem:basic-conservative}より従う．
    \end{enumerate}
\end{proof}



\subsection{tangent classifierの局所的構造}

層状化空間の局所的な特異性は\hyperref[def:dim-depth]{局所的深さ}によって特徴付けられる．これにより，局所的深さをとるという操作によって常に別の層状化を考えることができる．

集合
\begin{align}
    \mathbb{P} \coloneqq \bigl\{\, (k,\, n) \in \mathbb{Z} \times \mathbb{Z} \bigm| -1 \le k \le n \,\bigr\} 
\end{align}
の上の半順序 $\le \;\subset \mathbb{P} \times \mathbb{P}$を次のように定義する：
\begin{align}
    (k,\, n) \le (k',\, n') \DEF k \ge k' \AND n \ge n'
\end{align}

\begin{mylem}[label=lem:strat-depthdim]{局所的深さによる層状化}
    任意の\hyperref[def:Snglr-C0]{$C^0$ 級層状化空間} $(X \xrightarrow{s} P) \in \Obj{\SnglrC}$ に対して，以下の写像は連続になり，位相空間 $X$ の新しい\hyperref[def:stratified-space]{$\mathbb{P}$-層状化}を与える：
    \begin{align}
        \sigma_X \colon X &\lto \mathbb{P},\; x \lmto \bigl( \depth_x(X),\, \dim_x(X) \bigr)
    \end{align}

    \tcblower

    さらに，$(1,\, 1)$-圏 $\TOP$ における以下の図式を可換にするような，順序を保つ写像 $\textcolor{red}{\bar{s}} \colon P \lto \mathbb{P}$ が存在する：
    \begin{center}
        % https://tikzcd.yichuanshen.de/#N4Igdg9gJgpgziAXAbVABwnAlgFyxMJZABgBpiBdUkANwEMAbAVxiRAA0QBfU9TXfIRRkAjFVqMWbAArdeIDNjwEiI8uPrNWiEAB1dAWzo4AFgCMzwaV27iYUAObwioAGYAnCAaRkQOCEhqElpsCNQMdGYwDNL8ykIg7lgOJjhybp7eiL7+SABM1JpSOnAA+pw8GV6B1LmIBSAMWGDaIFB0cCb2IIWSrfpmdO7AcDbhkdGxSoJsSSlpXBRcQA
        \begin{tikzcd}
        X \arrow[d, "s"'] \arrow[r, "\sigma_X"] & \mathbb{P} \\
        P \arrow[ru, red, "\exists \bar{s}"', dashed]   &           
        \end{tikzcd}
    \end{center}
\end{mylem}

\begin{proof}
    $\sigma_X$ が連続であるような\hyperref[def:Snglr-C0]{$C^0$ 級層状化空間} $(X \to P)$ 全体が成す充満部分圏を $\Cat{T} \subset \SnglrC$ と書く．$\Cat{T}$ が条件 \hyperref[def:Snglr-C0]{\textsf{\textbf{(Snglr-1)}}-\textsf{\textbf{(Snglr-5)}}}を充たすことを示せば十分である\footnote{これらの条件を充たす最小の $(1,\, 1)$-圏のことを $\SnglrC$ と定義したため，$\Cat{T} \supset \SnglrC$ が言えるのである．}．
    \begin{description}
        \item[\textbf{(Snglr-1)}] $(\emptyset \to \emptyset) \in \Obj{\Cat{T}}$ は自明である．
        \item[\textbf{(Snglr-2)}] コンパクトな $(X \xrightarrow{s} P) \in \Obj{\Cat{T}}$ を与える．
        \hyperref[def:str-cone]{コーンの定義}から，$\forall [t,\, x] \in \Cone{X} \setminus \{\mathrm{pt}\}$ に対して $\sigma_{\Cone{X}} \bigl( [t,\, x] \bigr) = \bigl(\depth_x(X),\, \dim_x(X) + 1\bigr)$ が成り立つ．
        故に写像 $\sigma_{\Cone{X}}|_{\Cone{X} \setminus \{\mathrm{pt}\}} \colon \Cone{X} \setminus \{\mathrm{pt}\} \lto \mathbb{P}$ は連続写像の合成
        \begin{align}
            \Cone{X} \setminus \{\mathrm{pt}\} = \mathbb{R}_{> 0} \times X \xrightarrow{\mathrm{proj}_X} X \xrightarrow{\sigma_X} \mathbb{P} \xrightarrow{+(0,\, 1)} \mathbb{P}
        \end{align}
        と等しく，連続である．
        
        ところで，\exref{ex:depth-Cone}から $\mathrm{pt} \in \Cone{X}$ に対しては $\sigma_{\Cone{X}} (\mathrm{pt}) = \bigl(\dim (X) + 1,\, \dim (X) + 1\bigr)$ が成り立つ．
        故に $\mathbb{P}$ において $\sigma_{\Cone{X}} (\mathrm{pt}) \le \sigma_{\Cone{X}} \bigl( \Cone{X} \setminus \{\mathrm{pt}\} \bigr)$ が成り立つ．
        従って勝手な開集合 $Q \subset \mathbb{P}$ に対して，もし $\sigma_{\Cone{X}} (\mathrm{pt}) \notin Q$ ならば $\sigma_{\Cone{X}}^{-1} (Q) = (\sigma_{\Cone{X}}|_{\Cone{X} \setminus \{\mathrm{pt}\}})^{-1} (Q)$ が，
        もし $\sigma_{\Cone{X}} (\mathrm{pt}) \in Q$ ならば
        $\sigma_{\Cone{X}}^{-1} (Q) = \{\mathrm{pt}\} \cup \bigl( \Cone{X} \setminus \{\mathrm{pt}\} \bigr) = \Cone{X}$ 
        が成り立つので，$\sigma_{\Cone{X}}^{-1} (Q)$ が開集合であること，i.e. $\sigma_{\Cone{X}}$ が連続写像であることが示された．

        \item[\textbf{(Snglr-3)}] $\forall (X \to P) \in \Obj{\Cat{T}}$ を1つ固定する．
        このとき，\exref{ex:depth-Cone}から $\forall (x,\, t) \in X \times \mathbb{R}$ に対して $\sigma_{X \times \mathbb{R}} (x,\, t) = \bigl( \depth_x(X),\, \dim_x(X) + 1 \bigr)$ が成り立つ．
        i.e. 写像 $\sigma_{X \times \mathbb{R}} \colon X \times \mathbb{R} \lto \mathbb{P}$ は連続写像の合成
        \begin{align}
            X \times \mathbb{R} \xrightarrow{\mathrm{proj}_X} X \xrightarrow{\sigma_X} \mathbb{P} \xrightarrow{+(0,\, 1)} \mathbb{P}
        \end{align}
        に等しく，連続である．

        \item[\textbf{(Snglr-4)}] $\forall (X \xrightarrow{s} P) \in \Obj{\Cat{T}}$ および $\Hom{\StTOP} \bigl( (U \to P_U),\, (X \to P) \bigr) \neq \emptyset$ を充たす勝手な$(U \xrightarrow{u} P_U) \in \Obj{\SnglrC}$を固定する． 
        このとき\hyperref[def:strat-emb]{層状化開埋め込み} $(i,\, \bar{i}) \in \Hom{\StTOP} \bigl( (U \to P_U),\, (X \to P) \bigr) \neq \emptyset$ をとると，開埋め込みが\hyperref[def:dim-depth]{局所的次元}を保存すること\footnote{\hyperref[def:Snglr-C0]{$\SnglrC$ の定義から} $\SnglrC$ の対象はパラコンパクトなので．}から
        $\sigma_{U} \colon U \lto \mathbb{P}$ は連続写像の合成
        \begin{align}
            U \overset{i}{\hookrightarrow} X \xrightarrow{s} P \xrightarrow{\bar{s}} \mathbb{P}
        \end{align}
        に等しく，連続である．
        \item[\textbf{(Snglr-5)}] $\Cat{T}$ 内に開被覆 $\Familyset[\big]{(i_\lambda,\, \bar{i}_\lambda) \in \Hom{\StTOP} \bigl( (U_\lambda \xrightarrow{\sigma_\lambda} P_\lambda),\, (X \xrightarrow{s} P) \bigr) }{\lambda \in \Lambda}$ を持つような\hyperref[def:Snglr-C0]{$C^0$ 級層状化空間} $(X \xrightarrow{s} P) \in \Obj{\SnglrC}$ を与える．
        このとき，勝手な開集合 $Q \subset \mathbb{P}$ に対して
        \begin{align}
            \sigma_X^{-1} (Q)
            &= \bigcup_{\lambda \in \Lambda} i_\lambda \Bigl(U_\lambda \cap i_\lambda^{-1} \bigl( \sigma_X^{-1}(Q) \bigr) \Bigr)
        \end{align}
        が成り立つが，$i_\lambda \Bigl(U_\lambda \cap i_\lambda^{-1} \bigl( \sigma_X^{-1}(Q) \bigr) \Bigr) \subset X$ は開なので $\sigma_X^{-1} (Q)$ は $X$ の開集合である．i.e. $\sigma_X$ は連続写像である．
    \end{description}
    % \begin{description}
    %     \item[\textbf{$\bm{\sigma_X}$ の連続性}]　
        

    %     \item[\textbf{$\bm{\bar{s}}$ の一意存在}]　
        
    %     % 補題\ref{lem:strat-surj}より $s \colon X \lto P$ は全射なので明らか．
    % \end{description}
\end{proof}

補題\ref{lem:basic-conservative}および定理\ref{thm:basic-easy}より，補題\ref{lem:strat-depthdim}で定義した局所的深さによる層状化 $s_X \colon X \lto \mathbb{P}$ は，$X$ の各点における\hyperref[def:singularity-type]{特異性タイプ}を特徴づける．
しかるに，同じ深さだが異なる\hyperref[def:singularity-type]{特異性タイプ}を検出することができない．
従って，そもそも値域を $\OP{[\BSC]}$ そのものに持つような層状化 $\alpha_X \colon X \lto \OP{[\BSC]}$ があれば実用上便利である．
補題\ref{lem:strat-basic}においてそのような $\alpha_X$ を構成しよう．
直感的には，補題\ref{lem:strat-basic}の構成とは次のようなものである．

いま $X$ として，赤道上に\hyperref[def:dim-depth]{局所的深さ}がそれぞれ $1,\, 2$ であるような\hyperref[def:stratified-space]{層}を持つ2次元球面を考え，
3点 $x,\, y,\, z \in X$ を下の図のようにとる．
\begin{align}
    \label{ex:S2withDefect}
    X = \quad
    \begin{tikzpicture}[baseline={([yshift=-.5ex]current bounding box.center)}]
        \path coordinate (O)
        ;
        \shade[ball color = lightgray,opacity = .4] (O) circle (2);
        % 1-strata
        \halfDottedEllipse{(180:2)}{2}{0.5}
        % 0-strata
        \node[bullet] (pt) at ($(O)+(0,-0.5)$) {};
        % points
        \draw (60:1.5) node[cross=3pt,red] (x) {};
        \draw (-160:1.2) node[cross=3pt,blue] (y) {};
        \draw (pt.center) node[cross=3pt,green] (z) {};
        \node[above,red] at (x) {$x$};
        \node[above,blue] at (y) {$y$};
        \node[above right,DarkGreen] at (z) {$z$};
        % embedding
        \draw[rotate around={-30:(x)}, dashed] (x) ellipse (0.5 and 0.3);
        \draw[rotate around={90:(y)}, dashed] (y) ellipse (0.5 and 0.4);
        \draw[dashed] (z) circle (0.5);
    \end{tikzpicture}
\end{align}
このとき，$\alpha_X$ が各点に割り当てる特異性タイプが
\begin{align}
    \alpha_X (\textcolor{red}{x})
    &= \left[ 
        \begin{tikzpicture}[baseline={([yshift=-.5ex]current bounding box.center)}]
            \fill[gray!20,opacity=.4] (0,0) circle (1);
            \node[cross=3pt,red] (x) at (0,0) {};
            \node[above,red] at (x) {$0$};
        \end{tikzpicture}
     \right] = [\mathbb{R}^2], \\
     \alpha_X (\textcolor{blue}{y})
     &= \left[ 
        \begin{tikzpicture}[baseline={([yshift=-.5ex]current bounding box.center)}]
            \fill[gray!20,opacity=.4] (0,0) circle (1);
            \draw[thick] (180:1) -- (0:1);
            \node[cross=3pt,blue] (y) at (0,0) {};
            \node[above,blue] at (y) {$(0,\, \mathrm{pt})$};
        \end{tikzpicture}
      \right] = [\mathbb{R} \times \Cone{\bullet \amalg \bullet}], \\
      \alpha_X (\textcolor{DarkGreen}{z})
      &= \left[ 
        \begin{tikzpicture}[baseline={([yshift=-.5ex]current bounding box.center)}]
            \fill[gray!20,opacity=.4] (0,0) circle (1);
            \draw[thick] (180:1) -- (0:1);
            \node[bullet] at (0,0) {};
            \node[cross=3pt,green] (z) at (0,0) {};
            \node[above right,DarkGreen] at (z) {$\mathrm{pt}$};
        \end{tikzpicture}
       \right] = \left[ 
            \Cone{
                \begin{tikzpicture}[baseline={([yshift=-.5ex]current bounding box.center)}]
                    \draw[gray] (0,0) circle (.5);
                    % spoke
                    \foreach \q in {0,180} {
                        \node[vertex] at (\q:.5) {};
                    }
                \end{tikzpicture}
            }
        \right] 
\end{align}
となるように $\alpha_X$ を構成するのである．

\begin{mylem}[label=lem:strat-basic,breakable]{半順序集合 {$\OP{[\BSC]}$} による層状化}
    任意の\hyperref[def:c-smooth]{conically smoothな層状化空間} $(X \xrightarrow{s} P) \in \Obj{\Snglr}$ に対して，以下はwell-definedな連続写像になり，位相空間 $X$ の新しい\hyperref[def:stratified-space]{$\OP{[\BSC]}$-層状化}を与える：
    \begin{align}
        \alpha_X \colon X \lto \OP{[\BSC]},\; x \lmto \Bigl[\, \argmax_{U \in \mathcal{U}_x} \bigl(\, \depth \colon \mathcal{U}_x \lto \mathbb{Z}_{\ge -1} \bigr) \,\Bigr]
    \end{align}
    ただし，$\forall x \in X$ に対して集合 $\mathcal{U}_x \subset \Obj{\Bsc}$ を以下で定義した．
    \begin{align}
        \mathcal{U}_x \coloneqq \bigl\{\, U \in \Obj{\Bsc} \bigm| \exists \iota \in \Hom{\Snglr}(U,\, X),\; x = \iota(\bm{0}) \,\bigr\} 
    \end{align}
    さらに，$\alpha_X \colon X \lto \OP{[\BSC]}$ は補題\ref{lem:strat-depthdim}による層状化 $\sigma_X \colon X \lto \mathbb{P}$ を次のように分解する：
    \begin{center}
        % https://tikzcd.yichuanshen.de/#N4Igdg9gJgpgziAXAbVABwnAlgFyxMJZABgBpiBdUkANwEMAbAVxiRAA0QBfU9TXfIRQBGclVqMWbADrSAtnRwALAEYrgABS7deIDNjwEiZYePrNWiELIDyG4MlkAhAMoBhCtq7iYUAObwRKAAZgBOEHJIZCA4EEiiEhYy0th+CgD6nDwh4ZGI0bFIAEzU5lJWsoxoSnSZINQMdCowDBr8hkIgoVh+Sjg6ORHF1IWICWWWIAAUsrBoyqSypAAEs1hyAJT1II3Nre2CbN29-d5cQA
        \begin{tikzcd}
            X \arrow[r, "\sigma_X"] \arrow[d, "\alpha_X"']  & \mathbb{P} \\
            {\OP{[\BSC]}} \arrow[ru, "{(\depth,\, \dim)}"'] &           
        \end{tikzcd}
    \end{center}

    \tcblower

    任意の射 $f \in \Hom{\Snglr} (X,\, Y)$ に対して，$\alpha_X = \alpha_Y \circ f$ が成り立つ．
\end{mylem}

\begin{proof}
    余計な煩雑さを回避するために，以下の証明において\hyperref[def:Sm-induction]{conically smooth atlas}は常に省略する．
    
    $\forall x \in X$ を1つ固定する．$\forall U \in \Cat{U}_x$ に対して，\hyperref[def:dim-depth]{深さの定義}より $\depth_x (U) = \depth (U) \le \depth (X \xrightarrow{s} P) < \infty$ が成り立つので，$U_x \coloneqq \argmax_{U \in \mathcal{U}_x} (\depth) \in \mathcal{U}_x$ は存在する．
    補題\ref{lem:basic-conservative}より $\depth_x(U_x) = \depth_x(U'_x)$ ならば $[U_x] = [U'_x]$ であり，写像 $\alpha_X$ がwell-definedであることが示された．

    次に $\alpha_X$ の連続性を $\depth(X \to P)$ に関する数学的帰納法により示す．$\depth(X \to P) = -1$ のとき，$X = \emptyset$ であり自明である．
    $\forall (X \to P) \in \Obj{\Snglr_{\le d,\, \le \infty}}$ についての証明が完了しているとする．

    \begin{description}
        \item[\textbf{basicに関する連続性}]　

        まず，$\forall (U \xrightarrow{s} P) = (\mathbb{R}^n \xrightarrow{0} [0]) \times \Cone{Z \xrightarrow{r} Q} \in \Obj{\Bsc_{d+1,\, \le \infty}}$ について写像 $\alpha_U \colon U \lto [\BSC]$ が連続であることを示そう．
        
        　まず $\forall x \in \mathbb{R}^n \times \{\mathrm{pt}\} \subset \mathbb{R}^n \times \Cone{Z} = U$ に対しては
        $U \in \mathcal{U}_x$ であり，さらに $U_{s(x)} = s^{-1}\bigl(\{(0,\, -\infty)\}\bigr) = \mathbb{R}^n \times \{\mathrm{pt}\}$ なので
        \begin{align}
            \depth_x(U \xrightarrow{s} P) 
            &= \dim_x \bigl(\mathbb{R}^n \times \Cone{Z} \bigr)  - \dim_x (U_{s(x)}) \\
            &= n + \dim_{\mathrm{pt}} \bigl( \Cone{Z} \bigr) - n \\
            &= \depth_{\mathrm{pt}} \bigl( \Cone{Z \to Q} \bigr)\\
            &= \depth(U)
        \end{align}
        が成り立ち，$\alpha_U(x) = [U]$ と計算できる．
        一方で，$\forall x = \bigl(u,\, (t,\, z)\bigr) \in U \setminus (\mathbb{R}^n \times \{\mathrm{pt}\}) = \mathbb{R}^n \times (\mathbb{R}_{> 0} \times Z)$ に対して $\alpha_U(x)$ を計算しよう．
        $\forall (V \xrightarrow{v} P_V) \in \mathcal{U}_x$ に対して，層状化開埋め込み $(\iota,\, \bar{\iota}) \in \Hom{\Snglr}(V,\, U)$ とは
        % https://tikzcd.yichuanshen.de/#N4Igdg9gJgpgziAXAbVABwnAlgFyxMJZABgBpiBdUkANwEMAbAVxiRADUQBfU9TXfIRQBGclVqMWbAKrdeIDNjwEiZYePrNWiEAAUA+px58lgoqPXVNUnbu7iYUAObwioAGYAnCAFskZEBwIJFFAuiwGNgALCAgAaxArSW0QAB1U-Bw6OQ9vP0RQoKQAZmoGOgAjGAZdfmUhEAYYdxxEiS02BGMQL19-aiLEACYkjp0aNvKqmrqzHU8sJyjW7t78kcDgxFL2mzTUirpPYHTMui5Jyura0xV5xeX7LiA
        $(1,\, 1)$-圏 $\TOP$ の可換図式
        \begin{center}
            \begin{tikzcd}
                V \arrow[r, "\iota", hook] \arrow[d, "v"'] & U \arrow[d, "s"] \\
                P_V \arrow[r, "\bar{\iota}"']              & P               
            \end{tikzcd}
        \end{center}
        であって，$\forall p \in P_V$ に対して
        \begin{align}
            \iota|_{V_p} \colon V_p \lto U_{\bar{\iota}(p)}
        \end{align}
        が位相的開埋め込みになっているものであったことを思い出そう．したがって
        \begin{align}
            \iota|_{V_{v(x)}} \colon V_{v(x)} \lto U_{\bar{\iota} \circ v(x)} = U_{s(x)}
        \end{align}
        は開埋め込みであり
        \footnote{正確には，$V_{v(x)}$ と書いたときは $V_{v \circ \iota^{-1}(x)}$ を意味している．図式の可換性から $U_{\bar{\iota} \circ v \circ \iota^{-1}(x)} = U_{s(x)}$ が分かる．}，
        \begin{align}
            \dim_x(V_{v(x)}) = \dim_x(U_{s(x)})
        \end{align}
        が成り立つ．いま $s(x) = \bigl(0,\, r(z)\bigr)$ であるから
        $U_{s(x)} = s^{-1} \bigl( \{s(x)\} \bigr) = \mathbb{R}^n \times \mathbb{R}_{> 0} \times Z_{r(z)}$ であり
        \begin{align}
            \dim_x(U_{s(x)}) = n + 1 + \dim_z (Z_{r(z)})
        \end{align}
        と計算できる．そのうえ定理\ref{thm:basic-easy}-(2) より $\dim_x(V) \le \dim_x(U) = n + \dim_{[t,\, z]} \bigl( \Cone{Z} \bigr)$ が言えるので，
        \begin{align}
            \depth_x (V \xrightarrow{v} P_V) 
            &= \dim_x (V) - \dim_x (V_{v(x)}) \\
            &\le n + \dim_{[t,\, z]} \bigl( \Cone{Z} \bigr) - \bigl( n + 1 + \dim_z (Z_{r(z)}) \bigr) \\
            &< \dim_{[t,\, z]} \bigl( \Cone{Z} \bigr) - \dim_z (Z_{r(z)}) \\
            &= \depth_x(U)  \\
            &\le \depth(U)
        \end{align}
        が分かった．故に $\alpha_U(x) = [V_0]$ なる $V_0 \in \Obj{\Bsc}$ は常に $\depth(V_0) < \depth(U)$ を充たすことが分かった．
        
        　ここで，位相空間 $\OP{[\BSC]}$ の勝手な開集合 $\forall \mathcal{W} \subset \OP{[\BSC]}$ をとる．
        もし $[U] \in \mathcal{W}$ ならば，\hyperref[def:topo-poset]{半順序集合 $\OP{[\BSC]}$ の位相の定義}から $\forall [V] \in \OP{[\BSC]}$ に対して 
        \begin{align}
            [U] \le [V] \IFF \Hom{\Bsc}\bigl([V],\, [U]\bigr) \neq \emptyset \IMP [V] \in \mathcal{W}
        \end{align}
        が成り立つ\footnote{半順序\eqref{eq:relation-Bsc}の $\mathrm{op}$ をとっていることに注意！}．
        故に補題\ref{lem:basic-conservative}および定理\ref{thm:basic-easy}-(2) より，$\mathcal{W}$ は $\depth(V) \le \depth(U)$ を充たす全ての $[V] \in \OP{[\BSC]}$ を含む． 
        従って $\alpha_U^{-1}(\mathcal{W}) = U$ であり，$\alpha_U^{-1}(\mathcal{W})$ が $U$ の開集合であることが分かった．
        一方でもし $[U] \notin \mathcal{W}$ であるならば，\hyperref[def:topo-poset]{半順序集合 $\OP{[\BSC]}$ の位相の定義}からある $[V] \in \mathcal{W}$ が存在して
        $[U] \lneq [V] \IFF \Hom{\Bsc}\bigl([V],\, [U]\bigr) \neq \emptyset$ を充たす．そのような $[V]$ のうち最小のものを $[V_0]$ と書くと $\mathcal{W} = \bigl\{\, [V] \in \OP{[\BSC]} \bigm| [V_0] \le [V] \,\bigr\}$ であり，かつ定理\ref{thm:basic-easy}-(2) より $\depth (U) > \depth (V_0)$ でなくてはいけない．
        故に帰納法の仮定が使えて，$\alpha_U^{-1}(\mathcal{W}) \subset \mathbb{R}^n \times \mathbb{R}_{> 0} \times Z \subset U$ が $U$ の開集合であることが分かった．

        \item[\textbf{一般の $\bm{X \to P}$ に関する連続性}]　
        
        $(X \to P)$ の，\hyperref[def:Bsc-induction]{basic}による開被覆 $\Familyset[\big]{i_\lambda \in \Hom{\Snglr} (U_\lambda,\, X)}{\lambda \in \Lambda}$ をとる\footnote{\hyperref[def:Sm-induction]{conically smooth atlas}が少なくとも1つ存在するので，開被覆は存在する．}．
        $\forall \lambda \in \Lambda$ に対して，$i_\lambda$ が層状化開埋め込みなので $\alpha_{U_\lambda} = \alpha_X \circ i_\lambda \colon U_\lambda \lto \OP{[\BSC]}$ が成り立つことに注意する．

        　ここで位相空間 $\OP{[\BSC]}$ の勝手な開集合 $\forall \mathcal{W} \subset \OP{[\BSC]}$ をとる．すると
        \begin{align}
            \alpha_X^{-1} (\mathcal{W}) \cap i_\lambda (U_\lambda) 
            &= i_\lambda \bigl( (\alpha_X \circ i_\lambda)^{-1} (\mathcal{W}) \bigr) \\
            &= i_\lambda \bigl( (\alpha_{U_\lambda})^{-1} (\mathcal{W}) \bigr)
        \end{align}
        と書ける．ところで，basicの場合の証明より $\alpha_{U_\lambda} \colon U_\lambda \lto \OP{[\BSC]}$ は連続であるから，$(\alpha_{U_\lambda})^{-1} (\mathcal{W})$ は $U_\lambda$ の開集合である．
        そのうえ $i_\lambda$ が開埋め込みであることから $\alpha_X^{-1} (\mathcal{W}) \cap i_\lambda (U_\lambda)$ は $X$ の開集合であり，
        \begin{align}
            \alpha_X^{-1} (\mathcal{W}) = \bigcup_{\lambda \in \Lambda} \alpha_X^{-1} (\mathcal{W}) \cap i_\lambda (U_\lambda) 
        \end{align}
        が $X$ の開集合であることが示された．
        % 上述の考察より $\alpha_U^{-1}([U]) = \mathbb{R}^n \times \{\mathrm{pt}\}$ である．
    \end{description}
    
    $\alpha_X$ の定義から，$\forall x \in X$ に対して
    \begin{align}
        \sigma_X(x) 
        &= \bigl(\depth_x(X),\, \dim_x(X)\bigr) \\
        &= \bigl(\depth(\alpha_X(x)),\, \depth(\alpha_X(x)) \bigr)
    \end{align}
    が成り立つことは明らかである．
\end{proof}

\begin{mydef}[label=def:EntrU]{制限されたtangent classifier}
    $\forall (X \to P) \in \Obj{\Snglr}$ および\hyperref[def:Bsc-induction]{basic} $\forall U = \mathbb{R}^n \times \Cone{Z} \in \Obj{\Bsc}$ を与える．
    このとき，$(\infty,\, 1)$-圏 $\textcolor{red}{\BSC_{[U]}},\, \textcolor{red}{\ENTRU{[U]}{X}}$ および $(\infty,\, 1)$-圏の関手 $\textcolor{red}{(\tau_X)_{|}} \in \Hom{\sSet}\bigl(\textcolor{red}{\ENTRU{[U]}{X}},\, \textcolor{red}{\BSC_{[U]}}\bigr)$ を
    以下の\hyperref[def:pullback-infty]{$(\infty,\, 1)$-圏の引き戻し}によって次のように定義する
    ：
    \begin{center}
        % https://tikzcd.yichuanshen.de/#N4Igdg9gJgpgziAXAbVABwnAlgFyxMJZABgBpiBdUkANwEMAbAVxiRAB12BRAOQBUASsGQBVCgF9gADXEhxpdJlz5CKMgEYqtRizacAQgGUAwgH1hY2fMXY8BImQBMW+s1aIO7AEZYA5g04LCk4fXwAnTisFEAxbFSJ1UmdqV10PZAMTCjlo2OV7FETNFJ13TyNjHJt81WREyhK3PW5+IRk5LRgoX3giUAAzMIgAWyQAVmocCCQAFka0zxw6JlMpKpBBkdnJ6cQAZnmyjPZhgAsATzRTmDBs6w2h0cQyECnx+82nxNfdmY-HpCOHZIA6vOhYBhsU4QCAAa3WnyQ3zeiEc-y2z2BiG+qTKAApOEsVlIAJTmAA+smoDDoXhgDAACko7KoQGE-KccB1xEA
        \begin{tikzcd}
        {\textcolor{red}{\ENTRU{[U]}{X}}} \arrow[dr, phantom," " {pullback}, very near start, color=black] \arrow[r] \arrow[d, red, "(\tau_X)_{|}"'] & \ENTR{X} \arrow[d, "\tau_X"]   \\
        {\textcolor{red}{\BSC_{[U]}}} \arrow[dr, phantom," " {pullback}, very near start, color=black] \arrow[r] \arrow[d]                     &\BSC \arrow[d, "{[\, \mhyphen \,]}"] \\
        {\bigl\{[U]\bigr\}} \arrow[r, hook]                  & {\Ner\bigl([\BSC]\bigr)}                      
        \end{tikzcd}
    \end{center}
\end{mydef}

まず，\hyperref[def:infinity-1]{$(\infty,\, 1)$-圏} $\BSC_{[U]}$ の構造を詳細に調べよう．
定理\ref{thm:basic-easy}-(4) より $\BSC \xrightarrow{[\, \mhyphen \,]} \Ner \bigl( [\BSC] \bigr)$ は\hyperref[def:isofibration]{アイソファイブレーション}なので，補題\ref{lem:pullback-infty}, \ref{lem:pullback-iso}から $(\infty,\, 1)$-圏の引き戻し
\footnote{混乱防止のため，\hyperref[def:Strat-infty]{$\Kan$-豊穣圏 $\BSC$} およびその\hyperref[def:nerve-hc]{homotopy coherent nerve} $\hcNer (\BSC)$ を区別した．}
\begin{center}
    % https://tikzcd.yichuanshen.de/#N4Igdg9gJgpgziAXAbVABwnAlgFyxMJZABgBoBGAXVJADcBDAGwFcYkQAdDgERkZ3oA9YiAC+pdJlz5CKchWp0mrdlwByMAE5cARlgDmjABTIuAIQDKAYUq6DmgJRiJIDNjwEi84ooYs2iJwcllbOku4yRGQ+NH4qgebWAPrAyACqlKJiijBQ+vBEoABmmhAAtkhkIDgQSPLV9FiM7AAWEBAA1iCxygEg6ZTdIIz0OnwAClIesiCaBi04YSAl5UgAzDQ1SABM4sWlFYjbm7WI9XF9phykAARcZS0AnmgtMGB314N7ywfrJ5WiTJAA
    \begin{tikzcd}
    {\BSC_{[U]}} \arrow[r] \arrow[d] \arrow[dr, phantom," " {pullback}, very near start, color=black]  & \hcNer(\BSC) \arrow[d, "{[\, \mhyphen \,]}"] \\
    \Delta^0 \arrow[r, "{[U]}"', hook] & {\Ner\bigl([\BSC]\bigr)}            
    \end{tikzcd}
\end{center}
は，$(1,\, 1)$-圏 $\sSet$ の引き戻しとして計算できる．よって
\begin{align}
    (\BSC_{[U]})_0 = \bigl\{\, V \in \Obj{\Bsc} \bigm| V \in [U] \,\bigr\} 
\end{align}
である．さらに定理\ref{thm:Map-hcNer}より，$\BSC_{[U]}$ の\hyperref[def:Map]{射の空間}について\hyperref[def:htpy-equiv-Kan]{Kan複体のホモトピー同値}
\begin{align}
    \label{eq:Map-BscU}
    \inftyMap{\BSC_{[U]}} (V,\, W) \simeq \Hom{\BSC} (V,\, W) \simeq \inftyMap{\hcNer(\BSC)} (V,\, W)
\end{align}
が成り立つ．故に補題\ref{lem:basic-conservative}および定理\ref{thm:basic-easy}から，\hyperref[def:Kan-group]{群的なKan複体}の\hyperref[def:htpy-equiv-Kan]{ホモトピー同値}
\begin{align}
    \label{eq:htpy-equiv-BAut2BscU}
    \deloop{\Aut(U)} \overset{\simeq}{\hookrightarrow} \BSC_{[U]}
\end{align}
が存在することがわかる．ただし，群的なKan複体 \hyperref[def:Kan-Aut]{$\Aut(U)$} の\textbf{delooping}と呼ばれるKan複体 $\deloop{\Aut(U)}$ を次のように定義した
\footnote{より正確には，対象が1点で\hyperref[def:enriched]{Hom対象}が $\Aut(U) \in \Obj{\Kan}$ であるような $\Kan$-豊穣圏に\hyperref[def:nerve-hc]{homotopy coherent nerve}を適用して得られる $(\infty,\, 1)$-圏（特に，この場合は群的なKan複体）のことである．この構成は，\hyperref[def:delooping-groupoid]{群のdelooping groupoid}の構成を素朴に一般化したものだと見做せる．}：
\begin{align}
    \deloop{\Aut(U)} \coloneqq \hcNer \biggl( 
        % https://tikzcd.yichuanshen.de/#N4Igdg9gJgpgziAXAbVABwnAlgFyxMJZABgBpiBdUkANwEMAbAVxiRACoQBfKkGKAObwioAGYAnCAFskZEDgizqDCBDREAjAA4yoxnBi8GdAEYwGABUy58hRCHFYBACxwhq9Zq3sAdHwEEmHAAKAFUASm4KLiA
        \begin{tikzcd}[baseline={([yshift=-.5ex]current bounding box.center)}]
        * \arrow["\Aut(U)"', loop, distance=2em, in=125, out=55]
        \end{tikzcd}
     \biggr) 
\end{align}

次に，$(\infty,\, 1)$-圏 $\ENTRU{[U]}{X}$ の構造を詳細に調べる．
~\cite[\href{https://kerodon.net/tag/018F}{Tag 018F}]{kerodon}より $\tau_X$ が\hyperref[def:infty-fib]{右ファイブレーション}（したがって図\ref{fig:infty-fib}より\hyperref[def:isofibration]{アイソファイブレーション}）なので，補題\ref{lem:pullback-infty}, \ref{lem:pullback-iso}から $(\infty,\, 1)$-圏の引き戻し
\begin{center}
    % https://tikzcd.yichuanshen.de/#N4Igdg9gJgpgziAXAbVABwnAlgFyxMJZABgBpiBdUkANwEMAbAVxiRAB12BRAOQBUASsGQBVCgF9gADXEhxpdJlz5CKMgEYqtRizacAQgGUAwgH1hY2fMXY8BImQBMW+s1aIO7AEZYA5g04LCk4fXwAnTisFEAxbFSJ1UmdqV10PZAMTCjlo2OV7FETNFJ13TyNjHJt81WREyhK3PW5+IRk5LRgoX3giUAAzMIgAWyQAVmocCCQAFka0zxw6JlMpKpBBkdnJ6cQAZnmyjPZhgAsATzRTmDBs6w2h0cQyECnx+82nxNfdmY-HpCOHZIA6vOhYBhsU4QCAAa3WnyQ3zeiEc-y2z2BiG+qTKAApOEsVlIAJTmAA+smoDDoXhgDAACko7KoQGE-KccB1xEA
    \begin{tikzcd}
    {\ENTRU{[U]}{X}} \arrow[dr, phantom," " {pullback}, very near start, color=black] \arrow[r] \arrow[d, "(\tau_X)_{|}"'] & \ENTR{X} \arrow[d, "\tau_X"]   \\
    {\BSC_{[U]}} \arrow[r]                    & \hcNer(\BSC)
    \end{tikzcd}
\end{center}
は $(1,\, 1)$-圏 $\sSet$ の引き戻しとして計算できる．
故に\exref{ex:join-2}および定義\ref{def:nerve-hc}から，$\ENTRU{[U]}{X}$ の $n$-単体は
\begin{align}
    \ENTRU{[U]}{X}_n 
    &\cong (\BSC_{[U]}{}_{/X})_n \\
    &\cong \Bigl\{\, H \in \Hom{\sSet} \bigl(\Delta^{n+1},\, \hcNer(\SNGLR)\bigr) \Bigm| \substack{H_{[0]}\bigl(\{n+1\}\bigr) = X,\\ H_{[0]} \bigl( \{i\} \bigr) \in [U]\; (0\le \forall i \le n)}  \,\Bigr\} \\
    &= \Bigl\{\, H \in \Hom{\sCat} \bigl(\mathfrak{C}[n+1],\, \SNGLR\bigr) \Bigm| \substack{H_{0}(n+1) = X,\\ H_{0} (i) \in [U]\; (0\le \forall i \le n)}  \,\Bigr\} \label{eq:EntrU-simplex}
\end{align}
であり，$X$ を除く頂点に\hyperref[def:singularity-type]{特異性タイプ} $[U]$ の\hyperref[def:Bsc-induction]{basic}が載る形になる．
例えば $0$-単体は
\begin{align}
    \ENTRU{[U]}{X}_0 = \bigcup_{V \in [U]} \Hom{\Snglr} (V,\, X)
\end{align}
である．
$1$-単体は\exref{ex:hcNer-2}より
\begin{align}
    \ENTRU{[U]}{X}_1 \cong \left\{\, 
        \begin{tikzpicture}[baseline={([yshift=-.5ex]current bounding box.center)}]
            \path coordinate[bullet,label=below:$X$] (X)
            +(120:2) coordinate[bullet,label=above left:$V$] (V)
            +(60:2) coordinate[bullet,label=above right:$W$] (W)
            ;
            \begin{scope}[on background layer]
                \fill[gray!20] (V.center) -- (W.center) -- (X.center) --cycle;
            \end{scope}
            % \node[gray] at (barycentric cs:X=1,V=1,W=1) {$H$};
            \draw[->,gray,double] ($(W) + (-150:.3)$) -- node[midway,inner sep=1pt] {$H$} ($(W) + (-150:1.5)$);
            \draw[->-=.5] (V.center) -- node[midway,left] {$\iota_V$} (X.center);
            \draw[->-=.5] (W.center) -- node[midway,right] {$\iota_W$} (X.center);
            \draw[->-=.5] (V.center) -- node[midway,above] {$f$} (W.center);
        \end{tikzpicture}
     \middle| \substack{
        V,\, W \in [U],\\ 
        \iota_V \in \Hom{\Snglr} (V,\, X), \\ 
        \iota_W \in \Hom{\Snglr} (W,\, X), \\ 
        f \in \Hom{\BSC} (V,\, W)_0, \\   
        H \in \Hom{\SNGLR}(V,\, X)_1, \\ 
        H|_{\{(1,\, 0)\} \times V} = \iota_W \circ f,\; H|_{\{(0,\, 1)\} \times V} = \iota_V
        } 
    \,\right\} 
\end{align}
であることがわかる．i.e. $\ENTRU{[U]}{X}$ の1-射とは，以下のデータの組である：
\begin{itemize}
    \item 特異性タイプが $[U]$ の\hyperref[def:Bsc-induction]{basic} $V,\, W \in \Obj{\Bsc}$
    \item 3つの\hyperref[def:strat-emb]{層状化開埋め込み} $\iota_V \in \Hom{\Snglr}(V,\, X),\; \iota_W \in \Hom{\Snglr} (W,\, X),\; f \in \Hom{\Snglr} (V,\, W)$
    \item $\iota_W \circ f$ と $\iota_V$ を繋ぐアイソトピー\footnote{より正確には，$\mathrm{proj}_X \circ H \in \Hom{\Snglr} (\Delta^1_e \times V,\, X)$ がアイソトピーである．} $H \in \Hom{\Snglr} (\Delta^1_e \times V,\, \Delta^1_e \times X)$
\end{itemize}

以上の考察から，$(\infty,\, 1)$-圏 $\ENTRU{[U]}{X}$ とは直観的には，
\begin{itemize}
    \item $(X \to P)$ 上に与えられた\hyperref[def:Sm-induction]{conically smooth アトラス} $\Sm(X \to P)$ に属するチャート $(V,\, V \overset{\iota_V}{\hookrightarrow} X) \in \Sm(X \to P)$ のうち，特異性タイプが $[U]$ であるものを対象に持つ．
    \item \hyperref[def:singularity-type]{特異性タイプ}が $[U]$ であるチャート $V \overset{\iota_V}{\hookrightarrow} X$ の間のアイソトピーを1-射に持つ．
    \item アイソトピーの間のアイソトピーを2-射に持つ．
    \item $\cdots$
\end{itemize}
ような $(\infty,\, 1)$-圏であることが分かった．

\begin{mydef}[label=def:XU]{{$[U]$}-層}
    $\forall (X \to P) \in \Obj{\Snglr}$ および\hyperref[def:Bsc-induction]{basic} $\forall U = \mathbb{R}^n \times \Cone{Z} \in \Obj{\Bsc}$ を与える．
    このとき位相空間 $\bm{X_{[U]}} \in \Obj{\TOP}$ を，$(1,\, 1)$-圏 $\TOP$ における\hyperref[def:pullback-pushout]{引き戻し}によって次のように定義する：
    \begin{center}
        % https://tikzcd.yichuanshen.de/#N4Igdg9gJgpgziAXAbVABwnAlgFyxMJZARgBoAGAXVJADcBDAGwFcYkQANEAX1PU1z5CKMsWp0mrdsgA6MgEIBlAMKUefEBmx4CRcqTE0GLNohByARlgDmjOcGQBVSpZsAnOd3X9tQvRXFjKTMOAH0HZy9ucRgoa3giUAAzNwgAWyQAZhocCCQAJl5k1IzEfRBcpDIJE3YLUK4ikBT0rJy8sqaW0vz2qpz6LEZ2AAsICABrHkpuIA
        \begin{tikzcd}
        {X_{[U]}} \arrow[dr, phantom," " {pullback}, very near start, color=black] \arrow[d] \arrow[r]       & X \arrow[d, "\alpha_X"] \\
        {\bigl\{[U]\bigr\}} \arrow[r, hook] & {\OP{[\BSC]}}          
        \end{tikzcd}
    \end{center}
\end{mydef}

つまり，$X_{[U]}$ とは補題\ref{lem:strat-basic}で構成した $\OP{[\BSC]}$-層状化 $(X \xrightarrow{\alpha_X} \OP{[\BSC]})$ の，特異性タイプ $[U] \in \OP{[\BSC]}$ 上のファイバーである．
例えば
\begin{align}
    X = \quad
    \begin{tikzpicture}[baseline={([yshift=-.5ex]current bounding box.center)}]
        \path coordinate (O)
        ;
        \shade[ball color = lightgray,opacity = .4] (O) circle (2);
        % 1-strata
        \halfDottedEllipse{(180:2)}{2}{0.5}
        % 0-strata
        \node[bullet] (pt) at ($(O)+(0,-0.5)$) {};
    \end{tikzpicture}
\end{align}
の場合，空でない $X_{[U]}$ は
\begin{align}
    \left( 
        \begin{tikzpicture}[baseline={([yshift=-.5ex]current bounding box.center)}]
            \path coordinate (O)
            ;
            \shade[ball color = lightgray,opacity = .4] (O) circle (2);
            % 1-strata
            \halfDottedEllipse{(180:2)}{2}{0.5}
            % 0-strata
            \node[bullet] (pt) at ($(O)+(0,-0.5)$) {};
        \end{tikzpicture}
     \right)_{[\mathbb{R}^2]} 
    &= \quad
    \begin{tikzpicture}[baseline={([yshift=-.5ex]current bounding box.center)}]
        \path coordinate (N)
        +(0,-1.5) coordinate (S)
        ;
        % northern hemisphere
        \begin{scope}
            \clip ($(N)+(0:2)$) arc(0:180:2) ($(N)+(180:2)$) arc(-180:0:2 and .5) --cycle;
            \shade[ball color = lightgray,opacity = .4] (N) circle (2);
        \end{scope}
        \draw[gray,thin,dotted] ($(N)+(0:2)$) arc(0:180:2 and 0.5);
        \draw[gray,thin,dashed] ($(N)+(0:2)$) arc(0:-180:2 and 0.5);
        % southern hemisphere
        \begin{scope}
            \clip ($(S)+(0:2)$) arc(0:180:2 and .5) ($(S)+(180:2)$) arc(180:360:2) --cycle;
            \shade[ball color=lightgray,opacity=.4] (S) circle (2);
        \end{scope}
        \draw[gray,thin,dashed] (S) ellipse (2 and .5);
    \end{tikzpicture}
    \quad \approx \mathbb{R}^2 \amalg \mathbb{R}^2, \\
    \left( 
        \begin{tikzpicture}[baseline={([yshift=-.5ex]current bounding box.center)}]
            \path coordinate (O)
            ;
            \shade[ball color = lightgray,opacity = .4] (O) circle (2);
            % 1-strata
            \halfDottedEllipse{(180:2)}{2}{0.5}
            % 0-strata
            \node[bullet] (pt) at ($(O)+(0,-0.5)$) {};
        \end{tikzpicture}
     \right)_{\left[ 
        \begin{tikzpicture}[baseline={([yshift=-.5ex]current bounding box.center)},scale=.25]
            \fill[gray,opacity=.4] (0,0) circle (1);
            \draw[thick] (180:1) -- (0:1);
        \end{tikzpicture}
      \right]}
    &=\quad
    \begin{tikzpicture}[baseline={([yshift=-.5ex]current bounding box.center)}]
        \Ellipse{(0,0)}{2}{.5}
        \fill[white] (2,-.5) circle (3pt);
        \draw[dotted] (2,-.5) circle (3pt);
    \end{tikzpicture}
    \quad \approx \mathbb{R}, \\
    \left( 
        \begin{tikzpicture}[baseline={([yshift=-.5ex]current bounding box.center)}]
            \path coordinate (O)
            ;
            \shade[ball color = lightgray,opacity = .4] (O) circle (2);
            % 1-strata
            \halfDottedEllipse{(180:2)}{2}{0.5}
            % 0-strata
            \node[bullet] (pt) at ($(O)+(0,-0.5)$) {};
        \end{tikzpicture}
     \right)_{
        \left[ 
        \begin{tikzpicture}[baseline={([yshift=-.5ex]current bounding box.center)},scale=.25]
            \fill[gray,opacity=.4] (0,0) circle (1);
            \draw[thick] (180:1) -- (0:1);
            \node[vertex] at (0,0) {};
        \end{tikzpicture}
        \right]
    }
    &= \quad
    \begin{tikzpicture}[baseline={([yshift=-.5ex]current bounding box.center)}]
        \draw[white] (0,0) ellipse (2 and .5);
        \node[bullet] at (0,-.5) {};
    \end{tikzpicture}
\end{align}
の3つである．$X_{[U]}$ 自身は $C^\infty$ 多様体であることに注意しよう．

\begin{mylem}[label=lem:XU-mfd]{$X_{[U]}$ は $C^\infty$ 多様体}
    層状化空間 $(X_{[U]} \to \bigl\{[U]\bigr\}) \in \Obj{\Snglr}$ は $\dim (U) - \depth (U)$ 次元の $C^\infty$ 多様体である．
\end{mylem}

\begin{proof}
    ~\cite[Corollary]{AFT2014stratified}
\end{proof}

ここで，$\BSC_{[U]}$ の\hyperref[def:fullsub-infty]{部分 $(\infty,\, 1)$-圏} $\BSC^0_{[U]}$ を，対象は全く同一だが，\hyperref[def:Map]{射の空間}を原点を保つ層状化開埋め込みに制限したものとして定義する．
すると $\forall V,\, W \in (\BSC^0_{[U]})_0 = (\BSC_{[U]})_0$ に対して，定理\ref{thm:basic-easy}および\eqref{eq:Map-BscU}より\hyperref[def:htpy-equiv-Kan]{Kan複体のホモトピー同値}
\begin{align}
    \inftyMap{\BSC^0_{[U]}}(V,\, W) = \Hom{\BSC}^0 (V,\, W) \simeq \Hom{\BSC} (V,\, W) = \inftyMap{\BSC_{[U]}} (V,\, W)
\end{align}
が成り立つから，\hyperref[def:equiv-infty]{$(\infty,\, 1)$-圏同値}
\begin{align}
    \label{eq:equiv-BscU2BscU0}
    \BSC^0_{[U]} \overset{\simeq}{\hookrightarrow} \BSC_{[U]}
\end{align}
が成り立つ\footnote{包含写像 $\BSC^0_{[U]} \hookrightarrow \BSC_{[U]}$ に命題\ref{prop:equiv-FF-ES}を適用した．}．

さらに，$(\infty,\, 1)$-圏 $\ENTER^0_{[U]}(X)$ を $(\infty,\, 1)$-圏の引き戻し\footnote{\hyperref[def:infty-fib]{右ファイブレーション}は引き戻しについて安定なので $(\tau_X)_|$ もまた右ファイブレーションである．よって補題\ref{lem:pullback-infty}, \ref{lem:pullback-iso}から，この $(\infty,\, 1)$-圏の引き戻しは実際には $\sSet$ の引き戻しとして計算できる．}
\begin{center}
    % https://tikzcd.yichuanshen.de/#N4Igdg9gJgpgziAXAbVABwnAlgFyxMJZABgBoBGAXVJADcBDAGwFcYkQAdDgIQGUBhAHrEA+sGQBVSgF8Q00uky58hFOQrU6TVuy59+YyTLkKQGbHgJF1xTQxZtEnDgFEAcgBUAShPFTpwAAasvKKFipEZLY09jpOXO4eLl7Chv4AFIEAlHKaMFAA5vBEoABmAE4QALZIAMw0OBBIxKEgFdXNDU2I6iA49FiM7AAWEBAA1iAx2o7O2FUwAI5TIIz0AEYwjAAKSpaqIOVYBcM4JmWVNYgATF1IvbGz6Vz9zCLZYgA+IabtV-V9brXaQyIA
    \begin{tikzcd}
    {\ENTER^0_{[U]}(X)} \pullback \arrow[d] \arrow[r]   & {\ENTRU{[U]}{X}} \arrow[d, "(\tau_X)_{|}"] \\
    {\BSC^0_{[U]}} \arrow[r, "\simeq"', hook] & {\BSC_{[U]}}                              
    \end{tikzcd}
\end{center}
により定義する．引き戻しは $(\infty,\, 1)$-圏同値を保存するので，$\ENTER^0_{[U]}(X) \simeq \ENTRU{[U]}{X}$ が成り立つ．
\eqref{eq:EntrU-simplex}より，$\ENTER^0_{[U]}(X)$ の $n$-射は
\begin{align}
    \ENTER^0_{[U]}(X)_n 
    &\cong (\BSC^0_{[U]}{}_{/X})_n \\
    &\cong \Biggl\{\, H \in \Hom{\sSet} \bigl(\Delta^{n+1},\, \hcNer(\SNGLR)\bigr) \Biggm| \substack{H_{[0]}\bigl(\{n+1\}\bigr) = X,\\ H_{[0]} \bigl( \{i\} \bigr) \in [U]\; (0\le \forall i \le n), \\ H|_{\Delta^{n+1}\setminus \{n+1\}}\; \text{preserves}\; \bm{0}} \,\Biggr\} \\
    &= \Biggl\{\, H \in \Hom{\sCat} \bigl(\mathfrak{C}[n+1],\, \SNGLR\bigr) \Biggm| 
        \substack{
            H_{0}(n+1) = X,\\ 
            H_{0} (i) \in [U]\; (0\le \forall i \le n), \\
            H_{i,j} \bigl(\Ner(P_{ij})\bigr) \subset \Hom{\BSC}^0 \bigl( H_0(i),\, H_0(j) \bigr)\; (0 \le \forall i \le \forall j \le n)
            }  \,\Biggr\}
\end{align}
となっている．

ここで，$(\infty,\, 1)$-圏の関手
\begin{align}
    \mathrm{ev}_{\bm{0}} \colon \ENTER^0_{[U]} (X) \lto \Sing^{\text{sm}} (X)
\end{align}
を，$\forall n \ge 0$ に対して
\begin{align}
    \label{eq:eval0}
    \mathrm{ev}_{\bm{0}}{}_{[n]} \colon \ENTER^0_{[U]}(X)_n &\lto \Sing^{\text{sm}}(X) = \Hom{\Snglr} (\Delta^n_e,\, X), \\
    H &\lmto \bigl( x \mapsto \mathrm{proj}_{X} \circ H_{[n+1]}(\Id_{[n+1]})(x,\, \bm{0})\bigr)
\end{align}
と定義する．$H_{[n+1]}(\Id_{[n+1]}) \in \Hom{\SNGLR} \bigl( H_{[0]}(\{0\}),\, X \bigr)_{n} \subset  \Hom{\Snglr} \bigl(\Delta^n_e \times H_{[0]}(\{0\}),\; \Delta^n_e \times X\bigr)$ であることに注意する．
いま $H_{[0]}(\{0\}) \in [U]$ であるから，$\forall x \in \Delta^n_e$ に対して
\begin{align}
    \mathrm{proj}_X \circ H_{[n+1]}(\Id_{[n+1]}) (x,\, \mhyphen \,)
    &\in \Hom{\Snglr} \bigl( H_{[0]} (\{0\}),\, X \bigr) 
\end{align}
である．よって\hyperref[lem:strat-basic]{$\alpha_X$ の構成}より，$\mathrm{proj}_X \circ H_{[n+1]}(\Id_{[n+1]}) (x,\, \bm{0}) \in X_{[U]}$ が言える．
i.e. $(\infty,\, 1)$-圏の関手 $\mathrm{ev}_{\bm{0}}$ は $\Sing^{\text{sm}} (X_{[U]})$ に値をとる．

\begin{mytheo}[label=thm:tangentclassifier]{tangent classifierの局所的な構造}
    $\forall (X \to P) \in \Obj{\Snglr}$ および\hyperref[def:Bsc-induction]{basic} $\forall U = \mathbb{R}^n \times \Cone{Z} \in \Obj{\Bsc}$ を与える．
    このとき，
    \begin{align}
        \Bigl( \ENTRU{[U]}{X} \xrightarrow{(\tau_X)_{|}} \BSC_{[U]} \lto \BSC_{[\mathbb{R}^n]} \Bigr) 
        \simeq \Bigl( \Sing(X_{[U]}) \xrightarrow{\tau_{X_{[U]}}} \deloop{\LO(n)} \Bigr) 
    \end{align}
    が成り立つ．ただし，左辺では命題\ref{prop:Hom-htpy-equiv}を用いて得られるKan複体の\footnote{もしくは $\Kan$-豊穣圏の}ホモトピー同値 $\BSC_{[U]} \simeq \deloop{\Aut(U)} \simeq \deloop{\LO(U)}$ から誘導されるKan複体の射 $\BSC_{[U]} \lto \deloop{\LO(\mathbb{R}^n)}$ を用いた．
\end{mytheo}

\begin{proof}
    
\end{proof}

\subsection{接構造の具体例}

\begin{mydef}[label=def:sieve]{篩}
    \hyperref[def:Kan]{$\Kan$}-\hyperref[def:enriched]{豊穣圏} $\Cat{C}$ を与える．
    このとき充満部分豊穣圏 $\Cat{L} \subset \Cat{C}$ が\textbf{篩} (sieve) であるとは，
    $\forall x \in \Obj{\Cat{C}},\; \forall l \in \Obj{\Cat{L}}$ に対して
    \begin{align}
        \Hom{\Cat{C}} (x,\, l) \neq \emptyset \IMP x \in \Obj{\Cat{L}}
    \end{align}
    が成り立つことをいう．

    \tcblower

    \hyperref[def:infinity-1]{$(\infty,\, 1)$-圏} $\Cat{C}$ を与える．
    このとき\hyperref[def:fullsub-infty]{充満部分 $(\infty,\, 1)$-圏} $\Cat{L} \subset \Cat{C}$ が\textbf{篩} (sieve) であるとは，包含写像が\hyperref[def:infty-fib]{右ファイブレーション}であることを言う．

    もしくは同じ事だが，$\forall x \in \Cat{C}_0,\; \forall l \in \Cat{L}_0$ に対して
    \begin{align}
        \inftyMap{\Cat{C}} (x,\, l) \neq \emptyset \IMP x \in \Cat{L}_0
    \end{align}
    が成り立つことを言う．
\end{mydef}


\section{{$(\infty,\, 1)$}-オペラッドと対称モノイダル {$(\infty,\, 1)$-圏}}

\subsection{{$(\infty,\, 1)$}-オペラッドの定義}

本資料では，$(\infty,\, 1)$-圏を\hyperref[def:infinity-1]{quasi-category}として定義した．
この小節では，quai-categoryにおけるcolored operadを定義する．

\begin{mydef}[label=def:Fin]{{$(1,\, 1)$}-圏 {$\FIN,\; \FIN_*$}}
    $(1,\, 1)$-圏 $\FIN$ を以下で定義する：
    \begin{itemize}
        \item 有限集合および空集合を対象に持つ．
        \item $I,\, J \in \Obj{\FIN}$ の間の写像を射とする．
    \end{itemize}
    \tcblower
    $(1,\, 1)$-圏 $\FIN_*$ を以下で定義する：
    \begin{itemize}
        \item 基点付き有限集合
        \begin{align}
            \expval{n} \coloneqq \{*,\, 1,\, \dots,\, n\}
        \end{align}
        を対象に持つ．i.e.
        \begin{align}
            \Obj{\FIN_*} \coloneqq \bigl\{\, \expval{n} \bigm| n \in \mathbb{Z}_{\ge 0} \,\bigr\}.
        \end{align}
        \item $\forall \expval{m},\, \expval{n} \in \Obj{\FIN_*}$ に対して，それらの間の基点を保つ写像を射とする．i.e.
        \begin{align}
            \Hom{\FIN_*}\bigl( \expval{m},\, \expval{n} \bigr)  \coloneqq \bigl\{\, f \in \Hom{\SETS} \bigl( \expval{m},\, \expval{n} \bigr)  \bigm| f(*) = * \,\bigr\}.
        \end{align}
    \end{itemize}
\end{mydef}

\begin{mydef}[label=def:inert-active]{inert/active morphism}
    \begin{itemize}
        \item \hyperref[def:Fin]{圏 $\FIN_*$} における射 $f \in \Hom{\FIN_*} (\expval{m},\, \expval{n})$ が\textbf{inert}であるとは，$\forall i \in \expval{n} \setminus \{*\}$ に対して $f^{-1} \bigl( \{i\} \bigr) \subset \expval{m}$ が1点集合であることを言う．
        \item \hyperref[def:Fin]{圏 $\FIN_*$} における射 $f \in \Hom{\FIN_*} (\expval{m},\, \expval{n})$ が\textbf{active}であるとは，$f^{-1} \bigl( \{*\} \bigr) = \{*\} \subset \expval{m}$ であることを言う．
    \end{itemize}
\end{mydef}

\begin{myexample}[label=def:rho-i]{inertな射 {$\rho^i$}}
    $1 \le \forall i \le \forall n$ を1つ固定する．このとき，写像
    \begin{align}
        \rho^i \colon \expval{n} &\lto \expval{1}, \\
        j &\lmto 
        \begin{cases}
            1, &j=i \\
            *, &j\neq i
        \end{cases}
    \end{align}
    は圏 $\FIN_*$ における\hyperref[def:inert-active]{inertな射}である．$\rho^i$ は補題\ref{lem:Segal-cond}により，\hyperref[def:infty-operad]{$\infty$-オペラッド}の「射影」
    \begin{align}
        \rho^i_! \colon \Cat{O}^\otimes_{\expval{n}} \lto \Cat{O}^\otimes_{\expval{1}}
    \end{align}
    を規定する．
\end{myexample}

\begin{myexample}[label=def:alpha-n]{activeな射 {$\alpha_n$}}
    $\forall n \ge 1$ を1つ固定する．このとき，写像
    \begin{align}
        \alpha_n \colon \expval{n} &\lto \expval{1}, \\
        j &\lmto 
        \begin{cases}
            1, &j\neq * \\
            *, &j = *
        \end{cases}
    \end{align}
    は圏 $\FIN_*$ における\hyperref[def:inert-active]{activeな射}である．$\alpha_n$ は射の集合 $\Hom{\FIN_*} \bigl( \expval{n},\, \expval{1} \bigr)$ の元のうち，唯一のactiveな射である．
    $\alpha_n$ は，$n$ 点の「テンソル積」を特徴付ける．
\end{myexample}

\hyperref[def:nerve]{脈体の定義}において $[n] \in \Obj{\Delta}$ を $(1,\, 1)$-圏と見做した方法と同様にして，半順序集合 $\{n-1 \le n\}$ を $(1,\, 1)$-圏と見做す．
このとき，
\begin{align}
    \Ner \bigl( \{n-1 \le n\} \bigr) = 
    \left(
        \begin{tikzpicture}[baseline={([yshift=-.5ex]current bounding box.center)}]
            \path coordinate[bullet,label=below:$n-1$] (x)
            +(2,0) coordinate[bullet,label=below:$n$] (y)
            ;
            \draw[->-=.5] (x) -- (y);
        \end{tikzpicture}
     \right) \cong \Delta^1
\end{align}
と図示できる．同様に，
\begin{align}
    \Ner \bigl( \{0 \le 1\} \bigr) = 
    \left(
        \begin{tikzpicture}[baseline={([yshift=-.5ex]current bounding box.center)}]
            \path coordinate[bullet,label=below:$0$] (x)
            +(2,0) coordinate[bullet,label=below:$1$] (y)
            ;
            \draw[->-=.5] (x) -- (y);
        \end{tikzpicture}
     \right) \cong \Delta^1
\end{align}
である．

\begin{mydef}[label=def:p-Cart,breakable]{$p$-Cartesian morphism}
    $p \colon \Cat{E} \lto \Cat{B}$ を\hyperref[def:infty-fib]{内的ファイブレーション}とする．
    \begin{itemize}
        \item $\Cat{E}$ の1-射 $(x \xrightarrow{f} y) \in \Cat{E}_1$ が以下の条件を充たすとき，$f$ は\textbf{$\bm{p}$-Cartesian}であると言う：
        \begin{description}
            \item[\textbf{(Cartesian)}]
            
            $\forall n \ge 2$ に対して，以下の $\sSet$ の図式を可換にする $\Cat{E}$ の $n$-射 $\textcolor{red}{\bar{\varphi}} \in \Hom{\sSet} (\Delta^n,\, \Cat{E}) \cong \Cat{E}_n$ が存在する：
            \begin{center}
                % https://tikzcd.yichuanshen.de/#N4Igdg9gJgpgziAXAbVABwnAlgFyxMJZABgBpiBdUkANwEMAbAVxiRAB12BbOnACwBOXYADkAvpwBGWAOYMAFJ2BgAtAEYABJwYwNYThPbSZAgJRb2AYwIyLAERgMcdAHpqQY0uky58hFGqkalS0jCxsnADCvMAAomIeXiAY2HgERIEATCH0zKyIHOzROMAAQgme3ql+RGTZ1LnhBZwOTq6Elck+af4kQTlh+YUAMnRcklDtAPodITBQMvBEoABmAhBcSGQgOBBIACzUzlgMbHwQEADWiavrm4jbu0iBoXkRRnQCwCsVSWsbByOe0Qak6-3uAGYgUhMg1Bu8VhABIwGBZ6AI0HwsCBqAw6JJHAAFbo1AoCWR8HA3EDgwE7YFQnZ0E5nC7XMF3Z7QxCw15NZLU2mIRlPEG4rBgIaTOB8eY4vlDTgwAAeWDgODgFkkn2AnHRmKwFQoYiAA
                \begin{tikzcd}
                \Ner \bigl(\{n-1 \le n\}\bigr) \cong \Delta^1 \arrow[d, hook] \arrow[rd, "f"] &                        \\
                \Lambda^n_n \arrow[r,blue, "\forall \varphi_0"] \arrow[d, hook]                                                    & \Cat{E} \arrow[d, "p"] \\
                \Delta^n \arrow[r, blue,"\forall \varphi"'] \arrow[ru, red, "\exists \bar{\varphi}", dashed]       & \Cat{B}               
                \end{tikzcd}
            \end{center}
        \end{description}
        ただし，$\sSet$ の射 $f \colon \Ner \bigl(\{n-1 \le n\}\bigr) \lto \Cat{E}$ とは
        \begin{align}
            f_{[1]}
                \left( 
                   \begin{tikzpicture}[baseline={([yshift=-.5ex]current bounding box.center)}]
                       \path coordinate[bullet,label=below:$n-1$] (x)
                       +(1,0) coordinate[bullet,label=below:$n$] (y)
                       ;
                       \draw[->-=.5] (x) -- (y);
                   \end{tikzpicture}
                \right)
            =
                   \begin{tikzpicture}[baseline={([yshift=-.5ex]current bounding box.center)}]
                       \path coordinate[bullet,label=below:$x$] (x)
                       +(1,0) coordinate[bullet,label=below:$y$] (y)
                       ;
                       \draw[->-=.5] (x) --node[midway,above] {$f$} (y);
                   \end{tikzpicture}
        \end{align}
        を充たす唯一の自然変換である．
        \item $\Cat{E}$ の1-射 $(x \xrightarrow{f} y) \in \Cat{E}_1$ が以下の条件を充たすとき，$f$ は\textbf{$\bm{p}$-coCartesian}である\footnote{もしくは，なぜか $p$-coCartesian morphism とは呼ばずに $\bm{p}$\textbf{-coCartesian edge} と呼ぶことが多い．}と言う：
        \begin{description}
            \item[\textbf{(coCartesian)}] 
            
            $\forall n \ge 2$ に対して，以下の $\sSet$ の図式を可換にする $\Cat{E}$ の $n$-射 $\textcolor{red}{\bar{\varphi}} \in \Hom{\sSet} (\Delta^n,\, \Cat{E}) \cong \Cat{E}_n$ が存在する：
            \begin{center}
                % https://tikzcd.yichuanshen.de/#N4Igdg9gJgpgziAXAbVABwnAlgFyxMJZABgBpiBdUkANwEMAbAVxiRAB12BbOnACwBOXYADkAvpwBGWAOYMAFJ2BgAtAEYABJwYwNYThPbSZAgJRb2AYwIyLAERgMcdAHpqQY0uky58hFGqkalS0jCxsnADCvMAAomIeXiAY2HgERIEATCH0zKyIHOzROMAAQgme3ql+RGTZ1LnhBZwOTq6Elck+af4kQTlh+YUAMnRcklDtAPodITBQMvBEoABmAhBcSGQgOBBIACzUzlgMbHwQEADWiavrm4jbu0iBoXkRRnQCwCsVSWsbByOe0Qak6-3uAGYgUhMg1Bu8VhABIwGBZ6AI0HwsCBqAw6JJHAAFbo1AoCWR8HA3EDgwE7YFQnZ0E5nC7XMF3Z7QxCw15NZLU2mIRlPEG4rBgIaTOB8eY4vlDTgwAAeWDgODgFkkn2AnHRmKwFQoYiAA
                \begin{tikzcd}
                \Ner \bigl(\{0 \le 1\}\bigr) \cong \Delta^1 \arrow[d, hook] \arrow[rd, "f"] &                        \\
                \Lambda^n_0 \arrow[r,blue,"\forall \varphi_0"] \arrow[d, hook]                                                    & \Cat{E} \arrow[d, "p"] \\
                \Delta^n \arrow[r, blue,"\forall \varphi"'] \arrow[ru, red, "\exists \bar{\varphi}", dashed]       & \Cat{B}               
                \end{tikzcd}
            \end{center}
            ただし，$\sSet$ の射 $f \colon \Ner \bigl(\{0 \le 1\}\bigr) \lto \Cat{E}$ とは
            \begin{align}
                f_{[1]}
                    \left( 
                       \begin{tikzpicture}[baseline={([yshift=-.5ex]current bounding box.center)}]
                           \path coordinate[bullet,label=below:$n-1$] (x)
                           +(1,0) coordinate[bullet,label=below:$n$] (y)
                           ;
                           \draw[->-=.5] (x) -- (y);
                       \end{tikzpicture}
                    \right)
                =
                       \begin{tikzpicture}[baseline={([yshift=-.5ex]current bounding box.center)}]
                           \path coordinate[bullet,label=below:$x$] (x)
                           +(1,0) coordinate[bullet,label=below:$y$] (y)
                           ;
                           \draw[->-=.5] (x) --node[midway,above] {$f$} (y);
                       \end{tikzpicture}
            \end{align}
            を充たす唯一の自然変換である．
        \end{description}
    \end{itemize}
\end{mydef}

$n=2$ の場合の\hyperref[def:p-Cart]{\textsf{\textbf{(coCartesian)}}}の可換図式の意味を，系\ref{col:horn-coeq}を用いて解読しよう．
まず，包含 $\Ner \bigl( \{0 \le 1\} \bigr) \hookrightarrow \Lambda^2_0$ というのは，系\ref{col:horn-coeq}による角 $\Lambda^2_0$ の図示
\begin{center}
    \begin{tikzpicture} 
        \path coordinate[bullet,label=below left:$\{0\}$] (v0)
        +(2,0) coordinate[bullet,label=below right:$\{2\}$] (v2)
        +(60:2) coordinate[bullet,label=above:$\{1\}$] (v1)
        ;
        \draw[->-=.5] (v0) -- (v1);
        \draw[->-=.5] (v0) -- (v2);
        % \draw[->-=.5] (v1) -- (v2);
    \end{tikzpicture}   
\end{center}
のうち辺 $\{0\} \lto \{1\}$ への埋め込みであるから，可換図式の
\begin{center}
    % https://tikzcd.yichuanshen.de/#N4Igdg9gJgpgziAXAbVABwnAlgFyxMJZABgBpiBdUkANwEMAbAVxiRAB12BbOnACwBOXYADkAvpwBGWAOYMAFJ2BgAtAEYABJwYwNYThPbSZAgJRb2AYwIyLAERgMcdAHpqQY0uky58hFGqkalS0jCxsnADCvMAAomIeXiAY2HgERIEATCH0zKyIHOzROMAAQgme3ql+RGTZ1LnhBZwOTq6Elck+af4kQTlh+YUAMnRcklDtAPodITBQMvBEoABmAhBcSGQgOBBIACzUzlgMbHwQEADWiavrm4jbu0iBoXkRRnQCwCsVSWsbByOe0Qak6-3uAGYgUhMg1Bu8VhABIwGBZ6AI0HwsCBqAw6JJHAAFbo1AoCWR8HA3EDgwE7YFQnZ0E5nC7XMF3Z7QxCw15NZLU2mIRlPEG4rBgIaTOB8eY4vlDTgwAAeWDgODgFkkn2AnHRmKwFQoYiAA
    \begin{tikzcd}
    \Ner \bigl(\{0 \le 1\}\bigr) \cong \Delta^1 \arrow[d, hook] \arrow[rd, "f"] &                        \\
    \Lambda^2_0 \arrow[r,blue,"\forall \varphi_0"]                                                  & \Cat{E}
    \end{tikzcd}
\end{center}
の部分は勝手な角の図式 $\textcolor{blue}{\varphi_0} = (\bullet,\, f,\, \textcolor{blue}{\varphi_0{}_2}) \in \Cat{E}_1^{\times 2}$ を与えることに対応する．
図示すると
\begin{align}
    \label{eq:horn-coCart}
    \textcolor{blue}{\varphi_0}
    &= \begin{tikzpicture}[baseline={([yshift=-.5ex]current bounding box.center)}]
        \path coordinate[bullet,label=below left:$x$] (v0)
        +(2,0) coordinate[bullet,blue] (v2)
        +(60:2) coordinate[bullet,label=above:$y$] (v1)
        ;
        \draw[->-=.5] (v0) --node[midway,above left]{$f$} (v1);
        \draw[->-=.5,blue] (v0) --node[midway,below,blue]{$\varphi_0{}_2$} (v2);
        % \draw[->-=.5] (v1) -- (v2);
    \end{tikzpicture}
\end{align}
となる．従って，\hyperref[def:Cart-coCart]{\textsf{\textbf{(coCart-2)}}}の主張は次のような意味を持つ：

$(\infty,\, 1)$-圏 \underline{$\Cat{B}$ において}角の図式が2-射 $\textcolor{blue}{\varphi} \in \Hom{\sSet}(\Delta^2,\, \Cat{B}) \cong \Cat{B}_2$ によって
\begin{center}
    \begin{tikzpicture} 
        \path coordinate[bullet,label=below left:$p_{[0]}(x)$] (v0)
        +(2,0) coordinate[bullet,blue] (v2)
        +(60:2) coordinate[bullet,label=above:$p_{[0]}(y)$] (v1)
        ;
        \draw[->-=.5] (v0) --node[midway,above left]{$p_{[1]}(f)$} (v1);
        \draw[->-=.5,blue] (v0) --node[midway,below,blue]{$p_{[1]}(\varphi_0{}_2)$} (v2);
        \draw[->-=.5,blue] (v1) --node[midway,above right,blue]{$\partial^2_0(\varphi)$} (v2);

        \node[blue] at (barycentric cs:v0=1,v1=1,v2=1) {$\varphi$};

        \begin{scope}[on background layer]
            \fill[blue!20] (v0.center) -- (v1.center) -- (v2.center) --cycle;
        \end{scope}
        % \draw[->-=.5] (v1) -- (v2);
    \end{tikzpicture}   
\end{center}
と埋められているならば，\underline{$\Cat{E}$ において}角の図式\eqref{eq:horn-coCart}を
\begin{center}
    \begin{tikzpicture} 
        \path coordinate[bullet,label=below left:$x$] (v0)
        +(2,0) coordinate[bullet,blue] (v2)
        +(60:2) coordinate[bullet,label=above:$y$] (v1)
        ;
        \draw[->-=.5] (v0) --node[midway,above left]{$f$} (v1);
        \draw[->-=.5,blue] (v0) --node[midway,below,blue]{$\varphi_0{}_2$} (v2);
        \draw[->-=.5,red] (v1) --node[midway,above right,red]{$\partial^2_0(\bar{\varphi})$} (v2);
        \node[red] at (barycentric cs:v0=1,v1=1,v2=1) {$\bar{\varphi}$};

        \begin{scope}[on background layer]
            \fill[red!20] (v0.center) -- (v1.center) -- (v2.center) --cycle;
        \end{scope}
        % \draw[->-=.5] (v1) -- (v2);
    \end{tikzpicture}   
\end{center}
のように埋める2-射 $\textcolor{red}{\bar{\varphi}} \in \Hom{\sSet}(\Delta^2,\, \Cat{E}) \cong \Cat{E}_2$ が存在する．

\begin{mydef}[label=def:Cart-coCart,breakable]{デカルトファイブレーション}
    $p \colon \Cat{E} \lto \Cat{B}$ を\hyperref[def:infty-fib]{内的ファイブレーション}とする．
    \begin{itemize}
        \item $p$ が\textbf{デカルトファイブレーション} (Cartesian fibration) であるとは，
        \begin{itemize}
            \item $\Cat{B}$ の任意の1射 $(x \xrightarrow{f} y) \in \Cat{B}_1$ 
            \item $p_{[0]}(\bar{y}) = y$ を充たす $\Cat{E}$ の任意の対象 $\bar{y} \in \Cat{E}_0$
        \end{itemize}
        に対して，以下の条件を充たす $\Cat{E}$ の1-射 $(z \xrightarrow{\bar{f}} \bar{y}) \in \Cat{E}_1$ が存在することを言う：
        \begin{description}
            \item[\textbf{(Cart-1)}]　
            
            $\bar{f}$ は $f$ の持ち上げである．i.e. $p_{[1]}(\bar{f}) = f$ が成り立つ．

            \item[\textbf{(Cart-2)}]　
            
            $\bar{f}$ は\hyperref[def:p-Cart]{$p$-Cartesian}である．
        \end{description}
        特に，$\bar{f}$ のことを $f$ の $\bm{p}$\textbf{-Cartesian lift}と呼ぶ．
        \item $p$ が\textbf{余デカルトファイブレーション} (coCartesian fibration) であるとは，
        \begin{itemize}
            \item $\Cat{B}$ の任意の1射 $(x \xrightarrow{f} y) \in \Cat{B}_1$ 
            \item $p_{[0]}(\bar{x}) = x$ を充たす $\Cat{E}$ の任意の対象 $\bar{x} \in \Cat{E}_0$
        \end{itemize}
        に対して，以下の条件を充たす $\Cat{E}$ の1-射 $(\bar{x} \xrightarrow{\bar{f}} z) \in \Cat{E}_1$ が存在することを言う：
        \begin{description}
            \item[\textbf{(coCart-1)}]　
            
            $\bar{f}$ は $f$ の持ち上げである．i.e. $p_{[1]}(\bar{f}) = f$ が成り立つ．

            \item[\textbf{(coCart-2)}]　
            
            $\bar{f}$ は\hyperref[def:p-Cart]{$p$-coCartesian}である．
        \end{description}
        特に，$\bar{f}$ のことを $f$ の $\bm{p}$\textbf{-coCartesian lift}と呼ぶ．
    \end{itemize}
\end{mydef}

\begin{mydef}[label=def:infty-operad,breakable]{{$(\infty,\, 1)$}-オペラッド}
    $\bm{(\infty,\, 1)}$\textbf{-オペラッド} ($(\infty,\, 1)$-operad)\footnote{$\infty$-operadとも呼ばれる．} とは，
    \hyperref[def:infinity-1]{$(\infty,\, 1)$-圏の関手}
    \begin{align}
        p \colon \Cat{O}^{\otimes} \lto \Ner (\FIN_*)
    \end{align}
    であって以下の条件を充たすもののこと~\cite[Definition 2.1.1.10.]{Lurie2017HA}：
    \begin{description}
        \item[\textbf{(Op-1)}]　
        
        任意の\hyperref[def:inert-active]{inertな}射 $\alpha \in \Hom{\FIN_*} \bigl( \expval{m},\, \expval{n} \bigr)$ および $\forall x \in (\Cat{O}^\otimes_{\expval{m}})_0$ に対して，
        $\Cat{O}^\otimes$ における\hyperref[def:p-Cart]{$p$-coCartesian}な1-射 $(x \xrightarrow{\bar{\alpha}} y) \in (\Cat{O}^\otimes)_1$ が存在して $p_{[1]}(\bar{\alpha}) = \alpha$ を充たす．
        
        \item[\textbf{(Op-2)}]　
        
        $\forall \alpha \in \Hom{\FIN_*} \bigl( \expval{m},\, \expval{n} \bigr)$ および $\forall x \in (\Cat{O}^\otimes_{\expval{m}})_0,\; \forall y \in (\Cat{O}^\otimes_{\expval{n}})_0$ に対して，
        \exref{def:rho-i}のinertな射の族 $\Familyset[\big]{\rho^i \in  \Hom{\FIN_*}\bigl( \expval{n},\, \expval{1} \bigr)}{1 \le i \le n}$ に \textsf{\textbf{(Op-1)}} を適用して得られる\hyperref[def:p-Cart]{$p$-coCartesian}な1-射の族 $\Familyset[\big]{y \xrightarrow{\bar{\rho}^i} y_i \in (\Cat{O}^\otimes)_1}{1 \le i \le n}$ が誘導する\hyperref[def:infinity-1]{Kan複体の関手}
        \footnote{つまり，これは $(\bar{\rho}^1,\, \dots,\, \bar{\rho}^n) \in \prod_{i=1}^n \inftyMap{\Cat{O}^\otimes} (y,\, y_i)_0$ による\hyperref[def:Map-comp]{後合成}である．この関手は up to homotopy でしか定まらないのだった．}
        \begin{align}
            \inftyMap{\Cat{O}^\otimes} (x,\, y)_\alpha \lto \prod_{i=1}^n \inftyMap{\Cat{O}^\otimes} (x,\, y_i)_{\rho^i \circ \alpha}
        \end{align}
        は\hyperref[def:htpy-equiv]{ホモトピー同値}である．

        \item[\textbf{(Op-3)}]　
        
        $\forall x_1,\, \dots,\, x_n \in (\Cat{O}^\otimes_{\expval{1}})_0$ に対して，
        ある $x \in (\Cat{O}^\otimes_{\expval{n}})_0$ および \hyperref[def:p-Cart]{$p$-coCartesian}な1-射 $\widehat{\rho_i} \in \bigl(\inftyMap{\Cat{O}^\otimes} (x,\, x_i)_{\rho^i}\bigr)_0$ が存在する．
    \end{description}

    ただし，以下の記法を採用した：
    \begin{itemize}
        \item 
        点 $\expval{n} \in \Ner(\FIN_*)_0$ における\hyperref[def:infinity-1]{$(\infty,\, 1)$-圏の関手} $p \colon \Cat{O}^\otimes \lto \Ner (\FIN_*)$ の\textbf{ファイバー} $\bm{\Cat{O}}^{\bm{\otimes}}_{\expval{n}}$ を，
        \hyperref[prop:SimpSet-basic]{$(1,\, 1)$-圏 $\sSet$ における引き戻し}
        \begin{center}
            % https://tikzcd.yichuanshen.de/#N4Igdg9gJgpgziAXAbVABwnAlgFyxMJZABgBpiBdUkANwEMAbAVxiRAB12BhOnYAeQC+APWCcIeALbxBAfTHsYADzT0GwMIMEhBpdJlz5CKMgEYqtRizacefIcPFT4OvSAzY8BIqdLnqataIHOySvAAWAE6SwAByggAUnABiAJKxsgBUAJSu+p5GPuQWgazBnAAiMAw4dMLEOhYwUADmLiigAGaREJJIZCA4EEgAzLpdPX2IA0NIpuMg3b2j1LOIAEyrdFgMbOEQEADWIAFWZSHKqowa2gtLU76DwxunzOdoJyAMdABG1QAKBi8xhAkSwLXCOEagiAA
            \begin{tikzcd}
            \textcolor{red}{\Cat{O}^{\otimes}_{\expval{n}}} \arrow[dr, phantom," " {pullback}, very near start, color=black] \arrow[red, r] \arrow[red,d] & \Delta^0 \arrow[d, "\expval{n}", hook] \\
            \Cat{O}^\otimes \arrow[r, "p"']                    & \Ner (\FIN_*)                    
            \end{tikzcd}
        \end{center}
        と定義した
        \footnote{
            これ自体が $(\infty,\, 1)$-圏である．
        }．
        \item $\alpha \in \Hom{\FIN_*} \bigl( \expval{m},\, \expval{n} \bigr)$ および $x \in (\Cat{O}^\otimes_{\expval{m}})_0,\; y \in (\Cat{O}^\otimes_{\expval{n}})_0$ に対して，
        $\bm{\alpha}$ \textbf{上の射の空間} $\bm{\inftyMap{\Cat{O}^\otimes} (x,\, y)_\alpha}$ を，\hyperref[prop:SimpSet-basic]{$(1,\, 1)$-圏 $\sSet$ における引き戻し}
        \begin{center}
            % https://tikzcd.yichuanshen.de/#N4Igdg9gJgpgziAXAbVABwnAlgFyxMJZABgBpiBdUkANwEMAbAVxiRAB12swAzHATwCydNME4BhOjmAB5AL4A9ThDwBbeHIAEACgAepTqU38AlAH1OjNAAs6IOaXSZc+QijIBGKrUYs2nbj4hETF2SWl5JXYVLHU4OT0DdiNTe0cQDGw8AiIPcm96ZlZEDnYAERgGHDoFYjSnLNdc0i9qQr8SgN4BYVFOADkYACdtTgAxbjMAKhM5TgAjLABzBm00M2ASCgTdEySjdc3KBNMF5aGTe28YKCX4IlAeIYhVJAAmahwIJABmT7osAw2NYIBAANYgNq+YqlKy2eogJ4vJB5EBfX5QopsNCQkAMOjzSoABWc2TcICGy2sOARSNeiDIaO+iA8Dkez3pjPRiDecm2QA
            \begin{tikzcd}
            {\textcolor{red}{\inftyMap{\Cat{O}^\otimes} (x,\, y)_\alpha}} \arrow[dr, phantom," " {pullback}, very near start, color=black] \arrow[red,d] \arrow[red,r] & \Delta^0 \arrow[d, "\alpha", hook]                            \\
            {\inftyMap{\Cat{O}^\otimes}(x,\, y)} \arrow[r, "p"']             & {\inftyMap{\Ner(\FIN_*)} \bigl(\expval{m},\, \expval{n}\bigr)}
            \end{tikzcd}
        \end{center}
        と定義した
        \footnote{
            これ自体がKan複体である．なお，下の横向きの射は，定義\ref{def:FF-ES-infty}に登場したものである．
        }\cite[\href{https://kerodon.net/tag/01P2}{Tag 01P2}]{kerodon}．
        % $\bar{\alpha} \in \inftyMap{\Cat{O}^\otimes}(x,\, y)_{0}$ であって $p_{[1]}(\bar{\alpha}) = \alpha$ を充たすもの全体から定まる，$\inftyMap{\Cat{O}^\otimes} (x,\, y) \in \SPACES_0$
        % の
        % % \footnote{\hyperref[def:infty-1]{Kan複体} $\Hom{\Cat{O}^\otimes} (c,\, c') \in \SPACES_0$ は，\hyperref[def:infty-Hom]{$(\infty,\, 1)$-圏におけるHom関手} $\inftyMap{\Cat{O}^\otimes} \colon \OP{(\Cat{O}^\otimes)} \times \Cat{O}^\otimes \lto \SPACES$ による対象 $(c,\, c') \in (\OP{(\Cat{O}^\otimes)} \times \Cat{O}^\otimes)_0$ の像である．}
        % \hyperref[def:fullsub-infty]{充満部分Kan複体}を
        % \begin{align}
        %     \bm{\inftyMap{\Cat{O}^\otimes} (x,\, y)_\alpha} \hookrightarrow \inftyMap{\Cat{O}^\otimes} (x,\, y)
        % \end{align}
        % と書いた．
    \end{itemize}
\end{mydef}

\begin{mylem}[label=lem:Segal-cond]{Segal条件}
    \hyperref[def:infinity-1]{$(\infty,\, 1)$-圏の関手}
    \begin{align}
        p \colon \Cat{O}^{\otimes} \lto \Ner (\FIN_*)
    \end{align}
    であって条件\hyperref[def:infty-operad]{\textsf{\textbf{(Op-1)}}, \textsf{\textbf{(Op-2)}}}を充たすものを与える．
    このとき，条件\hyperref[def:infty-operad]{\textsf{\textbf{(Op-3)}}}は以下と同値である
    \footnote{一般にSegal条件というと $n$-fold pullbackのことだが，現在は $\Cat{O}^\otimes_{\expval{0}}$ がcontractibleなので $n$-fold productになっている．}
    ：
    \begin{description}
        \item[\textbf{(Segal)}]　
        
        $\forall n \ge 0$ に対して，\exref{def:rho-i}の\hyperref[def:inert-active]{inertな射}の族 $\Familyset[\big]{\rho^i \in \Hom{\FIN_*} \bigl( \expval{n},\, \expval{1} \bigr)}{1 \le i \le n}$ が条件\hyperref[def:infty-operad]{\textsf{\textbf{(Op-1)}}}により誘導する\hyperref[def:infinity-1]{$(\infty,\, 1)$-圏の関手}の族 $\Familyset[\big]{\rho^i_{!} \colon \Cat{O}^\otimes_{\expval{n}} \lto \Cat{O}^\otimes_{\expval{1}}}{1 \le i \le n}$ は，\hyperref[def:equiv-infty]{$(\infty,\, 1)$-圏同値}
        \begin{align}
            \label{eq:Segalmap}
            (\rho^1_{!},\, \dots,\, \rho^n_{!}) \colon \Cat{O}^\otimes_{\expval{n}} \xrightarrow{\simeq} \bigl( \Cat{O}^\otimes_{\expval{1}} \bigr)^{\times n}
        \end{align}
        を与える．
    \end{description}
\end{mylem}

\begin{proof}
    $\forall x \in (\Cat{O}^\otimes_{\expval{n}})_0$ は，定義から $x \in \Cat{O}^\otimes_0$ かつ $p_{[0]}(x) = \expval{n}$ を充たす．故に\hyperref[def:infty-operad]{\textsf{\textbf{(Op-1)}}}から，inertな射の族 $\Familyset[\big]{\rho^i \in  \Hom{\FIN_*}\bigl( \expval{n},\, \expval{1} \bigr)}{1 \le i \le n}$ の\hyperref[def:p-Cart]{$p$-coCartesian lift} $\Familyset[\big]{x \xrightarrow{\bar{\rho}^i_x} x_i \in (\Cat{O}^\otimes)_1}{1 \le i \le n}$ が存在する．

    まず，$(\infty,\, 1)$-圏の関手\eqref{eq:Segalmap}が具体的にどのように構成されるのかを見よう．

    \begin{description}
        \item[\textbf{0-射}]　
        
        $p_{[0]} (x_i) = p_{[0]} (\partial^1_0 \bar{\rho}^i) = \partial^1_0 p_{[1]}(\bar{\rho}^i) = \partial^1_0 \rho^i = \expval{1}$ より $x_i \in (\Cat{O}^\otimes_{\expval{1}})_0$ である．したがって
        \begin{align}
            \rho^i_{!}{}_{[0]} \colon (\Cat{O}^\otimes_{\expval{n}})_0 &\lto (\Cat{O}^\otimes_{\expval{1}})_0, \\
            x &\lmto x_i
        \end{align}
        と定義する．

        \item[\textbf{1-射}]　
        
        $\Cat{O}^\otimes_{\expval{n}}$ の定義から，$\forall (x \xrightarrow{f} y) \in (\Cat{O}^\otimes_{\expval{n}})_1 \subset \Cat{O}^\otimes_1$ に対して $p_{[1]} (f) = \Id_{\expval{n}}$ が成り立つ．
        i.e. $\rho^i = \bigl(\expval{n} \xrightarrow{\Id_{\expval{n}}} \expval{n} \xrightarrow{\rho^i} \expval{1}\bigr) \in \Ner(\FIN_*)_2$ について\footnote{\hyperref[def:nerve]{脈体}の $2$-射の記法\eqref{eq:nerve-spine}を用いた．}，$(1,\, 1)$-圏 $\sSet$ の可換図式
        % ところで，$\Cat{O}^\otimes$ は\hyperref[def:infinity-1]{$(\infty,\, 1)$-圏}だから，ある $\sigma_f \in \Cat{O}^\otimes_2$ が存在して以下の $\sSet$ の図式を可換にする：
        \begin{center}
            % https://tikzcd.yichuanshen.de/#N4Igdg9gJgpgziAXAbVABwnAlgFyxMJZABgBpiBdUkANwEMAbAVxiRAB12AZOgWwCModAHoAmAPoBGEAF9S6TLnyEUk8lVqMWbTgGE6OYAHkZwzhDy94s+SAzY8BImUkb6zVog7sAIjAY4IqKyGjBQAObWKKAAZgBOELxIZCA4EEii1IFYDGwAFhAQANY2sQlJiJmp6YhqIAx0-P4ACoqOKiBxWOF5OCDU7tpenDAAHlhwOHAABJzY4bx0pSDxiclZNXWDniAAFDGknKSz7PxMDAwwOIfsx5z8dHHAnHEFpljiAJ4AlCEyQA
            \begin{tikzcd}
            \Lambda^2_1 \arrow[d, hook] \arrow[r, "{(f,\, \bullet,\, \bar{\rho}^i_y)}"] & \Cat{O}^\otimes \ar[d,"p"]\\
            \Delta^2 \arrow[r, "{\rho^i}"']     &\Ner(\FIN_*)
            \end{tikzcd}
        \end{center}
        が成り立つ．そのうえ $p \colon \Cat{O}^\otimes \lto \Ner(\FIN_*)$ は\hyperref[def:infty-fib]{内的ファイブレーション}だから
        \begin{center}
            % https://tikzcd.yichuanshen.de/#N4Igdg9gJgpgziAXAbVABwnAlgFyxMJZABgBpiBdUkANwEMAbAVxiRAB12AZOgWwCModAHoAmAPoBGEAF9S6TLnyEUk8lVqMWbTgGE6OYAHkZwzhDy94s+SAzY8BImUkb6zVog7sAIjAY4IqKyGjBQAObWKKAAZgBOELxIZCA4EEii1IFYDGwAFhAQANY2sQlJiJmp6YhqIAx0-P4ACoqOKiBxWOF5OCDU7tpenDAAHlhwOHAABJzY4bx0pSDxiclZNXWDniAAFDGknKSz7PxMDAwwOIfsx5z8dHHAnHEFpljiAJ4AlCEyQA
            \begin{tikzcd}
            \Lambda^2_1 \arrow[d, hook] \arrow[r, "{(f,\, \bullet,\, \bar{\rho}^i_y)}"] & \Cat{O}^\otimes \ar[d,"p"]\\
            \Delta^2 \arrow[r, "{\rho^i}"']  \ar[ur, red, dashed, "\exists \sigma_f"]   &\Ner(\FIN_*)
            \end{tikzcd}
        \end{center}
        である．このとき $f^i \coloneqq \partial^2_1 (\sigma_f)$ とおくと，$p$ が\hyperref[def:nat]{自然変換}であることと\hyperref[def:SimpSet]{単体的恒等式}から
        \begin{align}
            p_{[1]} (f^i) 
            &= p_{[1]} \circ \partial^2_1 (\sigma_f)
            = \partial^2_1 \circ p_{[1]}(\sigma_f)
            = \rho^i, \\
            \partial^1_1 (f^i) 
            % &= \partial^1_1 \circ \partial^2_1 (f^i) \\
            % &= \partial^1_1 \circ \partial^2_2 (f^i) \\
            % &= \partial^1_1 (f) \\
            &= x, \\
            \partial^1_0 (f^i)
            &= y^i, \\
            \partial^1_1 \bigl( p_{[1]} (f^i) \bigr) 
            % &= p_{[0]} \circ \partial^1_1 \circ \partial^2_1 (\sigma_f) \\
            % &= p_{[0]} \circ \partial^1_1 \circ \partial^2_2 (\sigma_f) \\
            % &= p_{[0]} \circ \partial^1_1 (f) \\
            &= \expval{n}, \\
            \partial^1_0 \bigl( p_{[1]} (f^i) \bigr) 
            % &= p_{[0]} \circ \partial^1_0 \circ \partial^2_1 (\sigma_f) \\
            % &= p_{[0]} \circ \partial^1_0 \circ \partial^2_0 (\sigma_f) \\
            % &= p_{[0]} \circ \partial^1_0 (\bar{\rho}^i) \\
            &= \expval{1}
        \end{align}
        が成り立つ．
        i.e. $(\infty,\, 1)$-圏 $\Cat{O}^\otimes$ において
        \begin{center}
            \begin{tikzpicture} 
                \path coordinate[bullet,label=below left:$x$] (v0)
                +(2,0) coordinate[bullet,label=below right:$y_i$] (v2)
                +(60:2) coordinate[bullet,label=above:$x_i$] (v1)
                ;
                \draw[->-=.5] (v0) --node[midway,above left]{$\bar{\rho}^i_x$} (v1);
                \draw[->-=.5] (v0) --node[midway,below]{$f^i$} (v2);
                % \draw[->-=.5,red] (v1) --node[midway,above right,red]{$\partial^2_0(\bar{\varphi})$} (v2);
                % \node[red] at (barycentric cs:v0=1,v1=1,v2=1) {$\bar{\varphi}$};            
                % \begin{scope}[on background layer]
                %     \fill[red!20] (v0.center) -- (v1.center) -- (v2.center) --cycle;
                % \end{scope}
                % \draw[->-=.5] (v1) -- (v2);
            \end{tikzpicture}   
        \end{center}
        と図示できる．
        一方，$\varphi \coloneqq \bigl(\expval{n} \xrightarrow{\rho^i} \expval{1} \xrightarrow{\Id_{\expval{1}}} \expval{1}\bigr) \in \Ner(\FIN_*)_2$ は $(\infty,\, 1)$-圏 $\Ner(\FIN_*)$ において
        \begin{center}
            \begin{tikzpicture} 
                \path coordinate[bullet,label=below left:$\expval{n}$] (v0)
                +(2,0) coordinate[bullet,label=below right:$\expval{1}$] (v2)
                +(60:2) coordinate[bullet,label=above:$\expval{1}$] (v1)
                ;
                \draw[->-=.5] (v0) --node[midway,above left]{$p_{[1]}(\bar{\rho}^i_x) = \rho^i$} (v1);
                \draw[->-=.5] (v0) --node[midway,below]{$p_{[1]}(f^i) = \rho^i$} (v2);
                \draw[->-=.5] (v1) --node[midway,above right]{$\Id_{\expval{1}}$} (v2);
            
                \node[gray] at (barycentric cs:v0=1,v1=1,v2=1) {$\varphi$};
            
                \begin{scope}[on background layer]
                    \fill[gray!20] (v0.center) -- (v1.center) -- (v2.center) --cycle;
                \end{scope}
                % \draw[->-=.5] (v1) -- (v2);
            \end{tikzpicture}   
        \end{center}
        と図示できる．故に，$\bar{\rho}^i$ が\hyperref[def:p-Cart]{$p$-coCartesian}であることから $\textcolor{red}{\bar{\varphi}_f} \in \Cat{O}^\otimes_2$ が存在して
        \begin{center}
           \begin{tikzpicture} 
               \path coordinate[bullet,label=below left:$x$] (v0)
               +(2,0) coordinate[bullet,label=below right:$y^i$] (v2)
               +(60:2) coordinate[bullet,label=above:$x^i$] (v1)
               ;
               \draw[->-=.5] (v0) --node[midway,above left]{$\bar{\rho}^i_x$} (v1);
               \draw[->-=.5] (v0) --node[midway,below]{$f^i$} (v2);
               \draw[->-=.5,red] (v1) --node[midway,above right,red]{$\partial^2_0(\bar{\varphi}_f)$} (v2);
               \node[red] at (barycentric cs:v0=1,v1=1,v2=1) {$\bar{\varphi}_f$};

               \begin{scope}[on background layer]
                   \fill[red!20] (v0.center) -- (v1.center) -- (v2.center) --cycle;
               \end{scope}
               % \draw[->-=.5] (v1) -- (v2);
           \end{tikzpicture}   
        \end{center}
        となる．以上の構成より，
        \begin{align}
            \rho^i_{!}{}_{[1]} \colon (\Cat{O}^\otimes_{\expval{n}})_1 &\lto (\Cat{O}^\otimes_{\expval{1}})_1, \\
            (x \xrightarrow{f} z) &\lmto (x_i \xrightarrow{\partial^2_0(\bar{\varphi}_f)} y_i)
        \end{align}
        と定義する．この定義は $\sigma_f$ の選び方および $\bar{\varphi}$ の選び方の不定性がある．これらの不定性は
        それぞれ $\bar{\rho}^i_y$ の\hyperref[def:Map-comp]{後合成}および $\bar{\rho}^i_x$ の\hyperref[def:Map-comp]{前合成}が up to homotopy でしか定まらないという事実に対応している．

        \item[\textbf{$m \ge 1$-射}]　
        
        $\forall x,\, y \in (\Cat{O}^\otimes_{\expval{n}})_0$ を固定する．
        定義から，$\forall p \ge 1$ に対して
        \begin{align}
            \inftyMap{\Cat{O}^\otimes_{\expval{p}}} (x,\, y) = \inftyMap{\Cat{O}^\otimes} (x,\, y)_{\Id_{\expval{p}}}
        \end{align}
        が成り立つことに注意する．すると，\hyperref[def:Map-comp]{前合成および後合成}により，Kan複体の関手
        \begin{align}
            \label{def:Segal-map}
            \inftyMap{\Cat{O}^\otimes_{\expval{n}}} (x,\, y) = \inftyMap{\Cat{O}^\otimes}(x,\, y)_{\Id_{\expval{n}}} &\xrightarrow{\prod_{i=1}^n [\bar{\rho}^i_y]_*} \prod_{i=1}^n \inftyMap{\Cat{O}^\otimes} (x,\, y_i)_{\rho^i} \\ 
            &\xleftarrow{\prod_{i=1}^n [\bar{\rho}^i_x]^*} \prod_{i=1}^n \inftyMap{\Cat{O}^\otimes} (x_i,\, y_i)_{\Id_{\expval{1}}} = \prod_{i=1}^n \inftyMap{\Cat{O}^\otimes_{\expval{1}}} (x_i,\, y_i)
        \end{align}
        が up to homotopy で定まる．
        % \hyperref[def:infty-operad]{\textsf{\textbf{(Op-2)}}}より $\prod_{i=1}^n [\bar{\rho}^i_y]_*$ は\hyperref[def:htpy-equiv]{ホモトピー同値}である．
        \hyperref[def:infty-operad]{\textsf{\textbf{(Op-1)}}}より $\bar{\rho}_x^i$ が $p$-coCartesianなので，
        $[\bar{\rho}^i_x]^* \colon \inftyMap{\Cat{O}^\otimes} (x_i,\, y_i)_{\Id_{\expval{1}}} \lto \inftyMap{\Cat{O}^\otimes} (x,\, y_i)_{\rho^i}$
        は\hyperref[def:htpy-equiv]{ホモトピー同値}であり\cite[\href{https://kerodon.net/tag/01U8}{Tag 01U8}]{kerodon}，
        up to homotopy で向きを逆にできる．このようにして $\forall m \ge 1$ 射の対応が定まる．
        % このとき\exref{ex:Map-m1}より自然同型
        % \begin{align}
        %     \inftyMap{\Cat{O}^\otimes_{\expval{n}}} (x,\, y)_m \cong \Bigl\{\, (\sigma_0,\, \dots,\, \sigma_m) \in (\Cat{O}^\otimes_{\expval{n}})_{m+1}^{\times (m+1)} \Bigm| \substack{1 \le \forall k \le m,\; \partial^{m+1}_k (\sigma_k) = \partial^{m+1}_k (\sigma_{k-1}) \\ \partial^{m+1}_{m+1} (\sigma_m) = \sigma^{m-1}_0 \cdots \sigma^{0}_0(x),\; \partial^{m+1}_0 (\sigma_0) = \sigma^{m-1}_0 \cdots \sigma^{0}_0(y)} \,\Bigr\}
        % \end{align}
        % があるから，結局 $\sigma \in (\Cat{O}^\otimes_{\expval{n}})_{m+1}$ の行き先が指定できれば良い．
    \end{description}
    
    \begin{description}
        \item[\textbf{(OP-3)} $\Longrightarrow$ \textbf{(Segal)}]　
        
        　\hyperref[def:infty-operad]{\textsf{\textbf{(OP-3)}}}を仮定する．命題\ref{prop:equiv-FF-ES}より，$(\infty,\, 1)$-圏の関手\eqref{eq:Segalmap}が\hyperref[def:FF-ES]{忠実充満かつ本質的全射}であることを示せば良い．
        まず，\hyperref[def:infty-operad]{\textsf{\textbf{(OP-2)}}}により\eqref{def:Segal-map}は\hyperref[def:htpy-equiv]{ホモトピー同値}だから，\eqref{eq:Segalmap}は\hyperref[def:FF-ES-infty]{忠実充満}である．

        \eqref{eq:Segalmap}が\hyperref[def:FF-ES-infty]{本質的全射}であることは\hyperref[def:infty-operad]{\textsf{\textbf{(OP-3)}}}より従う．

        \item[\textbf{(OP-3)} $\Longleftarrow$ \textbf{(Segal)}]　
        明らか．
    \end{description}
\end{proof}


\begin{mydef}[label=def:infty-operad-map]{{$(\infty,\, 1)$}-オペラッドの射}
    \hyperref[def:infty-operad]{$(\infty,\, 1)$-オペラッド} $p \colon \Cat{O}^{\otimes} \lto \Ner (\FIN_*)$ を与える．
    このとき，$(\infty,\, 1)$-圏 $\Cat{O}^\otimes$ の1-射 $f \in (\Cat{O}^\otimes)_1$ が\textbf{inert}であるとは，以下の2条件を充たすことを言う：
    \begin{description}
        \item[\textbf{(inert-1)}] $p_{[1]}(f)$ は $(1,\, 1)$-圏 $\FIN_*$ における\hyperref[def:inert-active]{innerな射}である．
        \item[\textbf{(inert-2)}] $f$ は\hyperref[def:p-Cart]{$p$-coCartesian}な1-射である．
    \end{description}
    
    \tcblower
    
    2つの\hyperref[def:infty-operad]{$(\infty,\, 1)$-オペラッド} $p \colon \Cat{O}^{\otimes} \lto \Ner (\FIN_*),\; p' \colon \Cat{O}'{}^\otimes \lto \Ner (\FIN_*)$ を与える．
    このとき，\hyperref[def:infinity-1]{$(\infty,\, 1)$-圏の関手} $f \colon \Cat{O}^{\otimes} \lto \Cat{O}'{}^\otimes$ が\textbf{$\bm{(\infty,\, 1)}$-オペラッドの射} ($\infty$-operad map) であるとは，以下の2つの条件を充たすことを言う~\cite[Definition 2.1.2.7.]{Lurie2017HA}：
    \begin{description}
        \item[\textbf{(Opmap-1)}] $(1,\, 1)$-圏 $\sSet$ の図式
        \begin{center}
            % https://tikzcd.yichuanshen.de/#N4Igdg9gJgpgziAXAbVABwnAlgFyxMJZABgBpiBdUkANwEMAbAVxiRAB12BhOnYAeQC+APU4Q8AW3ghBpdJlz5CKAEzkqtRizacefIQHJgIsZOmz52PASIBGUrY31mrRB3YA5GACcABAApOADEASQ8AfQAqAEoZDRgoAHNpFFAAM28ICSQyEBwIJHtNFzY0mTkQDKyc6nykNWLtNzQQagY6ACMYBgAFBWtlEG8sRIALHHL0zOzEIrrEBucmkABHOMEgA
            \begin{tikzcd}
            \Cat{O}^\otimes \arrow[rr, "f"] \arrow[rd, "p"']    &               & \Cat{O}'{}^\otimes \arrow[ld, "p'"] \\
                                                                & \Ner (\FIN_*) &                                   
            \end{tikzcd}
        \end{center}
        は可換である．
        \item[\textbf{(Opmap-2)}] 1-射の間の写像 $f_{[1]} \colon (\Cat{O}^{\otimes})_1 \lto (\Cat{O}'{}^{\otimes})_1$ により，inertな1-射が保存される\footnote{条件 \textsf{\textbf{(Opmap-1)}} より，inertな1-射の $p_{[1]}$ による像が条件 \textsf{\textbf{(inert-1)}} を充たすことは明らかである．}．
    \end{description}
\end{mydef}

\begin{marker}
    以降では\hyperref[def:infty-operad]{$(\infty,\, 1)$-オペラッド} $\Cat{O}^\otimes$ に対して $\Cat{O} \coloneqq \Cat{O}^{\otimes}_{\expval{1}}$ とおき，$\Cat{O}^\otimes$ の \textbf{underlying} $\bm{(\infty,\, 1)}$\textbf{-category} と呼ぶ．
\end{marker}

\begin{mydef}[label=def:Alg]{{$\Cat{C}^\otimes$ 上の $\Cat{O}$-代数}}
    $\Cat{O}^\otimes \lto \Ner(\FIN_*),\, \Cat{C}^\otimes \lto \Ner(\FIN_*)$ を\hyperref[def:infty-operad]{$(\infty,\, 1)$-オペラッド}とする．
    \textbf{$\bm{\Cat{C}}$ 上の $\bm{\Cat{O}}$-代数}とは，\hyperref[def:infty-operad-map]{$(\infty,\, 1)$-オペラッドの射}が成す\hyperref[def:fullsub-infty]{充満部分 $(\infty,\, 1)$-圏}
    \begin{align}
        \bm{\Alg{\Cat{O}}{\Cat{C}}} \subset \FUN_{\Ner(\FIN_*)} (\Cat{O}^\otimes ,\, \Cat{C}^\otimes)
    \end{align}
    のこと~\cite[Definition 2.1.2.7.]{Lurie2017HA}．
\end{mydef}

\hyperref[def:Kan]{$\Kan$}-豊穣圏 $\Opinfty$ を次のように構成する：
\begin{itemize}
    \item \hyperref[def:infty-operad]{$(\infty,\, 1)$-オペラッド}を対象とする．
    \item $\Alg{\Cat{O}}{\Cat{C}}$ の\hyperref[def:core]{芯} $\Alg{\Cat{O}}{\Cat{C}}^{\cong} \in \Obj{\Kan}$ をHom対象とする．
\end{itemize}

\begin{mydef}[label=def:OPinfty]{{$(\infty,\, 1)$-オペラッドの成す $(\infty,\, 1)$-圏}}
    \hyperref[def:infinity-1]{$(\infty,\, 1)$-圏} $\OPinfty$ を以下で定義する~\cite[Definition 2.1.4.1.]{Lurie2017HA}：
    \begin{align}
        \OPinfty \coloneqq \hcNer(\Opinfty)
    \end{align}
\end{mydef}

\subsection{{$\Cat{O}$-モノイダル $(\infty,\, 1)$-圏}}

次に，\hyperref[redef:braided-monoidal]{対称モノイダル圏}の $(\infty,\, 1)$-圏論における対応物を考える．

\begin{mydef}[label=def:O-monoidal-infty]{{$\Cat{O}$-モノイダル $(\infty,\, 1)$-圏}}
    \begin{itemize}
        \item \hyperref[def:infinity-1]{$(\infty,\, 1)$-圏} $\Cat{C}^\otimes$
        \item \hyperref[def:infty-operad]{$(\infty,\, 1)$-オペラッド} $r \colon \Cat{O}^\otimes \lto \Ner(\FIN_*)$
    \end{itemize}
    を与える．
    \hyperref[def:infinity-1]{$(\infty,\, 1)$-圏の関手} $p \colon \Cat{C}^\otimes \lto \Cat{O}^\otimes$ が
    \textbf{$\bm{(\infty,\, 1)}$-オペラッドのcoCartesianファイブレーション} (coCartesian fibration of $(\infty,\, 1)$-operads) であるとは，以下の条件を充たすことを言う：
    \begin{description}
        \item[\textbf{(Op-coCart-1)}]　
        
        $p \colon \Cat{C}^\otimes \lto \Cat{O}^\otimes$ は\hyperref[def:Cart-coCart]{coCartesianファイブレーション}である．
        
        \item[\textbf{(Op-coCart-2)}]　
        
        $\Cat{C}^\otimes \xrightarrow{p} \Cat{O}^\otimes \xrightarrow{r} \Ner(\FIN_*)$ は\hyperref[def:infty-operad]{$(\infty,\, 1)$-オペラッド}である．
    \end{description}
    特にこのとき，$p$ は $\Cat{C}^\otimes$ を $\bm{\Cat{O}}$\textbf{-モノイダル $(\infty,\, 1)$-圏}にしている ($p$ exhibits $\Cat{C}^\otimes$ as $\Cat{O}$-monoidal $(\infty,\, 1)$-category) と言う\footnote{なお，この定義から $p$ は\hyperref[def:infty-operad-map]{$(\infty,\, 1)$-オペラッドの射}でもあることが分かる~\cite[Remark 2.1.2.14]{lurie2008higher}．}~\cite[Definition 2.1.2.13.]{Lurie2017HA}．
\end{mydef}

\begin{mydef}[label=def:sym-monoidal-infty]{対称モノイダル {$(\infty,\, 1)$-圏}}
    \hyperref[def:O-monoidal-infty]{$\Cat{O}$-モノイダル $(\infty,\, 1)$-圏} $\Cat{C}^\otimes \xrightarrow{p} \Cat{O}^\otimes$ において $\Cat{O}^\otimes = \Ner(\FIN_*)$ ならば，$\Cat{C}^\otimes$ は\textbf{対称モノイダル $\bm{(\infty,\, 1)}$-圏} (symmetric monoidal $(\infty,\, 1)$-category) であると言う．

    i.e. 対称モノイダル $(\infty,\, 1)$-圏とは，\hyperref[def:infty-operad]{$(\infty,\, 1)$-オペラッド} $\Cat{C}^\otimes \lto \Ner(\FIN_*)$ であって\hyperref[def:Cart-coCart]{coCartesianファイブレーション}でもあるもののこと~\cite[Example 2.1.2.18.]{Lurie2017HA}
\end{mydef}

\begin{mydef}[label=def:O-monoidal-functor]{{$\Cat{O}$-モノイダル $(\infty,\, 1)$-関手}}
    \begin{itemize}
        \item \hyperref[def:infty-operad]{$(\infty,\, 1)$-オペラッド} $\Cat{O}^\otimes \lto \Ner(\FIN_*)$
        \item \hyperref[def:O-monoidal-infty]{$(\infty,\, 1)$-オペラッドのcoCartesianファイブレーション} $p \colon \Cat{C}^\otimes \lto \Cat{O}^\otimes,\; q \colon \Cat{D}^\otimes \lto \Cat{O}^\otimes$
    \end{itemize}
    を与える．$f \in \Alg{\Cat{C}}{\Cat{D}}$ が\textbf{$\bm{\Cat{O}}$-モノイダル $\bm{(\infty,\, 1)}$-関手} ($\Cat{O}$-monoidal $(\infty,\, 1)$-functor) であるとは，
    それが全ての\hyperref[def:p-Cart]{$p$-coCartesianな1-射}を\hyperref[def:p-Cart]{$q$-coCartesianな1-射}へ移すことを言う．

    $\Cat{O}$-モノイダル $(\infty,\, 1)$-関手が成す\hyperref[def:fullsub-infty]{充満部分 $(\infty,\, 1)$-圏}を
    \begin{align}
        \bm{\FUN^\otimes_{\Cat{O}} (\Cat{C},\, \Cat{D})} \subset \FUN_{\Cat{O}^\otimes} (\Cat{C}^\otimes,\, \Cat{D}^\otimes)
    \end{align}
    と書く~\cite[Definition 2.1.3.7.]{Lurie2017HA}．
\end{mydef}

\begin{mydef}[label=def:sym-monoidal-functor]{対称モノイダル {$(\infty,\, 1)$}-関手}
    \hyperref[def:sym-monoidal-infty]{対称モノイダル $(\infty,\, 1)$-圏} $\Cat{C}^\otimes \lto \Ner(\FIN_*)$,\, $\Cat{D}^\otimes \lto \Ner(\FIN_*)$ の間の
    \textbf{対称モノイダル $(\infty,\, 1)$-関手} (symmetric monoidal $(\infty,\, 1)$-functor) とは，$f \in \Alg{\Cat{C}}{\Cat{D}}$ であって全ての\hyperref[def:p-Cart]{coCartesianな1-射}を保存するもののこと．
    対称モノイダル $(\infty,\, 1)$-関手が成す\hyperref[def:fullsub-infty]{充満部分 $(\infty,\, 1)$-圏}を
    \begin{align}
        \bm{\FUN^\otimes (\Cat{C},\, \Cat{D})} \subset \FUN_{\FIN_*} (\Cat{C}^\otimes,\, \Cat{D}^\otimes)
    \end{align}
    と書く~\cite[Definition 2.1.3.7.]{Lurie2017HA}．
\end{mydef}

\hyperref[def:Kan]{$\Kan$}-豊穣圏 $\CATinfty^{\otimes}{}^\Delta$ を次のように構成する：
\begin{itemize}
    \item \hyperref[def:sym-monoidal-infty]{対称モノイダル $(\infty,\, 1)$-圏}を対象とする．
    \item $\FUN^{\otimes}(\Cat{C},\, \Cat{D})$ の\hyperref[def:core]{芯} $\FUN^{\otimes}(\Cat{C},\, \Cat{D})^{\cong} \in \Obj{\Kan}$ をHom対象とする．
\end{itemize}

\begin{mydef}[label=def:Cat-infty-sym-mon]{対称モノイダル {$(\infty,\, 1)$}-圏が成す {$(\infty,\, 1)$-圏}}
    \hyperref[def:infinity-1]{$(\infty,\, 1)$-圏} $\CATinfty^\otimes$ を以下で定義する~\cite[Variant 2.1.4.13.]{Lurie2017HA}：
    \begin{align}
        \CATinfty^\otimes \coloneqq \hcNer(\CATinfty^\otimes{}^\Delta) \subset \OPinfty
    \end{align}
\end{mydef}


定義\ref{def:infty-operad}がオペラッドと呼ぶにふさわしいことを示すために，次の小節では $(1,\, 1)$-圏の文脈で対応物を考えよう．

\subsection{色付きオペラッドと {$(1,\, 1)$}-圏のcoCartesian fibration}

\begin{mydef}[label=def:c-operad,breakable]{colored operad}
    \textbf{色付きオペラッド}\footnote{いわゆる\textbf{対称色付きオペラッド} (symmetric colored operad) である．} (colored operad) $\Cat{O}$ は，以下の4つのデータから成る：
    \begin{itemize}
        \item \textbf{対象}\footnote{\textbf{色} (color) と呼ばれることもある．} (object) の集まり
        \begin{align}
            \bm{\Obj{\Cat{O}}}
        \end{align}
        \item $\forall I \in \Obj{\FIN},\; \forall \Familyset{x_i \in \Obj{\Cat{O}}}{i \in I},\;  \forall y \in \Obj{\Cat{O}}$ の3つ組に対して定まっている，
        $\Familyset{x_i}{i \in I}$ から $y$ へ向かう\textbf{多射} (multimorphism) の集合
        \begin{align}
            \bm{\Mul{\Cat{O}}} \bigl( \Familyset{x_i}{i \in I},\, y \bigr) \in \Obj{\SETS}
        \end{align}
        \item $\forall \alpha \in \Hom{\FIN} (I,\, J),\; \forall \Familyset[\big]{x_i \in \Obj{\Cat{O}}}{i \in I},\; \Familyset[\big]{y_j \in \Obj{\Cat{O}}}{j \in J},\; \forall z \in \Obj{\Cat{O}}$ の4つ組に対して定まっている，
        多射の\textbf{合成} (composition map) と呼ばれる写像
        \begin{align}
            \bm{\circ}_{\bm{\alpha}} \colon \Mul{\Cat{O}} \bigl( \Familyset{y_j}{j \in J},\, z \bigr) \times \prod_{j \in J} \Mul{\Cat{O}} \bigl( \Familyset{x_i}{i \in \alpha^{-1}(\{j\})},\, y_j \bigr) &\lto \Mul{\Cat{O}} \bigl( \Familyset{x_i}{i \in I},\, z \bigr), \\
            \bigl( G,\, (F_j)_{j \in J} \bigr) &\lmto G \circ_{\alpha} (F_j)_{j \in J}
        \end{align}
        \item \textbf{恒等射} (identitiy) と呼ばれる多射の族 $\Familyset[\Big]{\Id_x \in \Mul{\Cat{O}} \bigl(\{x\},\, x \bigr)}{x \in \Obj{\Cat{O}}}$
    \end{itemize}
    これらは以下の条件を充たさねばならない：
    \begin{description}
        \item[\textbf{(cOp-1)}]　
        
        恒等射は合成に関して単位元として振る舞う．
        
        \item[\textbf{(cOp-2)}]　
        
        多射の合成は結合則を充たす．i.e.
        $\forall \alpha \in \Hom{\FIN}(I,\, J),\; \forall \beta \in \Hom{\FIN}(J,\, K),\; \forall \Familyset[\big]{x_i \in \Obj{\Cat{O}}}{i \in I},\; \forall \Familyset[\big]{y_j \in \Obj{\Cat{O}}}{j \in J},\; \forall \Familyset[\big]{z_k \in \Obj{\Cat{O}}}{k \in K},\; \forall w \in \Obj{\Cat{O}}$ に対して，
        $(1,\, 1)$-圏 $\SETS$ の図式
        \small
        \begin{center}
            % https://tikzcd.yichuanshen.de/#N4Igdg9gJgpgziAXAbVABwnAlgFyxMJZARgBoAGAXVJADcBDAGwFcYkQAdDgWWceC4BhejmAB5AL4SuAIywBzRgAouAMXoBbLIwCecGKIBeAfQDWE4Ka5YwAAgDSE0l1K2A7rIUAnAJS2ueBrw-hxoXtDGliE2DhIhvPxCIuJSIXKKKhzqWrr6ojrGAFYWhdF2sgb0AHrAALTEEpmWXBI+Ti62JlYc6b4hgcFcYRHApdZ2AFJxXAkCHMKiktM9Cspqmtp6BsAAHsZYFlhlIUxoABbVdQ1NhS1tzhyuBbcr8r4gTuiYuPiEKABMpGI1DoTFY7BmfDmCxS0leayyG1y2y6Fm6MUcD1cHlefQCWCCcBCwygkVMx0c8ShSUWqU8GXWOS2oj2B2AR3GtkyMkqIQAxlgvHyToxzvQfDV6o0uM0OK12o9OmZ6e9PiAMNg8AQiOQgSCGCw2IhODxqfNkkt6Qjsps8sBniVOVMse4VX58YTieFSaNjlMqYlzbTlulrUjmbt9odjlxThdJdcZS95S7nm6PqQvprfkQyP99WCjSbZjTYVbMjbkSyo+zOQBJBWuQzpiQgmBQeTwIigABm4Q0SF1IBwECQgNBhohHAFQsiFRw9GWAyJXDrUBANEY9B5jAACt8tX8QF4FGccBne-3BzQR0gyBPwcbV1B+gTgsBMjO+XOOHHFz4yQpKQNxALcd33bNtWNRgYB7c81T7CAB0Qcdb0QABmGgDUfE0vx-P8JBAsCYD3A8c2NE95DPC8QEQ5D7zQzCHyLLg8LmHkF35QVhVjUULkI1sJCAA
            \begin{tikzcd}[column sep={10em,between origins},row sep=large]
                                                                                                                                                                                                    & {\Mul{\Cat{O}}\bigl(\Familyset{z_k}{k\in K},\, w\bigr) \times \prod_{k \in K} \Mul{\Cat{O}} \bigl(\Familyset{y_j}{j \in \beta^{-1}(\{k\})},\, z_k\bigr) \times \prod_{j \in J} \Mul{\Cat{O}} \bigl(\Familyset{x_i}{i \in \alpha^{-1}(\{j\})},\, y_j\bigr)} \arrow[ld, "\circ_{\beta} \times \Id"'] \arrow[ddr, "\Id \times {\bigl(\circ_{\alpha|_{\alpha^{-1}(\beta^{-1}(\{k\}))}}\bigr)_{k \in K}}"] &                                                                                                                                                                                                                  \\
            {\Mul{\Cat{O}}\bigl(\Familyset{y_j}{j\in J},\, w\bigr) \times \prod_{j \in J} \Mul{\Cat{O}} \bigl(\Familyset{x_i}{i \in \alpha^{-1}(\{j\})},\, y_j\bigr)} \arrow[ddr, "\circ_{\alpha}"'] &                                                                                                                                                                                                                                                                                                                                                          &                                                                                                                                                                                                                  \\
                                                                                                                                                                                                    &                                                                                                                                                                                                                                                                                                                                                          &{\Mul{\Cat{O}}\bigl(\Familyset{z_k}{k\in K},\, w\bigr) \times \prod_{k \in K} \Mul{\Cat{O}} \bigl(\Familyset{x_i}{i \in (\beta \circ \alpha)^{-1}(\{k\})},\, z_k\bigr)} \arrow[ld, "\circ_{\beta \circ \alpha}"]  \\
                                                                                                                                                                                                    & {\Mul{\Cat{O}}\bigl(\Familyset{x_i}{i\in I},\, z\bigr)}                                                                                                                                                                                                                                                                                                  &                                                                                                                                                                                                                 
            \end{tikzcd}
        \end{center}
        \normalsize
        は可換である．
    \end{description}
\end{mydef}

\hyperref[def:c-operad]{色付きオペラッド}の定義において，写像 $\alpha \in \Hom{\FIN} (I,\, J)$ が多射の合成の「型」を規定している．


\begin{mydef}[label=def:p-coCart-ord,breakable]{{$(1,\, 1)$}-圏におけるcoCartesian fibration}
    $p \colon \Cat{E} \lto \Cat{B}$ を $(1,\, 1)$-圏の関手とする．
    \begin{itemize}
        \item $\Cat{E}$ の射 $\bar{f} \in \Hom{\Cat{E}} (\bar{x},\, \bar{y})$ が以下の条件を充たすとき，$\bar{f}$ は\textbf{$\bm{p}$-coCartesian}であると言う：
        \begin{description}
            \item[\textbf{(coCart-ord)}] 
            $(1,\, 1)$-圏 $\Cat{B}$ の図式
            \begin{center}
                % https://tikzcd.yichuanshen.de/#N4Igdg9gJgpgziAXAbVABwnAlgFyxMJZABgBoBGAXVJADcBDAGwFcYkQ0AKAHW4CN6AJ2AAPAL4BKEGNLpMufIRTlSxanSat2XXgOEBPSdNkd5eAkQBMFdQxZtEHHvyHAAXkbHqYUAObwiUAAzQQgAWyQyEBwIJBUNe21nPWAgzxMQ8MiaGKRrBK1HHW4giEEmRgACXgZBNAALLAB9YikaRno+GEYABTNFdkEsX3qcY2DQiMR43MR8u0KQXlLyxiqaoQasaUoxIA
                \begin{tikzcd}
                                                                                        & p(\bar{y}) \arrow[rd, blue, "\forall \varphi"] &            \\
                p(\bar{x}) \arrow[ru, "p(\bar{f})"] \arrow[rr, blue, "p(\forall \varphi_0)"'] &                                          & p(\forall \textcolor{blue}{\bar{z}})
                \end{tikzcd}
            \end{center}
            を可換にする勝手な2つの射 $\textcolor{blue}{\varphi_0} \in \Hom{\Cat{E}} (\bar{x},\, \textcolor{blue}{\bar{z}}),\; \textcolor{blue}{\varphi} \in \Hom{\Cat{B}} \bigl( p(\bar{y}),\, p(\textcolor{blue}{\bar{z}}) \bigr)$ に対して，
            射 $\textcolor{red}{\bar{\varphi}} \in \Hom{\Cat{E}} (\bar{y},\, \textcolor{blue}{\bar{z}})$ が\underline{一意的に}存在して $(1,\, 1)$-圏 $\Cat{E}$ の図式
            \begin{center}
                % https://tikzcd.yichuanshen.de/#N4Igdg9gJgpgziAXAbVABwnAlgFyxMJZABgBoBGAXVJADcBDAGwFcYkQAdDgI3oCdgADwC+IYaXSZc+QinKli1Ok1bsuvAQE9R4ydjwEiAJgpKGLNok49+wAF46lMKAHN4RUADM+EALZIyEBwIJHllCzUbAU8dCRBvPwCaYKQTcNUrLgY+NAALLAB9YhAaRnpuGEYABSkDWRA+LBdcnDE4hP9EMJTENPMM6w1gLP48rEdhIA
                \begin{tikzcd}
                                                                       & \bar{y} \arrow[rd, red,"\exists ! \bar{\varphi}"] &         \\
                \bar{x} \arrow[ru, "\bar{f}"] \arrow[rr, blue, "\varphi_0"'] &                                     & \textcolor{blue}{\bar{z}}
                \end{tikzcd}
            \end{center}
            を可換にする．
        \end{description}
        \item $p$ が\textbf{coCartesian fibration}であるとは，
        \begin{itemize}
            \item $(1,\, 1)$-圏 $\Cat{B}$ の任意の射 $f \in \Hom{\Cat{B}} (x,\, y)$
            \item $p(\bar{x}) = x$ を充たす $(1,\, 1)$-圏 $\Cat{E}$ の対象 $\bar{x} \in \Obj{\Cat{E}}$
        \end{itemize}
        に対して，以下の条件を充たす $(1,\, 1)$-圏 $\Cat{E}$ の射 $\bar{f} \in \Hom{\Cat{E}} (\bar{x},\, z)$ が存在することを言う：
        \begin{description}
            \item[\textbf{(coCart-ord-1)}]　
            
            $\bar{f}$ は $f$ の持ち上げである．i.e. $p(\bar{f}) = f$ が成り立つ．

            \item[\textbf{(coCart-ord-2)}]　
            
            $\bar{f}$ は $p$-coCartesianである．
        \end{description}
        % \item $p$ が\textbf{isofibration}であるとは，
        % \begin{itemize}
        %     \item $(1,\, 1)$-圏 $\Cat{B}$ の任意の\underline{同型射} $f \in \Hom{\Cat{B}} (x,\, y)$
        %     \item $p(\bar{x}) = x$ を充たす $(1,\, 1)$-圏 $\Cat{E}$ の対象 $\bar{x} \in \Obj{\Cat{E}}$
        % \end{itemize}
        % に対して，条件 \textsf{\textbf{(coCart-ord-1)}}, \textsf{\textbf{(cocCart-ord-2)}} を充たす $(1,\, 1)$-圏 $\Cat{E}$ の射 $\bar{f} \in \Hom{\Cat{E}} (\bar{x},\, z)$ が存在することを言う．
    \end{itemize}
\end{mydef}

定義\ref{def:infty-operad}に合わせて，$(1,\, 1)$-圏の関手
\begin{align}
    \label{eq:operad-ord}
    p \colon \Cat{O}^\otimes \lto \FIN_*
\end{align}
であって以下の3条件を充たすものを考えてみる：
\begin{description}
    \item[\textbf{(OP-ord-1)}]　
    \label{cond:Op-ord-1}

    任意の\hyperref[def:inert-active]{inertな}射 $\alpha \in \Hom{\FIN_*} \bigl( \expval{m},\, \expval{n} \bigr)$ および $\forall x \in \Obj{\Cat{O^\otimes}_{\expval{m}}}$ に対して，
    $\Cat{O}^\otimes$ における\hyperref[def:p-coCart-ord]{$p$-coCartesian}な射 $\bar{\alpha} \in \Hom{\Cat{O}^\otimes} (x,\, y)$ が存在して $p(\bar{\alpha}) = \alpha$ を充たす．
        
    \item[\textbf{(OP-ord-2)}]　 
    \label{cond:Op-ord-2}
            
    $\forall \alpha \in \Hom{\FIN_*} \bigl( \expval{m},\, \expval{n} \bigr)$ および $\forall x \in \Obj{\Cat{O}^\otimes_{\expval{m}}},\, \forall y \in \Obj{\Cat{O}^\otimes_{\expval{n}}}$ に対して，
    inertな射 $\rho^i \in  \Hom{\FIN_*}\bigl( \expval{n},\, \expval{1} \bigr)$ に \hyperref[cond:Op-ord-1]{\textsf{\textbf{(Op-ord-1)}}} を適用して得られる\hyperref[def:p-coCart-ord]{$p$-coCartesian}な射の族 $\Familyset[\big]{\bar{\rho}^i_y \in \Hom{\Cat{O}^\otimes} (y,\, y_i)}{1 \le i \le n}$ が誘導する写像
    \begin{align}
        \Hom{\Cat{O}^\otimes} (x,\, y)_\alpha &\lto \prod_{i=1}^n \Hom{\Cat{O}^\otimes} (x,\, y_i)_{\rho^i \circ \alpha}, \\
        \varphi &\lmto (\bar{\rho}^1 \circ \varphi,\, \dots,\, \bar{\rho}^n \circ \varphi)
    \end{align}
    は，$(1,\, 1)$-圏 $\SETS$ における\hyperref[def:iso]{同型射}（i.e. 全単射）である．

    \item[\textbf{(OP-ord-3)}]　 
    \label{cond:Op-ord-3}

    $\forall x_1,\, \dots,\, x_n \in \Obj{\Cat{O}^\otimes_{\expval{1}}}$ に対して，
    ある $x \in \Obj{\Cat{O}^\otimes_{\expval{n}}}_0$ および \hyperref[def:p-Cart]{$p$-coCartesian}な射 $\widehat{\rho_i} \in \Hom{\Cat{O}^\otimes} (x,\, x_i)_{\rho^i}$ が存在する．
\end{description}
ここで，$\CAT$ における\hyperref[def:pullback-pushout]{引き戻し}
\begin{center}
    % https://tikzcd.yichuanshen.de/#N4Igdg9gJgpgziAXAbVABwnAlgFyxMJZABgBpiBdUkANwEMAbAVxiRAB12BhOnYAeQC+APWCcIeALbxBAfTHsYADzT0GwMIMEhBpdJlz5CKMgEYqtRizacefIcPFT4OvSAzY8BIqdLnqataIHOySvAAWAE6SwAByggAUnABiAJKxsgBUAJSu+p5GPuQWgazBnAAiMAw4dMLEOhYwUADmLiigAGaREJJIZCA4EEgAzLpdPX2IA0NIpuMg3b2j1LOIAEyrdFgMbOEQEADWIAFWZSHKqowa2gtLU76DwxunzOdoJyAMdABG1QAKBi8xhAkSwLXCOEagiAA
    \begin{tikzcd}
    \textcolor{red}{\Cat{O}^{\otimes}_{\expval{n}}} \arrow[red, r] \arrow[red,d] &* \arrow[d, "\expval{n}", hook] \\
    \Cat{O}^\otimes \arrow[r, "p"']                    & \FIN_*
    \end{tikzcd}
\end{center}
により $(1,\, 1)$-圏 $\Cat{O}^\otimes_{\expval{n}}$ を定義している．具体的には
\begin{align}
    \Obj{\Cat{O}^\otimes_{\expval{n}}} &= \bigl\{\, x \in \Obj{\Cat{O}^\otimes} \bigm| p(x) = \expval{n} \,\bigr\}, \\
    \Hom{\Cat{O}^\otimes_{\expval{n}}}(x,\, y) &= \bigl\{\, f \in \Hom{\Cat{O}^\otimes} \bigm| p(f) = \Id_{\expval{n}} \,\bigr\} 
\end{align}
である．さらに，$\forall \alpha \in \Hom{\FIN_*} \bigl( \expval{m},\, \expval{n} \bigr)$ に対して $\SETS$ の引き戻しにより
\begin{center}
    % https://tikzcd.yichuanshen.de/#N4Igdg9gJgpgziAXAbVABwnAlgFyxMJZABgBpiBdUkANwEMAbAVxiRAB12swAzHATwCydNME4BhOjmAB5AL4A9ThDwBbeHIAEACgAepTqU38AlAH1OjNAAs6IOaXSZc+QijIBGKrUYs2nbj4hETF2SWl5JXYVLHU4OT0DdiNTe0cQDGw8AiIPcm96ZlZEDnYAERgGHDoFYjSnLNdc0i9qQr8SgN4BYVFOADkYACdtTgAxbjMAKhM5TgAjLABzBm00M2ASCgTdEySjdc3KBNMF5aGTe28YKCX4IlAeIYhVJAAmahwIJABmT7osAw2NYIBAANYgNq+YqlKy2eogJ4vJB5EBfX5QopsNCQkAMOjzSoABWc2TcICGy2sOARSNeiDIaO+iA8Dkez3pjPRiDecm2QA
    \begin{tikzcd}
        {\textcolor{red}{\Hom{\Cat{O}^\otimes} (x,\, y)_\alpha}} \arrow[red,d] \arrow[red,r] & * \arrow[d, "\alpha", hook]                            \\
        {\Hom{\Cat{O}^\otimes}(x,\, y)} \arrow[r, "p"']             & {\Hom{\FIN_*}\bigl(\expval{m},\, \expval{n}\bigr)}
    \end{tikzcd}
\end{center}
と定義した．

\begin{mydef}[label=def:cat-of-operators]{Category of operators}
    いま，\hyperref[def:c-operad]{colored operad} $\Cat{O}$ が与えられたとする．このとき\textbf{category of operators}と呼ばれる $(1,\, 1)$-圏 $\Cat{O}^{\otimes}$ を次のように定義する：
    \begin{itemize}
        \item 2つ組 $\bigl( \expval{n};\, x_1,\, \dots,\, x_n\bigr) \in \Obj{\FIN_*} \times \Obj{\Cat{O}}$ を対象に持つ．
        \item $\forall \bigl( \expval{m};\, x_1,\, \dots,\, x_m\bigr),\, \bigl( \expval{n};\, y_1,\, \dots,\, y_n \bigr) \in \Obj{\Cat{O}^\otimes}$ に対して，Hom集合を
        \begin{align}
            &\Hom{\Cat{O}^\otimes} \Bigl( \bigl( \expval{m};\, x_1,\, \dots,\, x_m \bigr)  ,\, \bigl( \expval{n};\, y_1,\, \dots,\, y_n \bigr)  \Bigr) \\
            &\coloneqq \coprod_{\alpha \in \Hom{\FIN_*}\bigl( \expval{m},\, \expval{n} \bigr)} \prod_{j=1}^n \Mul{\Cat{O}} \bigl( \Familyset{x_i}{i \in \alpha^{-1}(\{j\})},\, y_j \bigr)
        \end{align}
        と定義する．
        \item 射の合成は，$\FIN_*$ における射の合成および $\Cat{O}$ における多射の合成によって定める．
    \end{itemize}
    \tcblower
    $\Cat{O}^\otimes$ から $\FIN_*$ への忘却関手を
    \begin{align}
        p \colon \Cat{O}^{\otimes} &\lto \FIN_*, \\
        \bigl( \expval{n};\, x_1,\, \dots,\, x_n\bigr) &\lmto \expval{n}, \\
        \bigl( \alpha;\, \phi_1,\, \dots,\, \phi_n \bigr) &\lmto \alpha
    \end{align}
    と定義する．
\end{mydef}

\begin{myprop}[label=prop:c-operad]{色付きオペラッドとcategory of operators}
    $(1,\, 1)$-圏の関手
    \begin{align}
        p \colon \Cat{O}^\otimes \lto \FIN_*
    \end{align}
    において，$\Cat{O}^\otimes$ がある\hyperref[def:c-operad]{色付きオペラッド} $\Cat{O}$ の\hyperref[def:cat-of-operators]{category of operators}と圏同値になる必要十分条件は，
    $p$ が条件\hyperref[cond:Op-ord-1]{\textsf{\textbf{(Op-ord-1)}}}-\hyperref[cond:Op-ord-1]{\textsf{\textbf{(Op-ord-3)}}}を充たすことである．
\end{myprop}

\begin{proof}
    ~\cite[Proposition 2.2.II.]{Haugseng2023iop}
\end{proof}

逆の対応を作ることもできる．

\begin{myprop}[label=prop:operad-ord,breakable]{色付きオペラッドの再構成}
    条件\hyperref[cond:Op-ord-1]{\textsf{\textbf{(Op-ord-1)}}}-\hyperref[cond:Op-ord-1]{\textsf{\textbf{(Op-ord-3)}}}を充たす $(1,\, 1)$-圏の関手\eqref{eq:operad-ord}から次のように構成されたデータの組み $\Cat{O}$ は\hyperref[def:c-operad]{色付きオペラッド}を成す：
    \begin{itemize}
        \item 対象の集まりを $\Obj{\Cat{O}} \coloneqq \Obj{\Cat{O}^\otimes_{\expval{1}}}$ と定義する．
        \item $\forall x_1,\, \dots,\, x_n,\, y \in \Obj{\Cat{O}}$ に対して，以下の3つ組全体が成す集合を $\Mul{\Cat{O}} \bigl( (x_1,\, \dots,\, x_n),\, y \bigr)$ と定義する．
        \begin{enumerate}
            \item $x_1,\, \dots,\, x_n$ に対して\hyperref[cond:Op-ord-1]{\textsf{\textbf{(Op-ord-3)}}}を適用することにより定まる $X \in \Obj{\Cat{O}^\otimes_{\expval{n}}}$ 
            \item $x_1,\, \dots,\, x_n$ に対して\hyperref[cond:Op-ord-1]{\textsf{\textbf{(Op-ord-3)}}}を適用することにより定まる\hyperref[def:p-Cart]{$p$-coCartesian}な射の族 $\Familyset[\big]{\widehat{\rho_i} \in \Hom{\Cat{O}^\otimes}(X,\, x_i)_{\rho^i}}{1 \le i \le n}$
            \item \exref{def:alpha-n}の\hyperref[def:inert-active]{active}な射
            $\alpha_n \in \Hom{\FIN_*} \bigl( \expval{n},\, \expval{1} \bigr)$ およびその上の $\Cat{O}^\otimes$ の射 $F \in \Hom{\Cat{O}^\otimes} (X,\, y)_{\alpha_n}$
        \end{enumerate}
        \item $n$ 個の多射
        \begin{align}
            (X^1,\, \Familyset{\widehat{\rho_i}^1}{1 \le i \le m_1},\, F^1) &\in \Mul{\Cat{O}} \bigl( (x^1_1,\, \dots,\, x^1_{m_1}),\, y_1 \bigr),\\ 
            &\vdots \\ 
            (X^n,\, \Familyset{\widehat{\rho_i}^n}{1 \le i \le m_n},\, F^n) &\in \Mul{\Cat{O}} \bigl( (x^n_1,\, \dots,\, x^n_{m_n}),\, y_n \bigr)
        \end{align}
        と1つの多射 $(Y,\, \Familyset{\widehat{\rho_j}}{1 \le j \le n},\, G) \in \Mul{\Cat{O}} \bigl( (y_1,\, \dots,\, y_n),\, z\bigr)$ の合成
        \begin{align}
            \bigl(X,\, \Familyset{\widehat{\rho_i}}{1 \le i \le m_1+\cdots +m_n},\, G \circ (F^1;\, \dots \, ;\, F^n)\bigr) \in \Mul{\Cat{O}} \bigl( (x^1_1,\, \dots,\, x^1_{m_1};\, \dots\, ;\, x^n_1,\, \dots,\, x^n_{m_n}),\, z \bigr) 
        \end{align}
        を次のように定義する：
        \begin{enumerate}
            \item $X \in \Obj{\Cat{O}^\otimes_{\expval{m_1 + \cdots + m_n}}}$ は，$x^1_1,\, \dots,\, x^n_{m_n} \in \Obj{\Cat{O}}$ に対して\hyperref[cond:Op-ord-1]{\textsf{\textbf{(Op-ord-3)}}}を適用することにより定める．
            \item \hyperref[def:p-Cart]{$p$-coCartesian}な射の族 $\Familyset[\big]{\widehat{\rho^j_{i}} \in \Hom{\Cat{O}^\otimes}(X,\, x^j_i)_{\rho^{i+m_1+\cdots+m_{j-1}}}}{\substack{1 \le j \le n \\ 1 \le i \le m_j}}$ は $x^1_1,\, \dots,\, x^n_{m_n} \in \Obj{\Cat{O}}$ に対して\hyperref[cond:Op-ord-1]{\textsf{\textbf{(Op-ord-3)}}}を適用することにより定める．
            \item $\Cat{O}^\otimes$ の射 $G \circ (F^1;\, \dots \, ;\, F^n) \in \Hom{\Cat{O}^\otimes} (X,\, z)_{\alpha_{m_1 + \cdots + m_n}}$ は以下の手順に従って構成する：
            \begin{description}
                \item[\textbf{(STEP-1)}]　
                
                    まず，$1 \le \forall j \le n$ に対して\hyperref[def:inert-active]{inertな射} $\pi_j \in \Hom{\FIN_*} \bigl( \expval{m_1 + \cdots + m_n},\, \expval{m_j} \bigr)$ を
                    \begin{align}
                        \pi_j \colon \expval{m_1 + \cdots + m_n} &\lto \expval{m_j}, \\
                        k &\lmto \begin{cases}
                            k - (m_1 + \cdots + m_{j-1}), &1 \le k - (m_1 + \cdots + m_{j-1}) \le m_j \\
                            *, &\text{otherwise}
                        \end{cases}
                    \end{align}
                    で定義する．
                \item[\textbf{(STEP-2)}]　
                
                    \hyperref[def:inert-active]{inertな射} $\pi_j \in \Hom{\FIN_*} \bigl( \expval{m_1 + \cdots + m_n},\, \expval{m_j} \bigr)$ に対して\hyperref[cond:Op-ord-1]{\textsf{\textbf{(Op-ord-1)}}}を適用することにより，\hyperref[def:p-coCart-ord]{$p$-coCartesianな射} 
                    $\bar{\pi}_j \in \Hom{\Cat{O}^\otimes} (X,\, X^j)_{\pi_j}$ を取得する．

                \item[\textbf{(STEP-3)}]　
                
                    \begin{align}
                        (F^1 \circ \bar{\pi}_1,\, \dots,\, F^n \circ \bar{\pi}_n) \in \prod_{j=1}^n \Hom{\Cat{O}^\otimes} (X,\, y_j)_{\alpha_{m_j} \circ \pi_j}
                    \end{align}
                    に対して\hyperref[cond:Op-ord-1]{\textsf{\textbf{(Op-ord-2)}}}を適用することで，対応する
                    \begin{align}
                        F \in \Hom{\Cat{O}^\otimes} (X,\, Y)_\pi
                    \end{align}
                    が一意的に定まる．これに $G \in \Hom{\Cat{O}^\otimes}(Y,\, z)_{\alpha_n}$ を合成して
                    \begin{align}
                        G \circ_{\pi} (F^1;\, \dots \, ;\, F^n) \coloneqq G \circ F \in \Hom{\Cat{O}^\otimes} (X,\, z)_{\alpha_{m_1 + \cdots + m_n}}
                    \end{align}
                    と定義する．
                    ただし，$\pi \in \Hom{\FIN_*} \bigl( \expval{m_1 + \cdots + m_n},\, \expval{n} \bigr)$ は
                    \begin{align}
                        \pi(k) \coloneqq 
                        \begin{cases}
                            j, &m_{j-1} < k \le m_j \\
                            *, &k=*
                        \end{cases}
                    \end{align}
                    と定義される\hyperref[def:inert-active]{activeな射}である\footnote{$1 \le \forall j \le n$ に対して $\rho^j \circ \pi = \alpha_{m_j} \circ \pi_j$ が成り立つ．}．
            \end{description}
        \end{enumerate}
    \end{itemize}
\end{myprop}

\begin{proof}
    
\end{proof}

\subsection{{$(\infty,\, 1)$}-オペラッドの構成}

% 話を\hyperref[def:infty-operad]{$(\infty,\, 1)$-オペラッド}に戻そう．
定理\ref{thm:Kan-enriched}のおかげで，\hyperref[def:nerve-hc]{homotopy coherent nerve}
\begin{align}
    \hcNer \colon \sCat \lto \sSet
\end{align}
を用いて $\Kan$-豊穣圏から非自明な\hyperref[def:infinity-1]{$(\infty,\, 1)$-圏}を作ることができるのだった．
同様に，$\hcNer$ を用いて非自明な\hyperref[def:infty-operad]{$(\infty,\, 1)$-オペラッド}を構成することができる．

\begin{mydef}[label=def:sc-operad,breakable]{simplicial colored operad}
    \textbf{単体的色付きオペラッド} (simplicial colored operad) $\Cat{O}$ は，以下の4つのデータから成る：
    \begin{itemize}
        \item \textbf{対象}\footnote{\textbf{色} (color) と呼ばれることもある．} (object) の集まり
        \begin{align}
            \bm{\Obj{\Cat{O}}}
        \end{align}
        \item $\forall I \in \Obj{\FIN},\; \forall \Familyset{x_i \in \Obj{\Cat{O}}}{i \in I},\;  \forall y \in \Obj{\Cat{O}}$ の3つ組に対して定まっている，
        $\Familyset{x_i}{i \in I}$ から $y$ へ向かう\textbf{多射} (multimorphism) が成す\hyperref[def:SimpSet]{単体的集合}
        \begin{align}
            \bm{\Mul{\Cat{O}}} \bigl( \Familyset{x_i}{i \in I},\, y \bigr) \in \Obj{\sSet}
        \end{align}
        \item $\forall \alpha \in \Hom{\FIN} (I,\, J),\; \forall \Familyset[\big]{x_i \in \Obj{\Cat{O}}}{i \in I},\; \Familyset[\big]{y_j \in \Obj{\Cat{O}}}{j \in J},\; \forall z \in \Obj{\Cat{O}}$ の4つ組に対して定まっている，
        多射の\textbf{合成} (composition map) と呼ばれる\hyperref[def:nat]{自然変換}
        \begin{align}
            \bm{\circ}_{\bm{\alpha}} \colon \Mul{\Cat{O}} \bigl( \Familyset{y_j}{j \in J},\, z \bigr) \times \prod_{j \in J} \Mul{\Cat{O}} \bigl( \Familyset{x_i}{i \in \alpha^{-1}(\{j\})},\, y_j \bigr) &\lto \Mul{\Cat{O}} \bigl( \Familyset{x_i}{i \in I},\, z \bigr), \\
            \bigl( G,\, (F_j)_{j \in J} \bigr) &\lmto G \circ_{\alpha} (F_j)_{j \in J}
        \end{align}
        \item \textbf{恒等射} (identitiy) と呼ばれる多射の族 $\Familyset[\Big]{\Id_x \in \Mul{\Cat{O}} \bigl(\{x\},\, x \bigr)}{x \in \Obj{\Cat{O}}}$
    \end{itemize}
    これらは以下の条件を充たさねばならない：
    \begin{description}
        \item[\textbf{(cOp-1)}]　
        
        恒等射は合成に関して単位元として振る舞う．
        
        \item[\textbf{(cOp-2)}]　
        
        多射の合成は結合則を充たす．i.e.
        $\forall \alpha \in \Hom{\FIN}(I,\, J),\; \forall \beta \in \Hom{\FIN}(J,\, K),\; \forall \Familyset[\big]{x_i \in \Obj{\Cat{O}}}{i \in I},\; \forall \Familyset[\big]{y_j \in \Obj{\Cat{O}}}{j \in J},\; \forall \Familyset[\big]{z_k \in \Obj{\Cat{O}}}{k \in K},\; \forall w \in \Obj{\Cat{O}}$ に対して，
        $(1,\, 1)$-圏 $\sSet$ の図式
        \small
        \begin{center}
            % https://tikzcd.yichuanshen.de/#N4Igdg9gJgpgziAXAbVABwnAlgFyxMJZARgBoAGAXVJADcBDAGwFcYkQAdDgWWceC4BhejmAB5AL4SuAIywBzRgAouAMXoBbLIwCecGKIBeAfQDWE4Ka5YwAAgDSE0l1K2A7rIUAnAJS2ueBrw-hxoXtDGliE2DhIhvPxCIuJSIXKKKhzqWrr6ojrGAFYWhdF2sgb0AHrAALTEEpmWXBI+Ti62JlYc6b4hgcFcYRHApdZ2AFJxXAkCHMKiktM9Cspqmtp6BsAAHsZYFlhlIUxoABbVdQ1NhS1tzhyuBbcr8r4gTuiYuPiEKABMpGI1DoTFY7BmfDmCxS0leayyG1y2y6Fm6MUcD1cHlefQCWCCcBCwygkVMx0c8ShSUWqU8GXWOS2oj2B2AR3GtkyMkqIQAxlgvHyToxzvQfDV6o0uM0OK12o9OmZ6e9PiAMNg8AQiOQgSCGCw2IhODxqfNkkt6Qjsps8sBniVOVMse4VX58YTieFSaNjlMqYlzbTlulrUjmbt9odjlxThdJdcZS95S7nm6PqQvprfkQyP99WCjSbZjTYVbMjbkSyo+zOQBJBWuQzpiQgmBQeTwIigABm4Q0SF1IBwECQgNBhohHAFQsiFRw9GWAyJXDrUBANEY9B5jAACt8tX8QF4FGccBne-3BzQR0gyBPwcbV1B+gTgsBMjO+XOOHHFz4yQpKQNxALcd33bNtWNRgYB7c81T7CAB0Qcdb0QABmGgDUfE0vx-P8JBAsCYD3A8c2NE95DPC8QEQ5D7zQzCHyLLg8LmHkF35QVhVjUULkI1sJCAA
            \begin{tikzcd}[column sep={10em,between origins},row sep=large]
                                                                                                                                                                                                    & {\Mul{\Cat{O}}\bigl(\Familyset{z_k}{k\in K},\, w\bigr) \times \prod_{k \in K} \Mul{\Cat{O}} \bigl(\Familyset{y_j}{j \in \beta^{-1}(\{k\})},\, z_k\bigr) \times \prod_{j \in J} \Mul{\Cat{O}} \bigl(\Familyset{x_i}{i \in \alpha^{-1}(\{j\})},\, y_j\bigr)} \arrow[ld, "\circ_{\beta} \times \Id"'] \arrow[ddr, "\Id \times {\bigl(\circ_{\alpha|_{\alpha^{-1}(\beta^{-1}(\{k\}))}}\bigr)_{k \in K}}"] &                                                                                                                                                                                                                  \\
            {\Mul{\Cat{O}}\bigl(\Familyset{y_j}{j\in J},\, w\bigr) \times \prod_{j \in J} \Mul{\Cat{O}} \bigl(\Familyset{x_i}{i \in \alpha^{-1}(\{j\})},\, y_j\bigr)} \arrow[ddr, "\circ_{\alpha}"'] &                                                                                                                                                                                                                                                                                                                                                          &                                                                                                                                                                                                                  \\
                                                                                                                                                                                                    &                                                                                                                                                                                                                                                                                                                                                          &{\Mul{\Cat{O}}\bigl(\Familyset{z_k}{k\in K},\, w\bigr) \times \prod_{k \in K} \Mul{\Cat{O}} \bigl(\Familyset{x_i}{i \in (\beta \circ \alpha)^{-1}(\{k\})},\, z_k\bigr)} \arrow[ld, "\circ_{\beta \circ \alpha}"]  \\
                                                                                                                                                                                                    & {\Mul{\Cat{O}}\bigl(\Familyset{x_i}{i\in I},\, z\bigr)}                                                                                                                                                                                                                                                                                                  &                                                                                                                                                                                                                 
            \end{tikzcd}
        \end{center}
        \normalsize
        は可換である．
    \end{description}
\end{mydef}

\begin{mydef}[label=def:sc-op-fibrant]{ファイブラントなsimplicial colored operad}
    \hyperref[def:sc-operad]{Simplicial colored operad} $\Cat{O}$ が\textbf{ファイブラント} (fibrant) であるとは，
    $\forall I \in \Obj{\FIN},\; \forall \Familyset{x_i \in \Obj{\Cat{O}}}{i \in I},\;  \forall y \in \Obj{\Cat{O}}$ に対して，
    $\Mul{\Cat{O}} \bigl( \Familyset{x_i}{i \in I},\, y \bigr) \in \Obj{\Kan}$ であることを言う．
\end{mydef}

次の構成は定義\ref{def:cat-of-operators}とほぼ同様である．

\begin{mydef}[label=def:scat-of-operators,breakable]{Simplicial category of operators}
    いま，\hyperref[def:sc-operad]{simplicial colored operad} $\Cat{O}$ が与えられたとする．このとき\textbf{category of operators}と呼ばれる\hyperref[def:SimpCat]{単体的豊穣圏} $\Cat{O}^{\otimes}$ を次のように定義する：
    \begin{itemize}
        \item 2つ組 $\bigl( \expval{n};\, x_1,\, \dots,\, x_n\bigr) \in \Obj{\FIN_*} \times \Obj{\Cat{O}}$ を対象に持つ．
        \item $\forall \bigl( \expval{m};\, x_1,\, \dots,\, x_m\bigr),\, \bigl( \expval{n};\, y_1,\, \dots,\, y_n \bigr) \in \Obj{\Cat{O}^\otimes}$ に対して，\hyperref[def:enriched]{Hom対象}を
        \begin{align}
            &\Hom{\Cat{O}^\otimes} \Bigl( \bigl( \expval{m};\, x_1,\, \dots,\, x_m \bigr)  ,\, \bigl( \expval{n};\, y_1,\, \dots,\, y_n \bigr)  \Bigr) \\
            &\coloneqq \coprod_{\alpha \in \Hom{\FIN_*}\bigl( \expval{m},\, \expval{n} \bigr)} \prod_{j=1}^n \Mul{\Cat{O}} \bigl( \Familyset{x_i}{i \in \alpha^{-1}(\{j\})},\, y_j \bigr) \in \Obj{\sSet}
        \end{align}
        と定義する．
        \item 射の合成は，$\FIN_*$ における射の合成および $\Cat{O}$ における多射の合成によって定める．
    \end{itemize}
    \tcblower
    $\Cat{O}^\otimes$ から $\FIN_*$ への忘却関手を
    \begin{align}
        p \colon \Cat{O}^{\otimes} &\lto \FIN_*, \\
        \bigl( \expval{n};\, x_1,\, \dots,\, x_n\bigr) &\lmto \expval{n}, \\
        \bigl( \alpha;\, \phi_1,\, \dots,\, \phi_n \bigr) &\lmto \alpha
    \end{align}
    と定義する．
\end{mydef}

\begin{myprop}[label=prop:operadic-nerve]{Operadic nerve}
    \hyperref[def:sc-op-fibrant]{ファイブラント}な\hyperref[def:sc-operad]{simplicial colored operad} $\Cat{O}$ に対して，
    その\hyperref[def:scat-of-operators]{simplicial category of operators} $\Cat{O}^\otimes \lto \FIN_*$ に\hyperref[def:nerve-hc]{homotopy coherent nerve}を作用させて得られる\hyperref[def:infinity-1]{$(\infty,\, 1)$-圏の関手}
    \begin{align}
        p \colon \hcNer (\Cat{O}^\otimes) \lto \Ner(\FIN_*)
    \end{align}
    は\hyperref[def:infty-operad]{$(\infty,\, 1)$-オペラッド}である．
\end{myprop}

\begin{proof}
    
\end{proof}

\begin{myprop}[label=prop:Op-infty-limit]{}
    忘却関手
    \begin{align}
        \text{forget} \colon \OPinfty \lto \CATinfty{}_{/\Ner(\FIN_*)}
    \end{align}
    は，\hyperref[def:lim-colim-infty]{極限}を生成する (creates limits)．
    i.e. 任意の小さな \hyperref[def:infinity-1]{$(\infty,\, 1)$-圏の図式} $D \colon \Cat{K} \lto \OPinfty$ に対して，$\Cat{K} \xrightarrow{D} \OPinfty \xrightarrow{\text{forget}} \CATinfty{}_{/\Ner(\FIN_*)}$ が\hyperref[def:lim-colim-infty]{極限} $(p \colon \Cat{O} \lto \Ner(\FIN_*)) \in (\CATinfty{}_{/\Ner(\FIN_*)})_0$ を持つ
    \footnote{正確に言うと，極限 $\lim_{\Cat{K}} (\mathrm{forget} \circ D) \in \Hom{\mathrm{forget} \circ D} (\leftCone{\Cat{K}},\, \CATinfty{}_{/\Ner(\FIN_*)})$ のコーンポイントのことを $p \colon \Cat{O} \lto \Ner(\FIN_*)$ と書いている．}
    ならば，
    $p$ は\hyperref[def:infty-operad]{$(\infty,\, 1)$-オペラッド}の構造を持ち，かつそれが $D$ の極限と一致する．
    ~\cite[Lemma 1.12.]{AFT2014FH}
 \end{myprop}

\begin{proof}
    忘却関手を $F \colon \OPinfty \lto \CATinfty{}_{/\Ner(\FIN_*)}$ と略記する．
    小さな \hyperref[def:infinity-1]{$(\infty,\, 1)$-圏の図式} $D \colon \Cat{K} \lto \OPinfty$ であって，
    $\Cat{K} \xrightarrow{D} \OPinfty \xrightarrow{F} \CATinfty{}_{/\Ner(\FIN_*)}$ が\hyperref[def:lim-colim-infty]{極限} $\lim_{\Cat{K}} (F \circ D) \in \Hom{F \circ D} (\leftCone{\Cat{K}},\, \CATinfty{}_{/\Ner(\FIN_*)})$ を持つものを1つ固定する．
    
    $\forall k \in \Cat{K}_0$ に対して $D_{[0]}(k) \eqqcolon \bigl(p_k \colon \Oper{\Cat{O}}_k \lto \Ner(\FIN_*)\bigr) \in (\OPinfty)_0$ とおく．
    さらに，
    $\leftCone{\Cat{K}} \xrightarrow{\lim_{\Cat{K}} (F \circ D)} \CATinfty{}_{/\Ner(\FIN_*)}$
    のコーンポイントを $p \colon \Cat{O} \lto \Ner(\FIN_*)$ と書き，コーンポイントと $D_{[0]}(k)$ を繋ぐ $\CATinfty{}_{/\Ner(\FIN_*)}$ の1-射を
    \begin{align}
        \sigma_k \coloneqq \lim_{\Cat{K}} (F \circ D)_{[1]}|_{\Delta^0_0 \times \Cat{K}_0} \bigl( \{0\} \to k \bigr) \in \inftyMap{\CATinfty{}_{/\Ner(\FIN_*)}} (p,\, p_k)_0
    \end{align}
    と書く．\eqref{eq:slice-1mor}から，これは $\CATinfty$ において
    \begin{align}
        \sigma_k = 
        \begin{tikzpicture}[baseline={([yshift=-.5ex]current bounding box.center)}]
        \path coordinate[bullet,label=below:$\Ner(\FIN_*)$] (x)
        +(120:3) coordinate[bullet,label=left:$\Cat{O}$] (y)
        +(60:3) coordinate[bullet,label=right:$\Oper{\Cat{O}}_k$] (z)
        ;
        \draw[->-=.5] (y) -- (z);
        \draw[->-=.5] (y) -- (x);
        \draw[->-=.5] (z) -- (x);
         \node[gray] at (barycentric cs:x=1,y=1,z=1) {$\sigma_k$};
                \begin{scope}[on background layer]
                    \fill[gray!20] (x.center) -- (y.center) -- (z.center);
                \end{scope}
        % \fill[red!20] (x) -- (y) -- (z) -- cycle;
    \end{tikzpicture}
    \in \CATinfty{}_2
    \end{align}
    と図示できる．故に，$f_k \coloneqq \partial^2_2 (\sigma_k) \in \bigl(\FUN(\Cat{O},\, \Oper{\Cat{O}}_k)^{\cong}\bigr)_0 \subset \Hom{\sSet} (\Cat{O},\, \Oper{\Cat{O}}_k)$ とおくと便利である．

    以下で $p$ が\hyperref[def:infty-operad]{$(\infty,\, 1)$-オペラッド}であることを示す．
    \begin{description}
        \item[\textbf{(Op-1)}]　
        
        勝手な\hyperref[def:inert-active]{inertな射} $\alpha \in \Hom{\FIN_*} \bigl( \expval{m},\, \expval{n} \bigr)$ および $x \in (\Cat{O}_{\expval{m}})_0$ を1つ固定する．
        このとき，$\Cat{O}$ における $\alpha$ の\hyperref[def:Cart-coCart]{$p$-coCartesian lift}全体が成す空間
        \begin{align}
            \inftyMap{\Cat{O}}^\mathrm{coCart} (x,\, \mhyphen\,)_\alpha \subset \FUN (\Delta^1,\, \Cat{O})
        \end{align}
        を構成しよう．

        　まず，\hyperref[def:infinity-1]{$(\infty,\, 1)$-圏の関手}
        \begin{align}
            \mathrm{coCart} \colon \Cat{K} \xrightarrow{D} \OPinfty \xrightarrow{F} \CATinfty{}_{/\Ner(\FIN_*)} \lto \CATinfty
        \end{align}
        を，$k \in \Cat{K}_0$ の行き先が $(1,\, 1)$-圏 $\sSet$ の引き戻し
        \begin{center}
            % https://tikzcd.yichuanshen.de/#N4Igdg9gJgpgziAXAbVABwnAlgFyxMJZABgBoBGAXVJADcBDAGwFcYkQAdDgW3pwAsATt2ABjCAGF6gnAAoAjgEoAviGWl0mXPkIoyAJmp0mrdl2AB5LqvWbseAkXKlDNBizaJOHC2hjAuKRxLZWUAfQBrYHCAjhgADzR3YG5QtQ0QDHsdJwojd1MvLgAxAFUAOQA9WJwE4PEpGVCAAi4AIywAc0ZZLgARGEYcekrnLlJWnz9YoJDwiPauwUUw2KY0fnobDKztRxR9F3yTT29YhKSmFOVrdLs93WRDgGZjjzMOcphBXo5igElymEAFSKO6ZLQOR7OV5uE4fXz+QJ8OaRcG7KG5YhvQpnLjrTa3WwQ7L7J55OHvIp-CrVLi1eL1STSHChX4DIYjMYcCZcREzFEWUKRMHEjE5A6kbGU3FcLBgABmOAAngBZehoWJfH4lQEglSTDrdX4XZKpUjjSamq5gG4cI3LNRGGBQTrwIigBWCCDcJDOEA4CBIQ4B+hYRjsfgQCARcFen1IZ40QPB4nx32IMgBoOIchp70ZrMpxDPfMJxAh4sANjLGYALMmcwBWZNhiNeKMxuMFpBVxtIFuh8OR6MRJg4bvlys5uu1xP9xAAdjnJYXAA4V4uFwBOVvDjujycZtc7mWnNAgGiMehtQYABUhEpAgi6-AnykoyiAA
            \begin{tikzcd}
                                                   & \{\alpha\} \arrow[r, hook]                                                                            & {\inftyMap{\Ner(\FIN_*)} \bigl(\expval{m},\, \expval{n}\bigr)}      \\
            \mathrm{coCart}_{[0]}(k) \arrow[d] \arrow[r] & {\FUN^{\text{coCart}} \bigl(\Delta^1,\, \Oper{\Cat{O}}_k\bigr)_{\alpha}} \arrow[d] \arrow[u] \arrow[r] & {\FUN^{\text{coCart}}(\Delta^1,\, \Oper{\Cat{O}}_k)} \arrow[u, "p_k"'] \\
            \{x\} \arrow[r, hook, "f_k"']                  & {\FUN^{\text{coCart}} \bigl(\{0\},\, \Oper{\Cat{O}}_k\bigr)_{\alpha}} = \Oper{\Cat{O}}_k{}_{\expval{m}} \arrow[d] \arrow[r]                                                    & \{\expval{m}\} \arrow[d, hook]                                      \\
                                                   & \Oper{\Cat{O}}_k \arrow[r, "p_k"']                                                                      & \Ner(\FIN_*)                                                       
            \end{tikzcd}
        \end{center}
        となるように構成する．特に\hyperref[def:equiv-infty]{$(\infty,\, 1)$-圏同値}
        \begin{align}
            \{O\} \times_{\Cat{O}_{\expval{m}}} \FUN (\Delta^1,\, \Cat{O})_\alpha \xrightarrow{\simeq}\lim_{k \in \Cat{K}} \Bigl( \{x\} \times_{\Oper{\Cat{O}_k}{}_{\expval{m}}} \FUN(\Delta^1,\, \Oper{\Cat{O}}_k)_{\alpha} \Bigr) 
        \end{align}
        が存在するので，
        \begin{align}
            \inftyMap{\Cat{O}}^\mathrm{coCart} (x,\, \mhyphen\,)_\alpha \coloneqq \lim_{\Cat{K}} \bigl( \mathcal{K} \xrightarrow{\mathrm{coCart}} \CATinfty \bigr)
        \end{align}
        と定義する．定義から明らかに $f \in \bigl(\inftyMap{\Cat{O}}^\mathrm{coCart} (x,\, \mhyphen\,)_\alpha\bigr)_0 \subset \Cat{O}_1$ は\hyperref[def:p-Cart]{$p$-coCartesian}である．

        \item[\textbf{(Op-2)}]　
        
        極限は有限積と交換するので明らか．

        \item[\textbf{(Op-3)}]　
        
        補題\ref{lem:Segal-cond}より，\hyperref[def:equiv-infty]{$(\infty,\, 1)$-圏同値}
        \begin{align}
            \Cat{O}_{\expval{n}} \xrightarrow{\simeq} (\Cat{O}_{\expval{1}})^{\times n}
        \end{align}
        を示せば良いが，これも極限が有限積と交換するので明らかである．
    \end{description}

\end{proof}

\subsection{{$\otimes$-sifted cocompleteness}}

$p \colon \Cat{V}^\otimes \lto \Ner(\FIN_*)$ を\hyperref[def:sym-monoidal-infty]{対称モノイダル $(\infty,\, 1)$-圏}とする．
このとき，$\forall n \ge 1$ に対して $\bm{n}$\textbf{-fold テンソル積}と呼ばれる \hyperref[def:infinity-1]{$(\infty,\, 1)$-圏の関手}が
\begin{align}
    \otimes \colon \Cat{V}^{\times n} \xleftarrow{\simeq} \Cat{V}^\otimes_{\expval{n}} \xrightarrow{\alpha_n{}_!} \Cat{V}
\end{align}
のようにして up to homotopy で定まる\footnote{Undelying $(\infty,\, 1)$-category $\Cat{V}$ を $\Cat{V} \coloneqq \Cat{V}^\otimes_{\expval{1}}$ と定義したのだった．}．ただし，$\Cat{V}^{\times n} \xleftarrow{\simeq} \Cat{V}^\otimes_{\expval{n}}$ は\hyperref[lem:Segal-cond]{条件 \textsf{\textbf{(Segal)}}}に登場した\hyperref[def:htpy-equiv]{ホモトピー同値}\eqref{eq:Segalmap}であり，
$\Cat{V}^\otimes_{\expval{n}} \xrightarrow{\alpha_n{}_!} \Cat{V}$ というのは\hyperref[def:alpha-n]{activeな射 $\alpha_n \in \Hom{\FIN_*} \bigl( \expval{n},\, \expval{1} \bigr)$} が誘導する\footnote{補題\ref{lem:Segal-cond}で $\rho^i$ の\hyperref[def:Cart-coCart]{$p$-coCartesian lift}が関手 $\Cat{O}^\otimes_{\expval{n}} \lto \Cat{O}$ を誘導したのと同じ方法である．} $(\infty,\, 1)$-圏の関手である．

\begin{mydef}[label=def:sifted-infty]{siftedな {$(\infty,\, 1)$-圏}}
    \hyperref[def:infinity-1]{$(\infty,\, 1)$-圏} $\Cat{K}$ が\textbf{sifted}であるとは，以下の2条件を満たすことを言う：
    \begin{description}
        \item[\textbf{(sifted-1)}] $\Cat{K} \neq \emptyset$
        \item[\textbf{(sifted-2)}] 対角関手 $\mathrm{diag} \colon \Cat{K} \lto \Cat{K} \times \Cat{K}$ が\textbf{cofinal}である．
        i.e. 任意の\hyperref[def:infty-fib]{右ファイブレーション} $q \colon \Cat{C} \lto \Cat{K} \times \Cat{K}$ に対して，$\mathrm{diag}$ の前合成
        \begin{align}
            \mathrm{diag}^* \colon \FUN_{/(\Cat{K} \times \Cat{K})} (\Cat{K} \times \Cat{K},\, \Cat{C}) \lto \FUN_{/(\Cat{K} \times \Cat{K})} (\Cat{K},\, \Cat{C})
        \end{align}
        が\hyperref[deftpy-equiv]{ホモトピー同値}である．
    \end{description}
    
\end{mydef}


\begin{mydef}[label=def:sifted-colim]{sifted {$(\infty,\, 1)$-colimit}}
    \hyperref[def:lim-colim-infty]{$(\infty,\, 1)$-圏の colimit} $\colim_{\Cat{K}} (\Cat{K} \xrightarrow{D} \Cat{C})$ が\textbf{sifted}であるとは，
    図式 $(\infty,\, 1)$-圏 $\Cat{K}$ が\hyperref[def:sifted-infty]{sifted}であることを言う．
\end{mydef}


\begin{mydef}[label=def:tensor-sifted]{{$\otimes$-sifted cocomplete}}
    \begin{itemize}
        \item \hyperref[def:sym-monoidal-infty]{対称モノイダル $(\infty,\, 1)$-圏} $p \colon \Cat{V}^\otimes \lto \Ner(\FIN_*)$
        \item 小さな\hyperref[def:infinity-1]{$(\infty,\, 1)$-圏} $\Cat{K}$
    \end{itemize}
    を与える．
    $\Cat{V}^\otimes$ の対称モノイダル構造が\textbf{形状 $\bm{\Cat{K}}$ のcolimitの上に分布する} (distributes over $\Cat{K}$-shaped colimits) とは，$\forall x \in \Cat{V}_0$ に対して
    \begin{align}
        x \otimes (\, \mhyphen \,) \colon \Cat{V} \cong \Delta^0 \times \Cat{V} \xrightarrow{x \times \Id_{\Cat{V}}} \Cat{V} \times \Cat{V} \xrightarrow{\otimes} \Cat{V}
    \end{align}
    が，形状 $\bm{\Cat{K}}$ の任意の\hyperref[def:lim-colim-infty]{colimit}を持つ\hyperref[def:infinity-1]{$(\infty,\, 1)$-圏の関手} $D \colon \Cat{K} \lto \Cat{V}$ と
    \begin{align}
        x \otimes (\colim_{\Cat{K}} D) = \colim_{\Cat{K}} (x \otimes D)
    \end{align}
    を充たすことを言う．

    \tcblower

    \hyperref[def:sym-monoidal-infty]{対称モノイダル $(\infty,\, 1)$-圏} $p \colon \Cat{V}^\otimes \lto \Ner(\FIN_*)$ が\textbf{ $\bm{\otimes}$-sifted cocomplete} であるとは，
    それが\hyperref[def:sifted-colim]{sifted-colimit}を持ち，かつ $\Cat{V}$ の対称モノイダル構造がsifted-colimitの上に分布することを言う．
\end{mydef}

\subsection{デカルト / 余デカルト対称モノイダル構造}



\subsection{{$\mathbb{E}_n$-オペラッド}}

\section{Stratified Disk algebras}



\end{document}